% Faithful transcription — original French. Transcribed from the original printing in the
% Journal de l'École Polytechnique, 2e série, 1er cahier (1895), printed pages 1–121:
% H. Poincaré, "Analysis situs". This is the complete memoir.
% Scan: Gallica, ark:/12148/bpt6k4337198 (views f7–f127), digitized by the Bibliothèque
% nationale de France from the copy of the Bibliothèque de l'École polytechnique.
%
% \origpage uses the PRINTED page numbers. Apparatus is NOT transcribed: the journal masthead
% on p. 1, the running heads, the page numbers, the signature lines ("J. E. P., 2e s. (C. no 1)"),
% library stamps, and the printer's ornaments closing the Introduction and p. 121.
%
% ORTHOGRAPHY. Only period glyph shapes are normalized: old-style figures are set as ordinary
% digits (this journal prints zero as a small "o" and one as a dotless "ı"), ligatures are
% expanded, and end-of-line hyphenation is dropped. The edition's own spelling is kept, including
% the unaccented capital "A" for "À" in running text and both "Équivalences" and "Equivalences"
% in the small headings (R12, R23). Headings set in capitals are normalized to sentence case,
% and the accents the capitals omitted are restored there ("Cas où p est impair").
%
% THE SCAN LOSES HAIRLINES. The Gallica images are 1749 px wide, and the thin minus of an index
% (p-1, q-2, n-p) and the period of a double index (\beta_{\lambda.p}) drop out of them
% systematically. Where a gap stands in their place and the index sequence requires them they
% are written in without a note (see notation.md). Full-size signs that lost strokes (=, \equiv,
% -, \pm) are written as the sign and carry an \ednote. Emphasis (\emph) was identified by eye,
% as a best-effort pass. The memoir has no figures.
%
% Notation and layout decisions shared across pages are recorded in notation.md.
\documentclass{article}
\usepackage{readmasters}
\begin{document}

\origpage{1}

\section*{Analysis situs;}

Par M.~H.~Poincaré.

\section*{Introduction.}

La Géométrie à $n$ dimensions a un objet réel; personne n'en doute aujourd'hui. Les êtres de l'hyperespace sont susceptibles de définitions précises comme ceux de l'espace ordinaire, et si nous ne pouvons nous les représenter, nous pouvons les concevoir et les étudier. Si donc, par exemple, la Mécanique à plus de trois dimensions doit être condamnée comme dépourvue de tout objet, il n'en est pas de même de l'Hypergéométrie.

La Géométrie, en effet, n'a pas pour unique raison d'être la description immédiate des corps qui tombent sous nos sens~: elle est avant tout l'étude analytique d'un groupe; rien n'empêche, par conséquent, d'aborder d'autres groupes analogues et plus généraux.

Mais pourquoi, dira-t-on, ne pas conserver le langage analytique et le remplacer par un langage géométrique, qui perd tous ses avantages dès que les sens ne peuvent plus intervenir. C'est que ce langage nouveau est plus concis; c'est ensuite que l'analogie avec la Géométrie ordinaire peut créer des associations d'idées fécondes et suggérer des généralisations utiles.
\origpage{2}

Peut-être ces raisons ne sont-elles pas suffisantes? Ce n'est pas assez, en effet, qu'une science soit légitime~: il faut que l'utilité ne puisse en être contestée. Tant d'objets divers sollicitent notre attention, que les plus importants ont seuls droit de l'obtenir.

Aussi y a-t-il des parties de l'Hypergéométrie auxquelles il n'y a pas lieu de beaucoup s'intéresser~: telles sont, par exemple, les recherches sur la courbure des surfaces dans l'espace à $n$ dimensions. On est sûr d'avance d'obtenir les mêmes résultats qu'en Géométrie ordinaire et l'on n'entreprend pas un long voyage pour retrouver des spectacles tout pareils à ceux que l'on rencontre chez soi.

Mais il y a des problèmes où le langage analytique serait tout à fait incommode.

On sait quelle est l'utilité des figures géométriques dans la théorie des fonctions imaginaires et des intégrales prises entre des limites imaginaires, et combien on regrette leur concours quand on veut étudier, par exemple, les fonctions de deux variables complexes.

Cherchons à nous rendre compte de la nature de ce concours; les figures suppléent d'abord à l'infirmité de notre esprit en appelant nos sens à son secours; mais ce n'est pas seulement cela. On a bien souvent répété que la Géométrie est l'art de bien raisonner sur des figures mal faites; encore ces figures, pour ne pas nous tromper, doivent-elles satisfaire à certaines conditions; les proportions peuvent être grossièrement altérées, mais les positions relatives des diverses parties ne doivent pas être bouleversées.

L'emploi des figures a donc avant tout pour but de nous faire connaître certaines relations entre les objets de nos études, et ces relations sont celles dont s'occupe une branche de la Géométrie que l'on a appelée \emph{Analysis Situs}, et qui décrit la situation relative des points des lignes et des surfaces, sans aucune considération de leur grandeur.

Il y a des relations de même nature entre les êtres de l'hyperespace; il y a donc une \emph{Analysis Situs} à plus de trois dimensions, comme l'ont montré Riemann et Betti.

Cette science nous fera connaître ce genre de relations, bien que cette connaissance ne puisse plus être intuitive, puisque nos sens nous font
\origpage{3}
défaut. Elle va ainsi, dans certains cas, nous rendre quelques-uns des services que nous demandons d'ordinaire aux figures de Géométrie.

Je me bornerai à trois exemples.

La classification des courbes algébriques en genres repose, d'après Riemann, sur la classification des surfaces fermées réelles, faite au point de vue de l'\emph{Analysis Situs}. Une induction immédiate nous fait comprendre que la classification des surfaces algébriques et la théorie de leurs transformations birationnelles sont intimement liées à la classification des surfaces fermées réelles de l'espace à cinq dimensions au point de vue de l'\emph{Analysis Situs}. M.~Picard, dans un Mémoire couronné par l'Académie des Sciences, a déjà insisté sur ce point.

D'autre part, dans une série de Mémoires insérés dans le \emph{Journal de Liouville}, et intitulés~: \emph{Sur les courbes définies par les équations différentielles}, j'ai employé l'\emph{Analysis Situs} ordinaire à trois dimensions à l'étude des équations différentielles. Les mêmes recherches ont été poursuivies par M.~Walther Dyck. On voit aisément que l'\emph{Analysis Situs} généralisée permettrait de traiter de même les équations d'ordre supérieur et, en particulier, celles de la Mécanique céleste.

M.~Jordan a déterminé analytiquement les groupes d'ordre fini contenus dans le groupe linéaire à $n$ variables. M.~Klein avait antérieurement, par une méthode géométrique d'une rare élégance, résolu le même problème pour le groupe linéaire à deux variables. Ne pourrait-on pas étendre la méthode de M.~Klein au groupe à $n$ variables, ou même à un \emph{groupe continu quelconque}? Je n'ai pu jusqu'ici y parvenir, mais j'ai beaucoup réfléchi à la question et il me semble que la solution doit dépendre d'un problème d'\emph{Analysis Situs} et que la généralisation du célèbre théorème d'Euler sur les polyèdres doit y jouer un rôle.

Je ne crois donc pas avoir fait une œuvre inutile en écrivant le présent Mémoire; je regrette seulement qu'il soit trop long; mais, quand j'ai voulu me restreindre, je suis tombé dans l'obscurité; j'ai préféré passer pour un peu bavard.
\origpage{4}

\section*{I.}

\subsection*{\S~1. --- Première définition des variétés.}

Soient $x_{1}, x_{2}, \ldots, x_{n}$ $n$ variables que je pourrai regarder comme les coordonnées d'un point dans l'espace à $n$ dimensions.

Jusqu'à nouvel ordre, je supposerai que ces $n$ variables demeurent toujours réelles.

Un système quelconque de valeurs de $n$ variables s'appellera \emph{un point}.

Considérons le système suivant formé de $p$ égalités et de $q$ inégalités~:
\[
\left\{
\begin{aligned}
&F_{1}(x_{1}, x_{2}, \ldots, x_{n}) = 0,\\
&F_{2}(x_{1}, x_{2}, \ldots, x_{n}) = 0,\\
&\ldots\ldots\ldots\ldots\ldots\ldots,\\
&F_{p}(x_{1}, x_{2}, \ldots, x_{n}) = 0,\\
&\varphi_{1}(x_{1}, x_{2}, \ldots, x_{n}) > 0,\\
&\varphi_{2}(x_{1}, x_{2}, \ldots, x_{n}) > 0,\\
&\ldots\ldots\ldots\ldots\ldots\ldots,\\
&\varphi_{q}(x_{1}, x_{2}, \ldots, x_{n}) > 0.
\end{aligned}
\right.
\tag{1}
\]

Je supposerai que les fonctions $F$ et $\varphi$ sont uniformes et continues et qu'elles ont des dérivées continues; je supposerai de plus que, si l'on forme le Tableau,
\[
\begin{matrix}
\frac{dF_{1}}{dx_{1}}, & \frac{dF_{1}}{dx_{2}}, & \ldots, & \frac{dF_{1}}{dx_{n}},\\[1ex]
\frac{dF_{2}}{dx_{1}}, & \frac{dF_{2}}{dx_{2}}, & \ldots, & \frac{dF_{2}}{dx_{n}},\\[1ex]
\ldots, & \ldots, & \ldots, & \ldots,\\[1ex]
\frac{dF_{p}}{dx_{1}}, & \frac{dF_{p}}{dx_{2}}, & \ldots, & \frac{dF_{p}}{dx_{n}},
\end{matrix}
\]
et que l'on forme les déterminants obtenus en prenant $p$ colonnes quelconques
\origpage{5}
dans ce Tableau, je supposerai, dis-je, que ces déterminants ne s'annulent jamais tous à la fois.

Je dirai que l'ensemble des points qui satisfont aux conditions (1), s'il y en a, ce que je suppose, forme une \emph{variété à $n-p$ dimensions}. (Si en particulier $p = 0$, de sorte qu'il n'y ait que des inégalités, j'aurai une variété à $n$ dimensions qui ne sera autre chose qu'une portion de l'espace à $n$ dimensions; je désignerai ordinairement cette variété sous le nom de \emph{domaine}.)

Deux cas peuvent se présenter. Soient $\alpha_{1}, \alpha_{2}, \ldots, \alpha_{n}$; $\beta_{1}, \beta_{2}, \ldots, \beta_{n}$ deux points satisfaisant aux conditions (1). Ou bien on pourra faire varier d'une manière continue $x_{1}, x_{2}, \ldots, x_{n}$ depuis $\alpha_{1}, \alpha_{2}, \ldots, \alpha_{n}$ jusqu'à $\beta_{1}, \beta_{2}, \ldots, \beta_{n}$, sans que les relations (1) cessent d'être satisfaites, et cela quelles que soient les valeurs de $\alpha_{1}, \alpha_{2}, \ldots, \alpha_{n}$ et de $\beta_{1}, \beta_{2}, \ldots, \beta_{n}$, pourvu que ces valeurs satisfassent aux conditions (1); ou bien cela ne sera pas toujours possible.

Dans le premier cas, nous dirons que la variété définie par les conditions (1) est \emph{continue}.

Dans ce qui va suivre, je ne considérerai que des variétés continues et, en ce qui concerne celles qui ne sont pas continues, je me bornerai à faire observer que l'on peut toujours les décomposer en un nombre fini ou infini de variétés continues.

Soit, par exemple, la variété
\[
x^{2}_{2} + x^{4}_{1} - 4x^{2}_{1} + 1 = 0.
\]

Ici $n = 2$ et le point $x_{1}, x_{2}$ est un point du plan; notre variété n'est alors autre chose qu'une courbe du 4ᵉ degré; mais, comme cette courbe se compose de deux branches de courbe fermées, notre variété n'est pas continue.

Mais nous pouvons la décomposer en deux autres, à savoir~:
\[
x^{2}_{2} + x^{4}_{1} - 4x^{2}_{1} + 1 = 0, \qquad x_{1} > 0
\]
et
\[
x^{2}_{2} + x^{4}_{1} - 4x^{2}_{1} + 1 = 0, \qquad x_{1} < 0
\]
\origpage{6}
et chacune d'elles, se réduisant à une seule branche de courbe fermée, est continue.

Je dirai qu'une variété est \emph{finie} si tous ses points satisfont à la condition
\[
x^{2}_{1} + x^{2}_{2} + \ldots + x^{2}_{n} < K^{2},
\]
$K$ étant une constante donnée.

Considérons maintenant le système de relations
\[
\left\{
\begin{aligned}
&F_{\alpha} = 0 &&(\alpha = 1, 2, \ldots, p),\\
&\varphi_{\beta} = 0,\\
&\varphi_{\gamma} > 0 &&(\gamma \gtrless \beta)
\end{aligned}
\right.
\tag{2}
\]
qui se compose de $p + 1$ égalités et de $q - 1$ inégalités.

Il peut arriver, ou bien qu'il n'y ait aucun point satisfaisant aux conditions (2), ou bien qu'il y en ait et, dans ce cas, ils formeront une variété, continue ou non, qui aura moins de $n - p$ dimensions.

L'ensemble des points qui satisferont à l'un des $q$ systèmes de relations~:
\[
\left\{
\begin{aligned}
&F_{\alpha} = 0, &&\varphi_{1} = 0, &&\varphi_{\gamma} > 0 &&(\gamma \gtrless 1),\\
&F_{\alpha} = 0, &&\varphi_{2} = 0, &&\varphi_{\gamma} > 0 &&(\gamma \gtrless 2),\\
&\ldots\ldots, &&\ldots\ldots, &&\ldots\ldots &&\ldots\ldots,\\
&F_{\alpha} = 0, &&\varphi_{q} = 0, &&\varphi_{\gamma} > 0 &&(\gamma \gtrless q)
\end{aligned}
\right.
\tag{3}
\]
s'appellera la \emph{frontière complète} de la variété définie par les conditions (1). Mais nous nous placerons quelquefois à un autre point de vue et nous ne considérerons comme une véritable variété frontière que celles qui auront $n - p - 1$ dimensions.

Il pourra se faire qu'il n'y ait aucune variété à $n - p - 1$ dimensions qui satisfasse à l'un des $q$ systèmes (3). Dans ce cas la variété définie par les conditions (1) sera dite \emph{illimitée}. Dans le cas contraire, elle sera dite \emph{limitée}.
\origpage{7}

Si une variété est à la fois finie, continue et illimitée, elle sera dite \emph{fermée}.

Pour abréger un peu le langage, nous donnerons le nom de \emph{surfaces} aux variétés à $n - 1$ dimensions, sauf dans le cas de $n = 2$; et nous donnerons, dans tous les cas, le nom de \emph{courbes} aux variétés à une dimension.

\subsection*{\S~2. --- Homéomorphisme.}

Considérons une substitution qui change $x_{1}, x_{2}, \ldots, x_{n}$ en $x_{1}, x'_{2}, \ldots, x'_{n}$\ednote{Printed $x_{1}$ without a prime; expected $x'_{1}$. The prime may be a misprint or lost ink.}, que j'assujettis seulement aux conditions suivantes.

On a
\[
x'_{i} = \varphi_{i}(x_{1}, x_{2}, \ldots, x_{n}) \qquad (i = 1, 2, \ldots, n).
\tag{4}
\]

Dans un certain domaine, les fonctions $\varphi_{i}$ sont uniformes, finies et continues; elles ont des dérivées continues, et leur déterminant fonctionnel ne s'annule pas.

Si l'on résout les équations (4) par rapport à $x_{1}, x_{2}, \ldots, x_{n}$, il vient
\[
x_{k} = \varphi'_{k}(x'_{1}, x'_{2}, \ldots, x'_{n}) \qquad (k = 1, 2, \ldots, n)
\]
et les fonctions $\varphi'_{k}$ satisfont aux mêmes conditions que les fonctions $\varphi_{i}$.

Il est clair que l'ensemble des substitutions qui satisfont à ces conditions forme un groupe, et ce groupe est un des plus généraux que l'on puisse imaginer. La Science dont l'objet est l'étude de ce groupe et de quelques autres analogues a reçu le nom d'\emph{Analysis Situs}.

Il est clair qu'une substitution de ce groupe transformera une variété de $m$ dimensions en une variété de $m$ dimensions, et que la variété transformée sera continue, ou finie, ou illimitée si la variété donnée l'est elle-même (et inversement).

Soient deux variétés d'un même nombre de dimensions définies $V$ et $V'$ respectivement par les conditions
\[
\left\{
\begin{aligned}
&F_{\alpha} = 0 &&(\alpha = 1, 2, \ldots, p),\\
&\varphi_{\beta} > 0 &&(\beta = 1, 2, \ldots, q),
\end{aligned}
\right.
\tag{1}
\]
\origpage{8}
et
\[
\left\{
\begin{aligned}
&F'_{\alpha} = 0 &&(\alpha = 1, 2, \ldots, p),\\
&\varphi'_{\beta} > 0 &&(\beta = 1, 2, \ldots, q').
\end{aligned}
\right.
\tag{1 \textit{bis}}
\]

Supposons que l'on puisse faire correspondre à un point $x_{1}, x_{2}, \ldots, x_{n}$ de la variété $V$ un point $x'_{1}, x'_{2}, \ldots, x'_{n}$ de la variété $V'$, de telle sorte que l'on ait
\[
x'_{k} = \psi_{k}(x_{1}, x_{2}, \ldots, x_{n}) \qquad (k = 1, 2, \ldots, n).
\tag{5}
\]

Je considère le domaine $D$ défini par les inégalités
\[
F_{\alpha} > -\varepsilon, \qquad F_{\alpha} < \varepsilon, \qquad \varphi_{\beta} > 0.
\]

La variété $V$ est évidemment contenue tout entière dans le domaine $D$.

Je suppose que \emph{dans le domaine} $D$ les fonctions $\psi_{k}$ sont finies, continues et uniformes, qu'elles ont des dérivées continues et que leur déterminant fonctionnel n'est jamais nul.

Résolvons maintenant les équations (5), nous trouverons
\[
x_{k} = \psi'_{k}(x'_{1}, x'_{2}, \ldots, x'_{n}) \qquad (k = 1, 2, \ldots, n).
\tag{6}
\]
\ednote{In the parenthesis of (6) the scan shows only a single bar where $=$ is expected, apparently the lower stroke of $=$ with the upper stroke unprinted.}

Je considère le domaine $D'$ défini par les inégalités
\[
F'_{\alpha} > -\varepsilon, \qquad F'_{\alpha} < \varepsilon, \qquad \varphi'_{\beta} > 0,
\]
et je suppose que \emph{dans le domaine} $D'$ les fonctions $\psi'_{k}$ sont finies, continues et uniformes, qu'elles ont des dérivées continues et que leur déterminant fonctionnel n'est pas nul.

Il résulte de ces hypothèses qu'à tout point de $V$ correspond un point de $V'$ et un seul et inversement; à toute variété $W$ contenue dans $V$ correspondra une variété $W'$, d'un même nombre de dimensions, contenue dans $V'$; si $W$ est continue, finie ou illimitée, il en sera de même de $W'$, et inversement.
\origpage{9}

Si toutes ces conditions sont remplies, nous dirons que les deux variétés $V$ et $V'$ sont équivalentes au point de vue de l'\emph{Analysis Situs}, ou, pour abréger le langage, qu'elles sont \emph{homéomorphes}, c'est-à-dire de forme pareille.

Je pourrai dire aussi que deux figures plus compliquées, composées d'un nombre quelconque de variétés, sont homéomorphes quand on passera de l'une à l'autre par une transformation de la forme (5).

Ainsi deux polygones d'un même nombre de côtés seront homéomorphes; deux polyèdres dont les faces seront en même nombre auront le même nombre de côtés et seront semblablement disposées, seront homéomorphes, etc.

\subsection*{\S~3. --- Deuxième définition des variétés.}

On peut définir les variétés d'une manière entièrement différente. Considérons $n$ équations
\[
\left\{
\begin{aligned}
&x_{1} = \theta_{1}(y_{1}, y_{2}, \ldots, y_{m}),\\
&x_{2} = \theta_{2}(y_{1}, y_{2}, \ldots, y_{m}),\\
&\ldots\ldots\ldots\ldots\ldots\ldots,\\
&x_{n} = \theta_{n}(y_{1}, y_{2}, \ldots, y_{m}).
\end{aligned}
\right.
\tag{8}
\]

Il est clair que ces équations (si les $y$ sont regardées comme des variables indépendantes) représentent une variété à $m$ dimensions.

On aura encore une variété à $m$ dimensions, si l'on adjoint aux équations (8) un certain nombre d'inégalités de la forme
\[
\psi(y_{1}, y_{2}, \ldots, y_{m}) > 0,
\]
qui limitent le champ des variations des variables $y$.

Je supposerai que les fonctions $\theta$ sont finies et continues; mais je puis faire une hypothèse de plus sans restreindre la généralité d'une façon essentielle (c'est d'ailleurs ce que je pouvais également faire tout aussi bien avec ma première définition des variétés).
\origpage{10}

Je puis supposer, dis-je, que les fonctions $\theta$ sont analytiques. Si, en effet, ces fonctions $\theta$ sont finies et continues, on pourra trouver des fonctions $\theta'$, qui seront analytiques et qui différeront des $\theta$ aussi peu que nous voudrons.

Je supposerai, en outre, que les déterminants fonctionnels de $m$ des fonctions $\theta$ par rapport aux $y$ ne s'annulent jamais tous à la fois.

On obtiendra une variété qui ne différera pas de la première, en remplaçant les $y$ dans les équations (8) par $m$ fonctions analytiques quelconques de $m$ variables $z_{1}, z_{2}, \ldots, z_{m}$.

La portée de cette nouvelle définition serait assez restreinte, si l'on ne pouvait l'augmenter par le procédé de la \emph{continuation analytique}.

Considérons deux variétés $V$ et $V'$ définies par des équations analogues à (8). Soient, par exemple, les équations
\[
x_{i} = \theta_{i}(y_{1}, y_{2}, \ldots, y_{m})
\tag{8}
\]
pour $V$, et les équations
\[
x_{i} = \theta'_{i}(y'_{1}, y'_{2}, \ldots, y'_{m})
\tag{8 \textit{bis}}
\]
\ednote{In (8 bis) the scan shows only a single bar at mid-height where $=$ is expected, apparently one stroke of $=$ with the other unprinted, as in (6) on p. 8.}
pour $V'$.

Il peut arriver que ces deux variétés aient une partie commune $V''$ ayant également $m$ dimensions; c'est-à-dire que tous les points de $V''$ appartiendront à la fois à $V$ et à $V'$.

Dans ce cas, à l'intérieur de $V''$, les $y$ seront des fonctions analytiques des $y'$ et inversement.

On dira alors que les deux variétés $V$ et $V'$ sont la \emph{continuation analytique} l'une de l'autre.

On peut former ainsi une chaîne de variétés
\[
V_{1}, \qquad V_{2}, \qquad \ldots, \qquad V_{n},
\]
de telle façon que chacune d'elles soit la continuation analytique de la
\origpage{11}
précédente et qu'entre deux variétés consécutives de la chaîne il y ait une partie commune. C'est ce que j'appellerai une \emph{chaîne continue}.

Il peut arriver aussi que la chaîne soit fermée, c'est-à-dire que $V_{n}$ ne diffère pas de $V_{1}$.

On peut avoir aussi un réseau de variétés, c'est-à-dire un ensemble de variétés telles que chacune d'elles soit la continuation de plusieurs autres et que l'on puisse passer d'une quelconque d'entre elles à une autre quelconque d'entre elles par continuation analytique; c'est ce que j'appellerai un \emph{réseau continu}.

On pourra alors considérer l'ensemble de toutes les variétés d'une même chaîne ou d'un même réseau comme formant une variété unique.

C'est là une définition plus étendue encore que la première.

Il y a, en effet, des variétés (et nous en verrons plus loin des exemples) qui peuvent être décomposées en un certain nombre de variétés partielles, formant une chaîne ou un réseau continus et telles que chacune d'elles puisse être définie par des équations de la forme (8) (des variétés qui, par conséquent, rentrent dans notre seconde définition) et qui cependant ne peuvent pas être définies par des relations de la forme (1) et ne rentrent pas dans notre première définition.

Au contraire, toute variété qui satisfait à la première définition satisfait aussi à la seconde.

Nous pouvons, en effet, d'après un théorème bien connu, si $y_{1}, y_{2}, \ldots, y_{n}$ sont définis par $n$ relations de la forme
\[
\left\{
\begin{aligned}
&y_{1} = F_{1}(x_{1}, x_{2}, \ldots, x_{n}),\\
&y_{2} = F_{2}(x_{1}, x_{2}, \ldots, x_{n}),\\
&\ldots\ldots\ldots\ldots\ldots\ldots,\\
&y_{n} = F_{n}(x_{1}, x_{2}, \ldots, x_{n}),
\end{aligned}
\right.
\tag{$\alpha$}
\]
si, dans un certain domaine, les $F$ sont des fonctions holomorphes des $x$, dont le déterminant fonctionnel ne s'annule pas; nous pouvons, dis-je, résoudre ces équations et en tirer
\[
x_{i} = \theta_{i}(y_{1}, y_{2}, \ldots, y_{n}),
\tag{$\beta$}
\]
\origpage{12}
les $\theta_{i}$ étant des fonctions des $y$ holomorphes dans un certain domaine.

Soit maintenant une variété $V$ satisfaisant à notre première définition, c'est-à-dire définie par les relations
\[
F_{\alpha} = 0, \qquad \varphi_{\beta} > 0.
\tag{1}
\]

Reprenons alors les équations ($\alpha$); prenons-y, pour $F_{1}, F_{2}, \ldots, F_{p}$, les premiers membres des $p$ égalités (1). Quant aux autres fonctions
\[
F_{p+1}, \qquad F_{p+2}, \qquad \ldots, \qquad F_{n},
\]
ce seront $n - p$ fonctions holomorphes quelconques des $x$. Je les assujettirai seulement à une condition.

Soit $x^{0}_{1}, x^{0}_{2}, \ldots, x^{0}_{n}$ un point quelconque $M_{0}$ de la variété $V$. Je m'arrangerai de telle façon que le déterminant fonctionnel des $n$ fonctions $F$ ne s'annule pas pour
\[
x_{i} = x^{0}_{i}.
\]

Cela est évidemment possible, puisque j'ai supposé que les déterminants fonctionnels des $p$ fonctions
\[
F_{1}, \qquad F_{2}, \qquad \ldots, \qquad F_{p},
\]
par rapport à $p$ quelconques des variables $x$, ne s'annulent jamais tous à la fois.

Je puis supposer encore que $F_{p+1}, F_{p+2}, \ldots, F_{n}$ s'annulent au point $M_{0}$, c'est-à-dire pour
\[
x_{i} = x^{0}_{i}.
\]

Alors, d'après le théorème cité plus haut, on peut résoudre les équations ($\alpha$) et l'on trouve que les $x_{i}$ s'expriment en séries ordonnées, suivant les puissances de
\[
y_{1}, \qquad y_{2}, \qquad \ldots, \qquad y_{n},
\]
et convergentes pourvu que ces quantités satisfassent à certaines inégalités.
\origpage{13}

Soient donc
\[
x_{i} = \theta_{i}
\tag{$\beta$}
\]
ces équations et
\[
\lambda_{k}(y_{1}, y_{2}, \ldots, y_{n}) > 0
\]
les conditions auxquelles les $y$ doivent satisfaire pour que les séries soient convergentes.

Faisons maintenant
\[
y_{1} = y_{2} = \ldots = y_{p} = 0;
\]
les $p$ premières équations ($\alpha$) ne différeront plus des $p$ équations (1); et les fonctions $\theta_{i}$, ne dépendant plus que de $n - p = m$ variables, seront de la forme (8).

Alors l'ensemble des relations
\[
x_{i} = \theta_{i}, \qquad \varphi_{\beta} > 0, \qquad \lambda_{k} > 0
\]
représentera une variété $v$, qui sera définie à la seconde manière et qui, ne différant pas de
\[
F_{\alpha} = 0, \qquad \varphi_{\beta} > 0, \qquad \lambda_{k} > 0,
\]
fera partie de $V$.

Le point $M_{0}$, qui est un point \emph{quelconque} de $V$, fait partie de $v$. On peut donc construire, autour d'un point quelconque de $V$, une variété analogue à $v$.

Le cas le plus simple est celui où les conditions de convergence $\lambda_{k} > 0$ sont une conséquence des inégalités $\varphi_{\beta} > 0$.

Alors $v$ ne diffère pas de $V$ et l'on peut se contenter, pour définir $v$, des égalités (8) et des inégalités
\[
\psi_{\beta} > 0,
\tag{9}
\]
\origpage{14}
que l'on obtient en substituant les $\theta$ à la place des $x$ dans les inégalités
\[
\varphi_{\beta} > 0.
\]

Remarquons en passant que tous les déterminants fonctionnels de $m$ des fonctions $\theta$ par rapport aux $y$ ne s'annuleront jamais à la fois.

Si les conditions $\lambda_{\beta} > 0$\ednote{Printed $\lambda_{\beta}$; the convergence conditions are written $\lambda_{k}$ on p. 13, so $\lambda_{k}$ is apparently meant.} n'étaient pas une conséquence des inégalités $\varphi_{\beta} > 0$, on partagerait la variété $V$ en variétés partielles; c'est ainsi, par exemple, que la variété $V$ peut évidemment être décomposée en deux variétés partielles
\[
F_{\alpha} = 0, \qquad \varphi_{\beta} > 0, \qquad \psi > 0
\]
et
\[
F_{\alpha} = 0, \qquad \varphi_{\beta} > 0, \qquad \psi < 0,
\]
$\psi$ étant une fonction quelconque de $x_{1}, x_{2}, \ldots, x_{n}$.

Cela posé, nous pourrons toujours partager $V$ en variétés partielles, ou mieux construire sur $V$ un certain nombre de variétés partielles, empiétant les unes sur les autres et assez petites pour que l'on puisse, pour chacune d'elles, trouver un système de variables auxiliaires, qui permette de représenter cette variété partielle par des égalités et des inégalités de la forme (8) et (9), satisfaisant à toutes les conditions que je viens d'énoncer. Tout point $M_{0}$ de $V$ appartiendra à l'une de ces variétés partielles, et l'ensemble de ces variétés formera un réseau continu. La première définition est ainsi ramenée à la seconde.

Pourtant, il reste une remarque à faire; il peut arriver qu'à deux systèmes différents de valeurs de $y_{1}, y_{2}, \ldots, y_{m}$ corresponde un même système de valeurs des fonctions $\theta_{1}, \theta_{2}, \ldots, \theta_{n}$ et, par conséquent, un même point de la variété $V$.

Dans ce cas, il convient d'adjoindre aux inégalités (9) d'autres inégalités
\[
\psi_{\gamma} > 0,
\tag{10}
\]
en les choisissant de telle sorte que des divers systèmes de valeurs des
\origpage{15}
variables $y$, qui correspondent à un même point de $V$, il y en ait toujours un et un seul qui satisfasse aux inégalités (9) et (10).

Soit, par exemple, un tore, dont l'équation sera
\[
(x^{2}_{1} + x^{2}_{2} + x^{2}_{3} + R^{2} - r^{2})^{2} - 4R^{2}(x^{2}_{1} + x^{2}_{2}) = 0.
\]
\ednote{The minus sign before $r^{2}$ prints as broken short strokes; read as a single $-$.}

Posons
\[
\begin{aligned}
&x_{1} = (R + r\cos y_{1})\cos y_{2},\\
&x_{2} = (R + r\cos y_{1})\sin y_{2},\\
&x_{3} = r\sin y_{1}.
\end{aligned}
\]

Nous voyons qu'à un même point du tore correspondront une infinité de systèmes de valeurs des $y$, compris dans la formule
\[
y_{1} + 2K_{1}\pi, \qquad y_{2} + 2K_{2}\pi,
\]
où les $K$ sont entiers.

Mais si nous assujettissons les $y$ aux conditions
\[
0 < y_{1} < 2\pi, \qquad 0 < y_{2} < 2\pi,
\]
nous n'aurons plus qu'un système de valeurs des $y$, qui correspondront à un point du tore.

\subsection*{\S~4. --- Variétés opposées.}

Avec l'une comme avec l'autre définition, il y a lieu de faire une distinction, dont on comprendra plus tard l'importance.

Supposons d'abord une variété $V$ définie de la première manière, c'est-à-dire par les relations
\[
F_{\alpha} = 0, \qquad \varphi_{\beta} > 0.
\]

On tiendra compte de l'ordre dans lequel sont rangées les équations $F_{\alpha} = 0$; si deux de ces équations sont permutées, on conviendra de dire
\origpage{16}
que le système de relations ne représente plus la variété $V$, mais une \emph{variété opposée} à $V$.

Nous pouvons maintenant remplacer les équations
\[
F_{\alpha} = 0 \qquad (\alpha = 1, 2, \ldots, p)
\]
par les suivantes~:
\[
\begin{aligned}
&\Phi_{1} = A_{1.1}F_{1} + A_{1.2}F_{2} + \ldots + A_{1.p}F_{p} = 0,\\
&\Phi_{2} = A_{2.1}F_{1} + A_{2.2}F_{2} + \ldots + A_{2.p}F_{p} = 0,\\
&\ldots\ldots\ldots\ldots\ldots\ldots,\\
&\Phi_{p} = A_{p.1}F_{1} + A_{p.2}F_{2} + \ldots + A_{p.p}F_{p} = 0,
\end{aligned}
\]
les $A_{i.k}$ étant des fonctions quelconques des $x$.

Si le déterminant $\Delta$ des coefficients $A_{i.k}$ ne s'annule pas dans le domaine considéré, les équations $\Phi_{\alpha} = 0$ seront équivalentes aux équations $F_{\alpha} = 0$, et, par conséquent, si on leur adjoint un certain nombre d'inégalités, elles représenteront encore la variété $V$ ou une variété opposée.

Nous conviendrons de dire que si $\Delta$ (qui a toujours le même signe, puisque j'ai supposé qu'il ne s'annule pas) est positif, ces équations représentent la variété $V$ et que, si $\Delta$ est négatif, elles représentent la variété opposée.

De même avec la seconde définition. Soit $V$ une variété définie par les équations
\[
x_{i} = \theta_{i}(y_{1}, y_{2}, \ldots, y_{m}).
\]

Il faudra tenir compte de l'ordre des variables $y$ et, si deux de ces variables sont permutées, nous conviendrons de dire que les équations ne représentent plus $V$, mais la variété opposée.

Imaginons que l'on remplace les $y$ par $m$ fonctions analytiques des $m$ variables nouvelles $z_{1}, z_{2}, \ldots, z_{m}$, de telle façon que l'on ait
\[
x_{i} = \theta'_{i}(z_{1}, z_{2}, \ldots, z_{m}),
\]
ces nouvelles équations représenteront encore $V$ ou la variété opposée.
\origpage{17}

Nous devons supposer que le déterminant fonctionnel $\Delta$ des $y$ par rapport aux $z$ ne s'annule pas, afin qu'à un système de valeurs des $y$ corresponde un seul système de valeurs des $z$. Il sera donc toujours de même signe.

Nous conviendrons de dire que, si $\Delta$ est positif, les nouvelles équations représentent encore $V$ et que si, $\Delta$ est négatif, elles représentent la variété opposée.

Voyons maintenant ce qui arrive lorsqu'on passe d'une définition à l'autre.

Soit une variété $V$ définie par
\[
F_{1} = F_{2} = \ldots = F_{p} = 0,
\]
et par certaines inégalités.

Adjoignons à ces $p$ équations les suivantes
\[
y_{1} = F_{p+1}, \qquad y_{2} = F_{p+2}, \qquad \ldots, \qquad y_{n-p} = F_{n},
\]
$F_{p+1}, F_{p+2}, \ldots, F_{n}$ étant $n - p$ fonctions quelconques des $x$.

Nous avons vu que si le déterminant fonctionnel $\Delta$ des $n$ fonctions $F_{1}, F_{2}, \ldots, F_{n}$ n'est pas nul, on peut résoudre ces $n$ équations par rapport aux $x$ et qu'on trouve ainsi $n$ équations
\[
x_{i} = \theta_{i}(y_{1}, y_{2}, \ldots, y_{n-p}),
\]
qui représentent une variété à $n - p$ dimensions. Mais on peut se demander si elles représenteront $V$ ou la variété opposée.

Nous conviendrons de dire qu'elles représentent $V$ si $\Delta$ est positif et la variété opposée si $\Delta$ est négatif.

Considérons la variété à $n - p$ dimensions
\[
F_{\alpha} = 0, \qquad \ldots, \qquad \varphi_{\beta} > 0,
\]
que j'appelle $V$. Nous avons vu que la variété à $n - p - 1$ dimensions
\origpage{18}
\[
F_{\alpha} = 0, \qquad \varphi_{\gamma} = 0, \qquad \varphi_{\beta} > 0, \qquad (\beta \gtrless \gamma),
\]
que j'appelle $v$, fait partie de la frontière de $V$.

Mais il importe de ranger les équations qui définissent $v$ dans l'ordre suivant~:
\[
F_{1} = 0, \qquad F_{2} = 0, \qquad \ldots, \qquad F_{p} = 0, \qquad \varphi_{\gamma} = 0,
\]
car si deux d'entre elles étaient permutées, nous conviendrions de dire que ces équations ne représentent plus la frontière de $V$ ou une partie de cette frontière, mais une variété opposée à cette frontière.

\subsection*{\S~5. --- Homologies.}

Considérons une variété $V$ à $p$ dimensions; soit maintenant $W$ une variété à $q$ dimensions ($q \leqq p$) faisant partie de $V$. Supposons que la frontière complète de $W$ se compose de $\lambda$ variétés continues à $q - 1$ dimensions
\[
v_{1}, \qquad v_{2}, \qquad \ldots, \qquad v_{\lambda}.
\]

Nous exprimerons ce fait par la notation
\[
v_{1} + v_{2} + \ldots + v_{\lambda} \backsim 0.
\]

Plus généralement la notation
\[
k_{1}v_{1} + k_{2}v_{2} \backsim k_{3}v_{3} + k_{4}v_{4},
\]
où les $k$ sont des entiers et les $v$ des variétés à $q - 1$ dimensions, signifiera qu'il existe une variété $W$ à $q$ dimensions faisant partie de $V$ et dont la frontière complète se composera de $k_{1}$ variétés peu différentes de $v_{1}$, de $k_{2}$ variétés peu différentes de $v_{2}$, de $k_{3}$ variétés peu différentes de la variété
\origpage{19}
opposée à $v_{3}$ et de $k_{4}$ variétés peu différentes de la variété opposée à $v_{4}$.

Les relations de cette forme pourront s'appeler des \emph{homologies}.

Les homologies peuvent se combiner comme des équations ordinaires. Nous emploierons aussi la notation suivante; supposons que l'on ait
\[
k_{1}v_{1} + k_{2}v_{2} + \ldots + k_{p}v_{p} \backsim w_{1} + w_{2} + \ldots + w_{q},
\]
et que les variétés $w_{1}, w_{2}, \ldots, w_{q}$ fassent partie de la frontière de $V$; nous écrirons quelquefois
\[
k_{1}v_{1} + k_{2}v_{2} + k_{p}v_{p} \backsim \varepsilon.
\]
\ednote{The terms $+ \ldots$ between $k_{2}v_{2}$ and $k_{p}v_{p}$, present in the preceding homology, are not printed here; an apparent misprint.}

\subsection*{\S~6. --- Nombres de Betti.}

Nous dirons que les variétés
\[
v_{1}, \qquad v_{2}, \qquad \ldots, \qquad v_{\lambda},
\]
d'un même nombre de dimensions et faisant partie de $V$, sont \emph{linéairement indépendantes}, si elles ne sont liées par aucune homologie à coefficients entiers.

S'il existe $P_{m} - 1$ variétés \emph{fermées} à $m$ dimensions faisant partie de $V$ et linéairement indépendantes et s'il n'en existe que $P_{m} - 1$, nous dirons que l'ordre de connexion de $V$ par rapport aux variétés à $m$ dimensions est égal à $P_{m}$.

Ainsi se trouvent définis, en ce qui concerne une variété $V$ à $m$ dimensions, $m - 1$ nombres que j'appellerai
\[
P_{1}, \qquad P_{2}, \qquad \ldots, \qquad P_{m-1},
\]
et qui sont les ordres de connexion de $V$ par rapport aux variétés de
\[
1, \qquad 2, \qquad \ldots, \qquad m - 1
\]
dimensions.
\origpage{20}

Je les appellerai dans la suite les \emph{nombres de Betti}.

Rendons ces définitions plus claires par un exemple~:

Soit $D$ un domaine faisant partie de l'espace ordinaire et limité par $n$ surfaces fermées
\[
S_{1}, \qquad S_{2}, \qquad \ldots, \qquad S_{n},
\]
qui ne se coupent pas.

Ce domaine est une variété à trois dimensions. Il admet donc deux nombres de Betti $P_{1}$ et $P_{2}$.

Cette variété est alors définie par les inégalités
\[
\varphi_{1} > 0, \qquad \varphi_{2} > 0, \qquad \ldots, \qquad \varphi_{n} > 0,
\]
si les équations des $n$ surfaces $S$ sont
\[
\varphi_{1} = 0, \qquad \varphi_{2} = 0, \qquad \ldots, \qquad \varphi_{n} = 0.
\]

Comme les surfaces ne se coupent pas, il n'y a pas de valeurs de $x_{1}, x_{2}, x_{3}$ qui satisfassent à la fois à deux de ces équations
\[
\varphi_{i} = 0, \qquad \varphi_{k} = 0.
\]

Les surfaces $S_{1}, S_{2}, \ldots, S_{n}$, ayant deux dimensions, n'auront qu'un seul nombre de Betti, qui sera l'ordre de connexion de Riemann; soient
\[
2Q_{1} + 1, \qquad 2Q_{2} + 1, \qquad \ldots, \qquad 2Q_{n} + 1,
\]
les ordres de connexion (qui sont impairs, puisque les surfaces sont fermées) des $n$ surfaces
\[
S_{1}, \qquad S_{2}, \qquad \ldots, \qquad S_{n},
\]
on aura
\[
P_{2} = n, \qquad P_{1} = Q_{1} + Q_{2} + \ldots + Q_{n} + 1.
\]

Ainsi, pour la région intérieure à une sphère
\[
P_{2} = 1, \qquad P_{1} = 1;
\]
\origpage{21}
pour la région comprise entre deux sphères
\[
P_{2} = 2, \qquad P_{1} = 1;
\]
pour la région intérieure à un tore
\[
P_{2} = 1, \qquad P_{1} = 2,
\]
pour la région comprise entre deux tores
\[
P_{2} = 2, \qquad P_{1} = 2.
\]

\subsection*{\S~7. --- Emploi des intégrales.}

Considérons une variété $V$ que l'on puisse représenter par les inégalités et égalités (8), (9) et (10) de façon à satisfaire à toutes les conditions énoncées plus haut.

On sait alors ce que l'on doit entendre par l'intégrale multiple d'ordre $m$
\[
\int F\,dy_{1}, dy_{2}, \ldots, dy_{m}
\]
étendue à la variété $V$; je désigne, bien entendu, par $F$ une fonction donnée des $y$. Il faut effectuer l'intégration successivement par rapport aux $m$ variables $y$, et les limites d'intégration seront définies par les inégalités (9) et (10).

Cela posé, je vais définir l'intégrale suivante~:
\[
\int \sum X_{\alpha_{1}\alpha_{2}\ldots\alpha_{m}}\,dx_{\alpha_{1}}\,dx_{\alpha_{2}} \ldots dx_{\alpha_{m}}. \tag{11}
\]

Les différentielles $dx_{\alpha_{1}}, dx_{\alpha_{2}}, \ldots, dx_{\alpha_{m}}$ sont $m$ quelconques des $n$ différentielles $dx_{1}, dx_{2}, \ldots, dx_{n}$. Les fonctions $X_{\alpha_{1}\alpha_{2}\ldots\alpha_{m}}$ sont des fonctions données de $x_{1}, x_{2}, \ldots, x_{n}$ et il y en a autant qu'il y a de combinaisons possibles des indices
\[
\alpha_{1}, \qquad \alpha_{2}, \qquad \ldots, \qquad \alpha_{m},
\]
\origpage{22}
c'est-à-dire qu'il y a de combinaisons de $n$ lettres $m$ à $m$. Il faut convenir que la fonction $X$ est nulle si deux de ses indices sont égaux et qu'elle change de signe quand on permute deux de ses indices.

Cela posé, l'intégrale (11) sera, par définition, égale à l'intégrale d'ordre $m$
\[
\int \sum X_{\alpha_{1}\alpha_{2}\ldots\alpha_{m}} \frac{\partial(x_{\alpha_{1}}, x_{\alpha_{2}}, \ldots, x_{\alpha_{m}})}{\partial(y_{1}, y_{2}, \ldots, y_{m})}\,dy_{1}, dy_{2}, \ldots, dy_{m}.
\]

Si maintenant la variété $V$ n'était pas susceptible d'être représentée par des relations de la forme (8), (9) et (10) satisfaisant à toutes les conditions énoncées, on décomposerait la variété $V$ en variétés partielles assez petites pour être susceptibles de ce mode de représentation, et l'intégrale (11), étendue à la variété totale $V$, serait par définition la somme des intégrales (11) étendues aux diverses variétés partielles.

Cette définition laisse toutefois subsister encore une ambiguité.

En effet, si l'on permute deux des lettres $y_{1}$ et $y_{2}$, l'intégrale change de signe; il importe donc de se donner l'ordre de ces lettres et la permutation de deux de ces lettres équivaudrait à un changement du sens de l'intégration dans l'étude des intégrales simples. Je dirai donc le sens de l'intégration pour parler de l'ordre dans lequel on convient de ranger les lettres $y_{1}, y_{2}, \ldots, y_{m}$.

J'ai eu l'occasion de m'occuper d'une question analogue dans un Mémoire \emph{sur les résidus des intégrales doubles}, inséré au Tome IX des \emph{Acta mathematica} et, en particulier, dans le \S~3 de ce Mémoire intitulé~: \emph{Conditions d'intégrabilité}.

J'ai recherché dans quels cas ces \emph{conditions d'intégrabilité} sont remplies, c'est-à-dire dans quels cas l'intégrale (11) est nulle toutes les fois qu'elle s'applique à une variété fermée.

Voici ce que j'ai trouvé; écrivons pour abréger l'écriture
\[
(\alpha_{1}, \alpha_{2}, \ldots, \alpha_{m})
\]
au lieu de $X_{\alpha_{1}, \alpha_{2}, \ldots, \alpha_{m}}$ et $[\alpha_{p}]$ au lieu de $x_{\alpha_{p}}$.
\origpage{23}

Nos conditions d'intégrabilité s'écriront
\[
\left\{
\begin{aligned}
&\frac{d(\alpha_{1}, \alpha_{2}, \ldots, \alpha_{m})}{d[\alpha_{m+1}]} \pm \frac{d(\alpha^{2}, \alpha^{3}, \ldots, \alpha_{m-1})}{d[\alpha_{1}]}\\
&\pm \frac{d(\alpha_{3}, \alpha_{1}, \ldots, \alpha_{m}, \alpha_{m+1}, \alpha_{1})}{d[\alpha_{2}]} \pm \ldots \pm \frac{d(\alpha_{m+1}, \alpha_{1}, \alpha_{2}, \ldots, \alpha_{m-1})}{d[\alpha_{m}]} = 0.
\end{aligned}
\right. \tag{12}
\]
\ednote{In the second term of (12) the print sets $\alpha^{2}, \alpha^{3}, \ldots, \alpha_{m-1}$, with superscripts and a final $\alpha_{m-1}$, where the cyclic pattern requires $\alpha_{2}, \alpha_{3}, \ldots, \alpha_{m+1}$; and the second index of the third term reads as $\alpha_{1}$ where $\alpha_{4}$ is expected, possibly a damaged 4. Apparent misprints, reproduced as printed.}

Voici la loi suivant laquelle doivent être choisis les signes $\pm$. On prendra toujours le signe $+$ si $m$ est pair, et alternativement le signe $+$ et le signe $-$ si $m$ est impair.

Il y aura autant d'équations (12) qu'il y a de systèmes possibles d'indices
\[
\alpha_{1}, \qquad \alpha_{2}, \qquad \ldots, \qquad \alpha_{m}, \qquad \alpha_{m+1},
\]
c'est-à-dire, puisque ces indices doivent être choisis parmi les lettres,
\[
1, \qquad 2, \qquad \ldots, \qquad n,
\]
autant qu'il y a de combinaisons de $n$ lettres $m + 1$ à $m + 1$.

Supposons maintenant que les conditions (12), au lieu d'être satisfaites pour toutes les valeurs possibles des $n$ variables
\[
x_{1}, \qquad x_{2}, \qquad \ldots, \qquad x_{n},
\]
le soient seulement pour certaines valeurs de ces variables. Par exemple, considérons une variété $V$ définie par les conditions
\[
F_{\alpha} = 0, \qquad \varphi_{\beta} > 0.
\]
Soit ensuite un domaine $D$ comprenant tous les points voisins de ceux de $V$ et défini, par exemple, par les conditions
\[
-\varepsilon < F_{\alpha} < \varepsilon, \qquad \varphi_{\beta} > -\varepsilon,
\]
$\varepsilon$ étant un très petit nombre positif.
\origpage{24}

Supposons que les conditions (12) soient satisfaites pour tous les points du domaine $D$.

En répétant le raisonnement du Mémoire cité et en le modifiant convenablement, on démontrera ce qui suit~: Soit une variété $V'$ à $m + 1$ dimensions faisant partie de $V$ (le nombre $m + 1$ doit donc être inférieur ou au plus égal au nombre des dimensions de $V$).

Supposons que la frontière complète de $V'$ se compose des $k$ variétés à $m$ dimensions
\[
W_{1}, \qquad W_{2}, \qquad \ldots, \qquad W_{k},
\]
de sorte que $W_{1} + W_{2} + \ldots + W_{k} \backsim 0$.

Alors si l'intégrale (11) satisfait aux conditions (12) dans le domaine $D$, la somme algébrique des intégrales (11) étendue aux variétés $W_{1}, W_{2}, \ldots, W_{k}$ est nulle. Il faut, bien entendu, pour chacune d'elles, faire attention au sens de l'intégration.

Les conditions (12) sont suffisantes pour qu'il en soit ainsi, mais elles ne sont pas nécessaires; ces conditions, nous l'avons vu, sont en nombre égal à celui des combinaisons de $n$ lettres $m + 1$ à $m + 1$; il suffirait, pour que le résultat que je viens d'énoncer fût encore exact, que l'intégrale (11) satisfît pour tous les points de $V$ à certaines conditions en nombre égal à celui des combinaisons de $n - p$ lettres $m + 1$ à $m + 1$, $n - p$ étant le nombre des dimensions de $V$.

Ces conditions seraient aisées à former, mais cela m'entraînerait trop loin de mon sujet.

Si alors une intégrale (11) satisfait aux conditions (12) dans le domaine $D$, les diverses valeurs que cette intégrale pourra prendre quand on l'étendra à diverses variétés fermées à $m$ dimensions faisant partie de $V$ seront des combinaisons linéaires à coefficients entiers d'un certain nombre d'entre elles que l'on pourra appeler les périodes de l'intégrale (11).

Le nombre maximum des périodes est égal à $P_{m} - 1$; car, si l'on considère $P_{m}$ variétés fermées quelconques à $m$ dimensions, il y aura toujours une variété à $m + 1$ dimensions qui admettra comme frontière complète
\origpage{25}
ces $P_{m}$ variétés ou quelques-unes d'entre elles. Il y aura donc toujours entre les $P_{m}$ intégrales correspondantes une relation linéaire à coefficients entiers. On pourrait d'ailleurs faire voir qu'il existe toujours des intégrales de la forme (11) pour lesquelles le nombre maximum des périodes est atteint. Cette manière de faire comprendre la définition des nombres de Betti a été employée par Betti lui-même pour le premier et le dernier de ces nombres, c'est-à-dire pour $P_{1}$ et $P_{m-1}$; mais nous venons de voir qu'il est aisé de faire de même pour les autres nombres de Betti.

\subsection*{\S~8. --- Variétés unilatères et bilatères.}

Considérons une variété $V$ définie à la seconde manière, c'est-à-dire formée d'une chaîne ou d'un réseau de variétés partielles dont chacune est définie elle-même par des relations de la forme (8) et (9).

Soit $v_{1}$ une variété partielle définie par les conditions
\[
\begin{aligned}
&x_{i} = \theta_{i}(y_{1}, y_{2}, \ldots, y_{m}),\\
&|y_{k}| < \beta_{k}.
\end{aligned}
\]

Soit $v_{2}$ une autre variété partielle et définie par les conditions
\[
\begin{aligned}
&x_{i} = \theta'_{i}(z_{1}, z_{2}, \ldots, z_{m}),\\
&|z_{k}| < \gamma_{k}.
\end{aligned}
\]

Supposons que ces deux variétés aient une partie commune $v'$ formant une variété continue. Je dis qu'à l'intérieur de cette variété le déterminant fonctionnel
\[
\Delta = \frac{\partial(y_{1}, y_{2}, \ldots, y_{m})}{\partial(z_{1}, z_{2}, \ldots, z_{m})}
\]
est toujours de même signe.

En effet, il ne pourrait changer de signe sans s'annuler ou sans devenir
\origpage{26}
infini. Or on a
\[
\Delta = \frac{\partial(x_{1}, x_{2}, \ldots, x_{m})}{\partial(z_{1}, z_{2}, \ldots, z_{m})} : \frac{\partial(x_{1}, x_{2}, \ldots, x_{m})}{\partial(y_{1}, y_{2}, \ldots, y_{m})},
\]
de sorte que $\Delta$ est lui-même le rapport de deux déterminants fonctionnels; comme ces deux déterminants fonctionnels sont essentiellement finis, $\Delta$ ne pourrait s'annuler que si l'on avait
\[
\frac{\partial(x_{1}, x_{2}, \ldots, x_{m})}{\partial(z_{1}, z_{2}, \ldots, z_{m})} = 0,
\]
et comme rien ne distingue les $m$ premières variables $x_{1}, x_{2}, \ldots, x_{m}$ des $n - m$ autres $x_{m+1}, x_{m+2}, \ldots, x_{n}$, il faudrait que le déterminant fonctionnel de $m$ quelconque des $x$ par rapport aux $z$ fût nul.

Tous les déterminants fonctionnels des fonctions $\theta'$ devraient donc s'annuler à la fois, ce qui n'a jamais lieu par hypothèse. $\Delta$ ne peut donc s'annuler, et l'on verrait absolument de même qu'il ne peut pas non plus devenir infini.

$\Delta$ est donc toujours de même signe et nous pourrons choisir l'ordre des variables $z$ de telle façon que ce signe soit toujours positif.

Une difficulté pourrait se présenter dans certains cas; supposons que la partie commune à $v_{1}$ et à $v_{2}$, au lieu de se réduire à une seule variété continue $v'$, se compose de plusieurs variétés continues $v'$, $v''$, $v'''$; dans chacune d'elles, le signe de $\Delta$ demeurera constant, mais il pourra changer lorsque l'on passera de l'une à l'autre. Dans ce cas, nous dirons que la variété $V$ est \emph{unilatère}.

Supposons que cette circonstance ne se présente pas, et considérons une suite de variétés partielles formant une chaîne fermée
\[
v_{1}, \qquad v_{2}, \qquad \ldots, \qquad v_{q}, \qquad v_{1}.
\]
Soient
\[
y^{i}_{1}, \qquad y^{i}_{2}, \qquad \ldots, \qquad y^{i}_{m}
\]
les $m$ variables qui jouent le rôle de $y_{1}, y_{2}, \ldots, y_{m}$ par rapport à $v_{i}$.
\origpage{27}

Supposons que $v_{1}$ et $v_{2}$ aient une partie commune $v'_{1}$, $v_{2}$ et $v_{3}$ une partie commune $v'_{2}$, \ldots, $v_{q-1}$ et $v_{q}$ une partie commune $v'_{q-1}$ et enfin $v_{q}$ et $v_{1}$ une partie commune $v'_{q}$.

Soit $\Delta_{i}$ le déterminant fonctionnel de
\[
y^{i+1}_{1}, \qquad y^{i+2}_{2}, \qquad \ldots, \qquad y^{i+1}_{m}
\]\ednote{Printed $y^{i+2}_{2}$ instead of $y^{i+1}_{2}$; an apparent misprint.}
par rapport à
\[
y^{i}_{1}, \qquad y^{i}_{2}, \qquad \ldots, \qquad y_{m}.
\]\ednote{Printed $y_{m}$ without the superscript; $y^{i}_{m}$ is expected. The superscript may have failed to print.}

Ce déterminant sera défini à l'intérieur de $v'_{i}$.

Je définis de même à l'intérieur de $v'_{q}$ le déterminant fonctionnel $\Delta_{q}$ de
\[
y^{1}_{1}, \qquad y^{1}_{2}, \qquad \ldots, \qquad y^{1}_{m}
\]
par rapport à
\[
y^{q}_{1}, \qquad q^{q}_{2}, \qquad \ldots, \qquad y^{q}_{m}.
\]\ednote{Printed $q^{q}_{2}$ instead of $y^{q}_{2}$; an apparent misprint.}

D'après ce que nous venons de voir, nous pouvons toujours choisir l'ordre des variables de telle façon que
\[
\Delta_{1}, \qquad \Delta_{2}, \qquad \ldots, \qquad \Delta_{q-1}
\]
soient toujours positifs. D'autre part, $\Delta_{q}$ est toujours de même signe, mais ce signe constant est-il le signe $+$ ou le signe $-$?

Si c'est le signe $-$, je dirai encore que la variété $V$ est \emph{unilatère}. Je pourrai dire aussi que les variétés $v_{1}, v_{2}, \ldots, v_{q}$ forment une chaîne unilatère.

Supposons maintenant que l'on ait construit un certain réseau continu de variétés partielles
\[
v_{1}, \qquad v_{2}, \qquad \ldots, \qquad v_{q}, \tag{4}
\]
de telle sorte que tout point de $V$ soit à l'intérieur (j'exclus la frontière) de l'une des variétés (4) ou de plusieurs d'entre elles. Si, dans la partie commune à deux quelconques des variétés (4), le déterminant $\Delta$ est positif, je dirai que la variété $V$ est \emph{bilatère}. Si elle ne l'était pas, il est clair
\origpage{28}
qu'on pourrait toujours, avec quelques-unes des variétés (4), former une chaîne unilatère et que la variété $V$ serait unilatère. Je pourrai dire aussi que le réseau des variétés (4) forme un système bilatère.

Mais pour justifier complètement notre définition, il faut faire voir que $V$ ne peut pas être à la fois unilatère et bilatère. Il est clair d'abord que, dans le système bilatère (4), on ne peut pas trouver de chaîne unilatère.

Mais il reste à montrer que le système (4) restera unilatère\ednote{Printed ``unilatère''; apparently for ``bilatère'', which is what the proof on p. 29 concludes.} quand on y adjoindra une variété quelconque $v_{q+1}$, faisant partie de $V$.

Soit $v'_{i}$ la partie commune à $v_{i}$ et à $v_{q+1}$; tout point de $v_{q+1}$ appartiendra au moins à l'une des variétés $v'_{i}$, et comme $v_{q+1}$ est continu, si je considère deux points $M_{1}$ et $M_{k}$ de $v_{q+1}$ appartenant respectivement à $v'_{1}$ et à $v'_{k}$, on pourra trouver des variétés intermédiaires que je pourrai appeler (puisque le numérotage reste arbitraire)~:
\[
v'_{1}, \qquad v'_{2}, \qquad \ldots, \qquad v'_{k},
\]
et qui formeront une chaîne.

Je pourrai toujours choisir l'ordre des variables analogues aux $y$ qui définissent la variété $v_{q+1}$, de telle façon que, dans $v'_{1}$, le déterminant analogue à $\Delta$ et relatif à $v_{1}$ et $v_{q+1}$ soit positif; je l'appellerai $\Delta_{1}$.

Appelons de même $\Delta_{i}$ le déterminant analogue à $\Delta$ et relatif à $v_{i}$ et $v_{q+1}$; $\Delta_{i}$ est défini à l'intérieur de $v'_{i}$.

Soit ensuite $\Delta'$ le déterminant relatif à $v_{1}$ et $v_{2}$; il sera défini dans toute la partie commune à ces deux variétés et, en particulier, dans la partie commune à $v'_{1}$ et $v'_{2}$; il sera positif puisque le système (4) est bilatère.

Alors, dans la partie commune à $v'_{1}$ et $v'_{2}$, $\Delta_{1}$ et $\Delta'$ seront positifs, il en résulte que leur rapport
\[
\Delta_{2} = \frac{\Delta_{1}}{\Delta'}
\]
est positif, et comme il conserve toujours le même signe, il sera positif pour tous les points de $v'_{2}$.

On démontrerait de proche en proche que $\Delta_{3}, \ldots, \Delta_{k}$ sont également positifs.
\origpage{29}

Le système restera donc bilatère après l'adjonction de $v_{q+1}$; il le restera encore si, dans l'expression des équations d'une des variétés $v_{i}$, on remplace les variables $y_{1}, y_{2}, \ldots, y_{m}$ par des fonctions holomorphes de variables nouvelles $y'_{1}, y'_{2}, \ldots, y'_{m}$ dont le déterminant fonctionnel ne s'annule jamais (ce qui est nécessaire si l'on veut qu'à un système de valeurs des $y$ corresponde un seul système de valeurs des $y'$). Il faut, bien entendu, choisir l'ordre des nouvelles variables $y'$ de telle façon que ce déterminant fonctionnel soit positif.

Le système restant toujours bilatère, on n'y pourra construire de chaîne unilatère, de sorte qu'une variété ne peut pas être à la fois bilatère et unilatère. C.~Q.~F.~D.

Tout le monde connaît l'exemple de surface unilatère que l'on obtient en ployant un rectangle de papier $ABCD$ (dont les côtés opposés sont $AB$ et $CD$ d'une part, $BC$ et $DA$ d'autre part), et en collant ensuite ensemble les côtés $AB$ et $CD$ de façon que $A$ soit collé avec $C$ et $B$ avec $D$.

Les exemples de variétés bilatères sont plus aisés encore à former. Ainsi dans l'espace à $n$ dimensions~:

1° Tout domaine à $n$ dimensions est bilatère;

2° Toute courbe à 1 dimension est bilatère;

3° Toute surface \emph{fermée} à $n - 1$ dimensions est bilatère.

Mais on peut aller plus loin.

Considérons une variété $V$ définie \emph{à la première manière}, c'est-à-dire par des égalités et des inégalités de la forme (1). Je dis qu'elle sera toujours bilatère.

Soient, en effet,
\[
F_{1} = F_{2} = \ldots = F_{p} = 0
\]
les $p$ égalités qui, avec quelques inégalités que nous n'écrirons pas, définissent $V$. Partageons $V$ en un certain nombre de variétés partielles $v$, définies par des relations de la forme (8) et (9). Soient
\[
v_{1}, \qquad v_{2}, \qquad \ldots, \qquad v_{q}, \qquad v_{1}
\]
\origpage{30}
un certain nombre de variétés partielles formant une chaîne continue, c'est-à-dire telles que chacune d'elles ait avec la suivante une partie commune. Je dis que cette chaîne est toujours bilatère.

En effet, nous avons supposé plus haut que les déterminants fonctionnels des $p$ fonctions $F$ par rapport à $p$ quelconques des variables $x$ ne s'annulent jamais tous à la fois. Je puis donc supposer que les variétés $v$ sont assez petites pour que, à l'intérieur de $v_{1}$, l'un des déterminants fonctionnels (que j'appellerai $\Delta_{1}$) ne s'annule pas; pour que, à l'intérieur de $v_{2}$, un autre de ces déterminants (que j'appellerai $\Delta_{2}$) ne s'annule pas, et ainsi de suite.

S'il n'en était pas ainsi, on n'aurait qu'à subdiviser chacune des variétés $v$ en variétés plus petites.

Soit, par exemple,
\[
\begin{aligned}
&\Delta_{1} = \frac{\partial(F_{1}, F_{2}, \ldots, F_{p})}{\partial(x_{\alpha_{1}}, x_{\alpha_{2}}, \ldots, x_{\alpha_{p}})},\\
&\Delta_{2} = \frac{\partial(F_{1}, F_{2}, \ldots, F_{p})}{\partial(x_{\beta_{1}}, x_{\beta_{2}}, \ldots, x_{\beta_{p}})},\\
&\ldots\ldots\ldots\ldots\ldots\ldots,
\end{aligned}
\]

Comme l'ordre des lettres $x_{\alpha_{1}}, x_{\alpha_{2}}, \ldots, x_{\alpha_{p}}$ reste arbitraire et que $\Delta_{1}$ conserve le même signe à l'intérieur de $v_{1}$, je pourrai toujours supposer que $\Delta_{1}$ est positif à l'intérieur de $v_{1}$. De même $\Delta_{2}$ sera positif à l'intérieur de $v_{2}$; et ainsi de suite.

Posons alors à l'intérieur de $v_{1}$
\[
\left\{
\begin{aligned}
&F_{1} = 0, &&F_{2} = 0, &&\ldots, &&F_{p} = 0,\\
&y^{1}_{1} = x'_{\alpha_{1}}, &&y^{1}_{2} = x'_{\alpha_{2}}, &&\ldots, &&y^{1}_{m} = x'_{\alpha_{m}}.
\end{aligned}
\right.
\tag{13}
\]

Ici $m = n - p$; les variables $x'_{\alpha_{i}}$ sont les $m = n - p$ variables $x$ qui restent quand on a supprimé les $p$ variables $x_{\alpha_{i}}$.

Posons de même à l'intérieur de $v_{2}$
\[
\left\{
\begin{aligned}
&F_{1} = 0, &&F_{2} = 0, &&\ldots, &&F_{p} = 0,\\
&y^{2}_{1} = x'_{\beta_{1}}, &&y^{2}_{2} = x'_{\beta_{2}}, &&\ldots, &&y^{2}_{m} = x'_{\beta_{m}}.
\end{aligned}
\right.
\tag{14}
\]
\origpage{31}

Les variables $x'_{\beta_{i}}$ sont les $m = n - p$ variables $x$ qui restent quand on a supprimé les $p$ variables $x_{\beta_{i}}$.

Et ainsi de suite.

Je supposerai que l'ordre des variables $x'_{\alpha_{i}}$ ait été choisi de telle sorte que l'on puisse passer de l'ordre normal des variables $x$, c'est-à-dire de l'ordre
\[
x_{1}, \qquad x_{2}, \qquad \ldots, \qquad x_{n}
\]
à l'ordre
\[
x_{\alpha_{1}}, \qquad x_{\alpha_{2}}, \qquad \ldots, \qquad x_{\alpha_{p}}, \qquad x'_{\alpha_{1}}, \qquad x'_{\alpha_{2}}, \qquad \ldots, \qquad x'_{\alpha_{m}},
\]
par une substitution du groupe alterné.

Alors, en résolvant les équations (13), on verrait qu'à l'intérieur de $v_{1}$ les $x$ sont des fonctions holomorphes de
\[
y^{1}_{1}, \qquad y^{1}_{2}, \qquad \ldots, \qquad y^{1}_{m}
\]
et ainsi de suite. En général, à l'intérieur de $v_{i}$, les $x$ seront des fonctions holomorphes de
\[
y^{i}_{1}, \qquad y^{i}_{2}, \qquad \ldots, \qquad y^{i}_{m}.
\]

Il s'agit alors de calculer le déterminant
\[
\frac{\partial(y^{2}_{1}, y^{2}_{2}, \ldots, y^{2}_{m})}{\partial(y^{1}_{1}, y^{1}_{2}, \ldots, y^{1}_{m})},
\]
ou, pour abréger,
\[
\frac{\partial y^{2}_{i}}{\partial y^{1}_{i}}
\]
à l'intérieur de la partie commune à $v_{1}$ et à $v_{2}$.

Pour cela, remplaçons les équations (13) et (14) par les équations plus générales
\[
\left\{
\begin{aligned}
&F_{1} = \lambda_{1}, &&F_{2} = \lambda_{2}, &&\ldots, &&F_{p} = \lambda_{p},\\
&y^{1}_{i} = x'_{\alpha_{i}}
\end{aligned}
\right.
\tag{13 \textit{bis}}
\]
\origpage{32}
et
\[
F_{k} = \lambda_{k}, \qquad y^{2}_{i} = x'_{\beta_{i}}.
\tag{14 \textit{bis}}
\]

Nous ferons ensuite $\lambda_{k} = 0$.

Résolvons les équations (13~\emph{bis}), nous aurons les $x$ en fonctions des $\lambda_{k}$ et des $y^{1}_{i}$ holomorphes à l'intérieur de $v_{1}$ et le déterminant fonctionnel de
\[
x_{1}, \qquad x_{2}, \qquad \ldots, \qquad x_{n}
\]
par rapport à
\[
\lambda_{1}, \qquad \lambda_{2}, \qquad \ldots, \qquad \lambda_{p}, \qquad y^{1}_{i}, \qquad y^{2}_{i}, \qquad \ldots, \qquad y^{m}_{i}
\]
\ednote{Printed $y^{1}_{i}, y^{2}_{i}, \ldots, y^{m}_{i}$; the variables of $v_{1}$ are apparently meant, i.e. $y^{1}_{1}, y^{1}_{2}, \ldots, y^{1}_{m}$ as in the denominator of the determinant below. Apparently a transposition of the indices.} sera évidemment
\[
\frac{1}{\Delta_{1}}.
\]

Résolvons de même les équations (14~\emph{bis}), nous aurons les $x$ en fonctions des $\lambda_{k}$ et des $y^{2}_{i}$, holomorphes à l'intérieur de $v_{2}$ et le déterminant fonctionnel sera $\frac{1}{\Delta_{2}}$.

Il viendra alors, pour les points communs à $v_{1}$ et à $v_{2}$,
\[
\frac{\partial y^{2}_{i}}{\partial y^{1}_{i}} = \frac{\partial(\lambda_{1}, \lambda_{2}, \ldots, \lambda_{p}, y^{2}_{1}, y^{2}_{2}, \ldots, y^{2}_{m})}{\partial(\lambda_{1}, \lambda_{2}, \ldots, \lambda_{p}, y^{1}_{1}, y^{1}_{2}, \ldots, y^{1}_{m})} = \frac{\Delta_{1}}{\Delta_{2}}.
\]

Il ne reste plus ensuite qu'à faire $\lambda_{k} = 0$ dans les expressions des $x$.

On trouverait de même, dans la partie commune à $v_{h}$ et $v_{h+1}$,
\[
\frac{\partial y^{h+1}_{i}}{\partial y^{h}_{i}} = \frac{\Delta_{h}}{\Delta_{h+1}},
\]
et dans la partie commune à $v_{q}$ et à $v_{1}$
\[
\frac{\partial y^{1}_{i}}{\partial y^{q}_{i}} = \frac{\Delta_{q}}{\Delta_{1}}.
\]

Nous avons vu plus haut que l'on peut toujours supposer que $\Delta_{i}$ est,
\origpage{33}
à l'intérieur de $v_{i}$, positif. Tous ces rapports sont donc positifs et la chaîne est bilatère. C.~Q.~F.~D.

Toutes les variétés qui rentrent dans la première définition sont donc bilatères et, comme j'ai cité plus haut un exemple de variété unilatère satisfaisant à la deuxième définition, nous devons conclure que les variétés qui satisfont à la deuxième définition ne satisfont pas toutes à la première. C'est ce que j'avais annoncé plus haut.

Toute variété unilatère est par définition \emph{opposée à elle-même}.

\subsection*{\S~9. --- Intersection de deux variétés.}

Soient $V$ et $V'$ deux variétés définies à la seconde manière, ayant l'une $p$ dimensions, l'autre $n-p$ dimensions, et soient $M_{0}$ et $M'_{0}$ deux points appartenant respectivement à $V$ et à $V'$ et ayant respectivement pour coordonnées
\[
x^{0}_{1}, \qquad x^{0}_{2}, \qquad \ldots, \qquad x^{0}_{n}
\]
et
\[
x'^{0}_{1}, \qquad x'^{0}_{2}, \qquad \ldots, \qquad x'^{0}_{n}.
\]

Autour du point $M_{0}$, construisons sur $V$ une variété partielle $v$, analogue à celles que nous avons envisagées dans le \no 3, de telle façon qu'en un point $x_{1}, x_{2}, \ldots, x_{n}$ de $v$ on ait
\[
x_{i} = \theta_{i}(y, y_{2}, \ldots, y_{p}) \qquad (i = 1, 2, n).
\]
\ednote{The first argument is printed $y$ with no index, apparently for $y_{1}$; the index may have failed to print. The range is printed $i = 1, 2, n$, apparently for $i = 1, 2, \ldots, n$.}

De même autour de $M'_{0}$, construisons sur $V'$ une variété $v'$ telle qu'en un point $x'_{1}, x'_{2}, \ldots, x'_{n}$ de $v'$ on ait
\[
x'_{i} = \theta'_{i}(y'_{1}, y'_{2}, \ldots, y'_{n-p}) \qquad (i = 1, 2, n).
\]
\ednote{The range is printed $i = 1, 2, n$, apparently for $i = 1, 2, \ldots, n$.}
\origpage{34}

Envisageons le déterminant
\[
\begin{vmatrix}
\frac{dx_{1}}{dy_{1}} & \frac{dx_{2}}{dy_{1}} & \ldots & \frac{dx_{n}}{dy_{1}} \\[1ex]
\frac{dx_{1}}{dy_{2}} & \frac{dx_{2}}{dy_{2}} & \ldots & \frac{dx_{n}}{dy_{2}} \\[1ex]
\ldots & \ldots & \ldots & \ldots \\[1ex]
\frac{dx_{1}}{dy_{p}} & \frac{dx_{2}}{dy_{p}} & \ldots & \frac{dx_{n}}{dy_{p}} \\[1ex]
\frac{dx'_{1}}{dy'_{1}} & \frac{dx'_{2}}{dy'_{1}} & \ldots & \frac{dx'_{n}}{dy'_{1}} \\[1ex]
\ldots & \ldots & \ldots & \ldots \\[1ex]
\frac{dx'_{1}}{dy'_{n-p}} & \frac{dx'_{2}}{dy'_{n-p}} & \ldots & \frac{dx'_{n}}{dy'_{n-p}}
\end{vmatrix}.
\]

C'est une fonction de $y_{1}, y_{2}, \ldots, y_{p}, y'_{1}, y'_{2}, \ldots, y'_{n-p}$; cette fonction dépendra donc de la position des points $M$ et $M'$ de coordonnées $x_{1}, x_{2}, \ldots, x_{n}$ et $x'_{1}, x'_{2}, \ldots, x'_{n}$ sur les variétés $v$ et $v'$. Je l'appellerai donc
\[
f(M, M').
\]

Je puis changer de variables en remplaçant $y_{1}, y_{2}, \ldots, y_{p}$ par des fonctions holomorphes de $p$ variables nouvelles $z_{1}, z_{2}, \ldots, z_{p}$ choisies de telle sorte qu'à un système de valeurs de $y$ corresponde un seul système de valeurs des $z$. Il faut pour cela que le déterminant fonctionnel des $y$ par rapport aux $z$ ne s'annule pas et je puis toujours supposer, en rangeant les $z$ dans un ordre convenable, que ce déterminant est positif.

La fonction $f(M, M')$ est alors multipliée par ce déterminant fonctionnel et, par conséquent, conserve son signe.

Il en est encore de même quand on fait un changement de variables analogues sur la variété $v'$.

Considérons maintenant une autre variété $v_{1}$, analogue à $v$, construite sur $V$, et une variété $v'_{1}$, analogue à $v'$, construite sur $V'$.

Soit $M_{1}$ un point de $v_{1}$ et $M'_{1}$ un point de $v'_{1}$. Nous pourrons construire une fonction analogue à $f(M, M')$ et que j'appellerai $f_{1}(M_{1}, M'_{1})$.
\origpage{35}

Supposons maintenant que les variétés $v_{1}$ et $v$ (de même que les variétés $v'_{1}$ et $v'$) aient une partie commune, et que les points $M$ et $M_{1}$ se confondent, de même que les points $M'$ et $M'_{1}$.

Comparons les deux fonctions
\[
f(M, M'), \qquad f_{1}(M, M').
\]

Supposons que les deux variétés $V$ et $V'$ soient bilatères. Nous pourrons supposer alors (en choisissant convenablement l'ordre des variables $y$ relatives à $v_{1}$), que le déterminant analogue à celui que nous avons appelé $\Delta$ dans le \no 8, et relatif à $v$ et $v_{1}$, est positif dans la partie commune à $v$ et $v_{1}$, de même le déterminant analogue à $\Delta$ et relatif à $v'$ et $v'_{1}$ sera également positif.

Alors $f$ et $f_{1}$ seront de même signe.

Soit alors
\[
S(M, M')
\]
une fonction qui sera égale à $+1$, à $-1$ ou à 0, selon que $f(M, M')$ sera positif, négatif ou nul. Cette fonction sera alors parfaitement déterminée, elle ne dépendra que de la position des points $M$ et $M'$ sur $V$ et $V'$; elle ne dépendra pas de la manière dont les variétés $v$ et $v'$ auront été construites autour de $M$ et de $M'$.

\emph{Cela ne serait plus vrai si l'une des variétés $V$ et $V'$ était unilatère, ce que je ne suppose pas}.

Supposons, en particulier, que les points $M$ et $M'$ se confondent et, par conséquent, que le point $M$ soit un point d'intersection de $V$ et de $V'$, la considération de la fonction $S(M, M')$ présente alors un grand intérêt.

Supposons que l'on forme la fonction $S(M, M)$ pour tous les points d'intersection $M$ de $V$ et de $V'$, et faisons la somme de toutes les fonctions $S$ ainsi obtenues; je désignerai cette somme par
\[
N(V, V').
\]
\origpage{36}

Cette définition de $N(V, V')$ s'applique quand les variétés $V$ et $V'$ ont été définies comme au \no 3. Mais on peut la simplifier quand $V$ et $V'$ ont été définies comme au \no 1.

Soient alors
\[
F_{\alpha} = 0, \qquad \varphi_{\beta} > 0 \qquad (\alpha = 1, 2, \ldots, p; \beta = 1, 2, \ldots, q)
\]
les conditions qui définissent $V$, et soient
\[
F'_{\gamma} = 0, \qquad \varphi'_{\delta} > 0 \qquad (\gamma = 1, 2, \ldots, n-p; \delta = 1, 2, \ldots, q'),
\]
celles qui définissent $V'$. $V$ aura $n-p$ dimensions et $V'$ en aura $p$.

Considérons un point $M$ de $V$ ayant pour coordonnées $x_{1}, x_{2}, \ldots, x_{n}$ et un point $M'$ de $V'$ ayant pour coordonnées $x'_{1}, x'_{2}, \ldots, x'_{n}$ et formons le déterminant
\[
\begin{vmatrix}
\frac{dF_{1}}{dx_{1}} & \frac{dF_{2}}{dx_{1}} & \ldots & \frac{dF_{p}}{dx_{1}} & \frac{dF'_{1}}{dx'_{1}} & \frac{dF'_{2}}{dx'_{1}} & \ldots & \frac{dF'_{n-p}}{dx'_{1}} \\[1ex]
\ldots & \ldots & \ldots & \ldots & \ldots & \ldots & \ldots & \ldots \\[1ex]
\frac{dF_{1}}{dx_{n}} & \frac{dF_{2}}{dx_{n}} & \ldots & \frac{dF_{p}}{dx_{n}} & \frac{dF'_{1}}{dx'_{n}} & \frac{dF'_{2}}{dx'_{n}} & \ldots & \frac{dF'_{n-p}}{dx'_{n}}
\end{vmatrix}
\]
que j'appellerai $\psi(M, M')$. Quelle relation y a-t-il entre $\psi(M, M')$ et $S(M, M')$?

Pour nous en rendre compte, posons
\[
x'_{i} = x_{i} + h_{i} \qquad (i = 1, 2, \ldots, n);
\]
faisons ensuite varier à la fois les $x'_{i}$ et les $x_{i}$, mais de telle façon que leur différence $h_{i}$ demeure constante.

Posons
\[
\begin{aligned}
&F_{\alpha}(x_{i}) = y_{\alpha},\\
&F'_{\gamma}(x_{i} + h_{i}) = y'_{\gamma}.
\end{aligned}
\]

Si le déterminant $\psi(M, M')$ n'est pas nul, nous pourrons résoudre ces
\origpage{37}
équations par rapport aux $x_{i}$ et nous en tirerons
\[
x_{i} = \theta_{i}(y_{\alpha}, y'_{\gamma}).
\]

Le déterminant fonctionnel des $\theta_{i}$ par rapport aux $y_{\alpha}$ et aux $y'_{\gamma}$ sera alors $\frac{1}{\psi(M, M')}$.

Si nous faisons $y_{\alpha} = 0$, nous aurons
\[
x_{i} = \theta_{i}(0, y'_{\gamma});
\]
ce sera là l'équation d'une variété à $n-p$ dimensions qui fera partie de $V$; de même en faisant
\[
x'_{i} = h_{i} + \theta_{i}(y_{\alpha}, 0),
\]
on aura l'équation d'une variété à $p$ dimensions qui fera partie de $V'$.

Si l'on définit de la sorte les variétés $v$ et $v'$ du commencement de ce numéro, on voit que le déterminant fonctionnel des $\theta_{i}$ n'est autre chose que $f(M, M')$, d'où
\[
f(M, M')\psi(M, M') = 1.
\]

Supposons maintenant que les deux variétés $V$ et $V'$ fassent toutes deux parties d'une variété $W$ d'un plus grand nombre de dimensions. Soient, par exemple,
\[
F_{\alpha} = 0, \qquad \varphi_{\beta} > 0 \qquad (\alpha = 1, 2, \ldots, p; \beta = 1, 2, \ldots, q) \tag{1}
\]
les relations qui définissent $W$ qui aura ainsi $n-p$ dimensions.

Soient
\[
F'_{\gamma} = 0, \qquad \varphi'_{\delta} > 0 \qquad (\gamma = 1, 2, \ldots, p; \delta = 1, 2, \ldots, q') \tag{2}
\]
\ednote{The range of $\gamma$ is printed as ending at $p$; apparently for $p'$, since $V$ is said to have $n-p-p'$ dimensions and the next system runs to $p''$.} les relations qui jointes à (1) définissent $V$ qui aura ainsi $n-p-p'$ dimensions.

Soient enfin
\[
F''_{\varepsilon} = 0, \qquad \varphi''_{\zeta} > 0 \qquad (\varepsilon = 1, 2, \ldots, p''; \zeta = 1, 2, \ldots, q'') \tag{3}
\]
\origpage{38}
les relations qui jointes à (1) définiront $V'$ qui aura ainsi $n-p-p''$ dimensions. Je suppose $p + p' + p'' = n$.

Soit $M$ un point d'intersection de $V$ et de $V'$.

Je formerai le déterminant fonctionnel des $F_{\alpha}$, des $F'_{\gamma}$ et des $F''_{\varepsilon}$ par rapport aux $x$, et ce sera ce déterminant que j'appellerai $\psi(M)$.

Alors $S(M)$ sera par définition égal à $+1$ ou à $-1$, selon que $\psi(M)$ sera positif ou négatif.

Je désignerai ensuite par
\[
N(V, V')
\]
la somme de toutes les quantités $S(M)$ relatives à tous les points d'intersection de $V$ et de $V'$.

Il est clair que, si l'une des variétés $V$ ou $V'$ est remplacée par la variété opposée, l'expression $N(V, V')$ change de signe.

Soit maintenant $V$ une variété fermée à 1 dimension et soit $W$ un domaine à $n$ dimensions dont la frontière complète se compose des variétés à $n-1$ dimensions
\[
V_{1}, \qquad V_{2}, \qquad \ldots, \qquad V_{k},
\]
de telle sorte que
\[
V_{1} + V_{2} + \ldots + V_{k} \backsim 0.
\]

Je dis que
\[
N(V, V_{1}) + N(V, V_{2}) + \ldots + N(V, V_{k}) = 0.
\]

En effet, la variété $V$ va se composer d'un certain nombre de variétés continues à 1 dimension, c'est-à-dire d'un certain nombre de lignes fermées. Les équations d'une de ces lignes pourront se mettre sous la forme
\[
x_{i} = \theta_{i}(t),
\]
$t$ étant une variable que nous ferons croître de $-\infty$ à $+\infty$.

Comme la ligne est fermée, nous aurons
\[
\theta_{i}(-\infty) = \theta_{i}(+\infty).
\]
\origpage{39}

Envisageons maintenant un point $M$ d'intersection d'une de ces lignes avec la frontière de $W$; si, quand on fait croître $t$, cette ligne passe de l'intérieur de $W$ à l'extérieur, on aura
\[
S(M) = 1.
\]
Si elle passe de l'extérieur à l'intérieur $S(M)$ sera égal à $-1$. Mais la ligne étant fermée reviendra à son point de départ, de sorte que la somme de tous les $S(M)$ devra être nulle. C.~Q.~F.~D.

Supposons maintenant que toutes les variétés que nous envisageons et, par exemple, $V$, $W$, $V_{1}, V_{2}, \ldots, V_{k}$ fassent partie d'une variété $U$ et que ce soit cette variété $U$ dont nous étudions les ordres de connexion et par rapport à laquelle nous prenons les homologies.

(Voici ce que je veux dire; quand, au \no 5, j'ai défini les homologies, il y avait une certaine variété $V$ qui jouait dans cette définition un rôle important; eh bien, c'est ici $U$ qui jouera ce rôle.)

Je suppose que $U$ ait $h$ dimensions, $V$ une dimension, $W$ $h$ dimensions, $V_{1}, V_{2}, \ldots, V_{k}$ $h-1$ dimensions.

Il n'y a rien à changer à ce qui précède et l'on voit que, si $V$ est fermé et à une dimension et que
\[
V_{1} + V_{2} + \ldots + V_{k} \backsim 0,
\]
on aura
\[
N(V, V_{1}) + N(V, V_{2}) + \ldots + N(V, V_{k}) = 0.
\]

Je dis que cela sera encore vrai si la variété à une dimension $V$ n'est plus fermée, mais a ses deux extrémités sur la frontière de $U$ et si $W$ n'a aucun point sur la frontière de $U$.

En effet, $U$ se trouve partagé en deux régions continues ou non, à savoir~: $W$ et la région $R$ extérieure à $W$. Par hypothèse, la frontière de $U$ appartient tout entière à $R$.

La ligne $V$ n'étant pas fermée, le point initial correspondant à $t = -\infty$ ne coïncidera plus avec le point final correspondant à $t = +\infty$; mais,
\origpage{40}
par hypothèse, ces deux points appartiennent tous deux à la frontière de $U$ et, par conséquent, à $R$.

Notre ligne passera donc précisément autant de fois de $R$ dans $W$ que de $W$ dans $R$; donc la somme des $S(M)$ est nulle. C.~Q.~F.~D.

Je vais maintenant établir la réciproque.

Soient $V_{1}, V_{2}, \ldots, V_{k}$, $k$ variétés fermées à $h-1$ dimensions situées dans $U$, $U$ n'ayant aucun point commun avec la frontière de $V_{i}$. Supposons que l'on n'ait pas
\[
V_{1} + V_{2} + \ldots + V_{k} \backsim 0,
\]
je dis que l'on pourra toujours trouver une ligne $V$, située dans $U$, qui sera fermée, on\ednote{Printed ``on''; apparently for ``ou'', a turned letter.} aura ses deux extrémités sur la frontière de $U$ et qui sera telle que
\[
N(V, V_{1}) + N(V, V_{2}) + \ldots + N(V, V_{k}) \gtrless 0.
\]

En effet, ou bien les variétés $V_{1}, V_{2}, \ldots, V_{k}$ ne partageront pas $U$ en plusieurs régions distinctes. Alors on pourra aller d'un point quelconque de $U$ à un autre point quelconque de $U$ sans rencontrer aucune des variétés $V_{i}$.

Considérons alors une ligne très petite rencontrant $V_{1}$ en un point $M$ et ne rencontrant aucune autre des variétés $V_{i}$; nous pourrons joindre les deux extrémités de cette ligne par une autre ligne continue qui ne rencontrera aucune des variétés $V_{i}$; nous aurons construit ainsi une ligne fermée qui ne rencontrera les variétés $V_{i}$ qu'en un seul point. La somme des $S(M)$ ne pourra donc être nulle.

Supposons maintenant que les variétés $V_{i}$ partagent $U$ en plusieurs régions $R$; mais cependant qu'elles soient linéairement indépendantes. Dans ce cas, la frontière complète d'aucune des régions $R$ ne pourra être constituée par une partie des variétés $V_{i}$, sans quoi il y aurait entre ces variétés une homologie, ce que nous ne supposons pas.

La frontière complète de $R$ se composera donc de quelques-unes des
\origpage{41}
variétés $V_{i}$ et d'une partie de la frontière de $U$. Il en résulte que d'un point quelconque de $R$ on pourra aller à la frontière de $U$ sans sortir de $R$ et sans rencontrer aucune des variétés $V_{i}$.

Considérons alors une des variétés $V_{i}$, $V_{1}$, par exemple, et envisageons une petite ligne coupant $V_{1}$ en un point $M$; soient $R$ et $R'$ les régions auxquelles appartiennent les deux extrémités $\mu$ et $\mu'$ de cette petite ligne. De $\mu$ on pourra aller par une ligne continue à la frontière de $U$ sans sortir de $R$; de $\mu'$ on pourra aller par une troisième ligne continue à la frontière de $U$ sans sortir de $R'$.

L'ensemble de ces trois lignes continues ira de la frontière de $U$ à la frontière de $U$, et si je l'appelle $V$, on aura
\[
N(V, V_{1}) = 1, \qquad N(V, V_{i}) = 0 \qquad (i > 1).
\]

On peut donc choisir $V$ de telle façon que tous les $N(V, V_{i})$ soient nuls, excepté un d'entre eux qui sera égal à 1.

Supposons enfin qu'il y ait entre les $V_{i}$ un certain nombre d'homologies, 3, par exemple,
\[
\Sigma k_{i}V_{i} \backsim \Sigma k'_{i}V_{i} \backsim \Sigma k''_{i}V_{i} \backsim 0. \tag{$\alpha$}
\]

Nous pourrons alors, parmi les $k$ variétés $V_{i}$, en trouver $k-3$ qui soient linéairement indépendantes; soit, par exemple, $V_{1}, V_{2}, \ldots, V_{k-3}$.

Nous pouvons choisir $V$ de telle façon que tous les $N(V, V_{i})$, où $i = 1, 2, \ldots, k-3$, soient nuls à l'exception de l'un quelconque d'entre eux qui sera égal à 1. Les valeurs de
\[
N(V, V_{k-2}), \qquad N(V, V_{k-1}), \qquad N(V, V_{k})
\]
s'en déduiront à l'aide des relations
\[
\Sigma k_{i}N(V, V_{i}) = \Sigma k'_{i}N(V, V_{i}) = \Sigma k''_{i}N(V, V_{i}) = 0 \tag{$\beta$}
\]
qui sont une conséquence nécessaire de nos homologies ($\alpha$).
\origpage{42}

Il ne peut pas arriver alors que l'on ait
\[
\Sigma N(V, V_{i}) = 0, \tag{$\gamma$}
\]
quelle que soit celle des $k-3$ quantités
\[
N(V, V_{i}) \qquad (i = 1, 2, \ldots, k-3)
\]
qui soit égale à 1. Cela ne pourrait se faire, en effet, que si l'équation ($\gamma$) était une conséquence nécessaire des équations $\beta$; mais alors l'homologie
\[
\Sigma V_{i} \backsim 0 \tag{$\delta$}
\]
serait une conséquence nécessaire des homologies ($\alpha$). Or nous avons supposé au début que cette homologie ($\delta$) n'avait pas lieu.

On pourra donc toujours choisir $V$ de manière que l'égalité $\gamma$ n'ait pas lieu. C.~Q.~F.~D.

Nous désignerons sous le nom de \emph{coupure} toute variété contenue dans $U$~: 1° si elle est fermée; 2° si elle n'est pas fermée, mais si sa frontière fait partie de la frontière complète de $U$.

Nous pouvons alors énoncer le théorème suivant qui résume la discussion précédente~:

\emph{La condition nécessaire et suffisante} (si les variétés $V_{i}$ sont fermées et à $h-1$ dimensions) \emph{pour que l'on puisse choisir la coupure} $V$ \emph{de telle sorte que l'égalité}
\[
\Sigma k_{i}N(V, V_{i}) = 0
\]
\emph{n'ait pas lieu, c'est que l'homologie}
\[
\Sigma k_{i}V_{i} \backsim 0
\]
\emph{n'ait pas lieu}.
\origpage{43}

Cherchons maintenant à étendre ce théorème au cas où $V$ a $p$ dimensions et où $V_{1}, V_{2}, \ldots, V_{k}$ en ont $h-p$.

Soient
\[
F_{\alpha} = 0, \qquad \varphi_{\beta} > 0 \qquad (\alpha = 1, 2, \ldots, n-h)
\]
les équations de $U$; soient
\[
F_{\alpha} = 0, \qquad F'_{\gamma} = 0, \qquad \varphi_{\beta} > 0 \qquad (\gamma = 1, 2, \ldots, h-p)
\]
celles de $V$.

Quant à $V_{1}, V_{2}, \ldots, V_{k}$ nous les définirons de la manière suivante. Nous pourrons toujours trouver $p-1$ équations
\[
\Phi_{1} = \Phi_{2} = \ldots = \Phi_{p-1} = 0
\]
auxquelles satisfont tous les points de $V_{1}, V_{2}, \ldots, V_{k}$; c'est ainsi que si $h = 3$, $p = 2$, si $U$ est l'espace ordinaire et si $V_{1}, V_{2}, \ldots, V_{k}$ sont des courbes, on pourra toujours, par ces courbes, quelles qu'elles soient, faire passer une surface.

Pour définir $V_{i}$ nous y adjoindrons une $p^{\text{ième}}$ égalité
\[
F''_{i} = 0.
\]

Soit alors $U'$ la variété à $h-p+1$ dimensions
\[
F_{\alpha} = 0, \qquad \Phi_{\nu} = 0, \qquad \varphi_{\beta} > 0,
\]
et $V'$ la variété à 1 dimension
\[
F_{\alpha} = F'_{\gamma} = \Phi_{\nu} = 0, \qquad \varphi_{\beta} > 0.
\]

On aura d'abord
\[
N(V', V_{i}) = N(V, V_{i}).
\]

D'autre part, si l'homologie
\[
\Sigma k_{i}V_{i} \backsim 0
\]
\origpage{44}
a lieu par rapport à $U'$, elle aura lieu par rapport à $U$. Il est vrai que la réciproque n'est pas vraie et que cette homologie peut avoir lieu par rapport à $U$ sans avoir lieu par rapport à $U'$; mais, si elle a lieu par rapport à $U$, on pourra toujours, en choisissant convenablement les fonctions $\Phi$, trouver une variété $U'$ à $h-p+1$ dimensions par rapport à laquelle elle ait lieu.

\emph{Nous devons donc conclure que le théorème est encore vrai quand} $V$ \emph{a plus d'une dimension}.

Si donc $V_{1}$ et $V_{2}$ sont deux variétés fermées à $h-p$ dimensions, si pour toutes les coupures $V$ à $p$ dimensions, on a
\[
N(V, V_{1}) = N(V, V_{2}),
\]
on aura aussi
\[
V_{1} \backsim V_{2}
\]
et inversement.

Bornons-nous au cas où $U$ est fermée. Alors $U$ n'a pas de frontière et toutes les coupures sont fermées.

Le nombre des coupures de $p$ dimensions linéairement indépendantes est $P_{p} - 1$.

Soient
\[
C_{1}, \qquad C_{2}, \qquad \ldots, \qquad C_{\lambda} \qquad (\lambda = P_{p} - 1)
\]
ces coupures.

Soient ensuite
\[
V_{i} \qquad (i = 1, 2, \ldots, \mu)
\]
$\mu$ variétés fermées à $h-p$ dimensions.

La condition nécessaire et suffisante pour que l'on ait entre elles une homologie
\[
\Sigma k_{i}V_{i} \backsim 0,
\]
c'est que l'on ait~:
\[
\Sigma k_{i}N(C_{1}, V_{i}) = \Sigma k_{i}N(C_{2}, V_{i}) = \ldots = \Sigma k_{i}N(C_{\lambda}, V_{i}) = 0.
\]
\origpage{45}

Or, si le nombre $\mu$ des $V_{i}$ est supérieur à $\lambda$, on pourra toujours trouver des entiers $k_{i}$ satisfaisant à ces conditions, puisque nous aurons $\mu$ arbitraires $k_{i}$ et $\lambda$ conditions à remplir. Si donc les $V_{i}$ sont linéairement indépendantes, on aura
\[
\mu \leqq \lambda.
\]

Donc
\[
P_{h-p} \leqq P_{p};
\]
mais, en changeant $p$ en $h-p$, on trouve,
\[
P_{p} \leqq P_{h-p}.
\]

Donc
\[
P_{p} = P_{h-p}.
\]

\emph{Par conséquent, pour une variété fermée, les nombres de Betti également distants des extrêmes sont égaux}.

Ce théorème n'a, je crois, jamais été énoncé; il était cependant connu de plusieurs personnes qui en ont même fait des applications.

Voyons maintenant ce qui se passe quand $h$ est pair pour le nombre moyen $P_{\frac{h}{2}}$. Supposons que $h$ soit multiple de $4 + 2$, de telle façon que $\frac{h}{2}$ soit impair.

On sait que, quand, dans un déterminant, on fait un nombre impair de permutations de lignes, il change de signe. Quand on permutera $V$ et $V_{1}$, qui ont $\frac{h}{2}$ dimensions, le déterminant $f(M, M)$\ednote{Printed $f(M, M)$ with no prime; apparently for $f(M, M')$, as on pp. 34--37.} changera de signe; on aura donc
\[
N(V, V_{1}) = -N(V_{1}, V);
\]
on en déduit
\[
N(V, V) = 0.
\]

Le symbole $N(V, V)$ n'a, par lui-même, aucun sens, puisque, les deux variétés $V$ et $V$ coïncidant, le nombre des points d'intersection est infini;
\origpage{46}
il faut donc précisément\ednote{Printed ``précisément''; apparently for ``préciser''.} la définition; nous poserons donc
\[
N(V, V) = N(V, V'),
\]
où
\[
V \backsim V'.
\]

Cela posé, je dis que $P_{\frac{h}{2}}$ est impair. Supposons, en effet, qu'il soit pair et soient
\[
V_{1}, \qquad V_{2}, \qquad \ldots, \qquad V_{\mu}
\]
les $\mu$ variétés linéairement indépendantes à $h-1$\ednote{Printed $h-1$; apparently for $\frac{h}{2}$, the dimension of the varieties under discussion.} dimensions où
\[
\mu = P_{\frac{h}{2}} - 1
\]
est impair.

Formons le déterminant où le $i^{\text{ième}}$ terme de la $k^{\text{ième}}$ colonne est $N(V_{i}, V_{k})$. Ce déterminant serait gauche; c'est-à-dire que les termes de la diagonale principale seraient nuls et que deux termes symétriques par rapport à cette diagonale seraient égaux et de signe contraire. Comme le nombre des colonnes serait impair, ce déterminant serait nul. Il en résulterait, contrairement à l'hypothèse, que les variétés $V_{1}, V_{2}, \ldots, V_{\mu}$ ne seraient pas linéairement indépendantes.

Donc $P_{\frac{h}{2}}$ est impair.

Cela ne serait plus vrai ni si $h$ était multiple de 4, ni si la variété $U$ était unilatère; car tous nos raisonnements supposent les variétés bilatères. Nous en verrons plus loin des exemples.

\subsection*{\S~10. --- Représentation géométrique.}

Il y a une manière de se représenter les variétés à trois dimensions situées dans l'espace à quatre dimensions, manière qui en facilite singulièrement l'étude.

Considérons, dans l'espace ordinaire, un certain nombre de polyèdres
\origpage{47}
\[
P_{1}, \qquad P_{2}, \qquad \ldots, \qquad P_{n}.
\]

Nous pouvons supposer qu'il existe, dans l'espace à quatre dimensions, un certain nombre de variétés à trois dimensions
\[
Q_{1}, \qquad Q_{2}, \qquad \ldots, \qquad Q_{n}
\]
respectivement homéomorphes à $P_{1}, P_{2}, \ldots, P_{n}$.

Soient $F_{1}$ une face du polyèdre $P_{1}$ et $\Phi_{1}$ l'ensemble des points de la frontière de $Q_{1}$ qui correspondent aux divers points de $F_{1}$. Soient de même $F_{2}$ une face de $P_{2}$ et $\Phi_{2}$ l'\emph{image} de cette face sur la frontière de $Q_{2}$.

Il peut se faire que $\Phi_{1}$ coïncide avec $\Phi_{2}$. Dans ce cas les deux variétés $Q_{1}$ et $Q_{2}$ sont contiguës, et l'on passe de l'intérieur de l'une à l'intérieur de l'autre en franchissant $\Phi_{1}$.

Dans ce cas, nous dirons que les faces $F_{1}$ et $F_{2}$ sont \emph{conjuguées}.

Il pourra arriver que les faces $F_{1}$ et $F_{2}$ appartiennent à un \emph{même} polyèdre $P_{1}$. Alors la variété à deux dimensions $\Phi_{1}$, qui ne diffère pas de la variété à deux dimensions $\Phi_{2}$, séparera l'une de l'autre deux portions de la variété $Q_{1}$.

On le comprendra mieux si je compare avec un exemple se passant dans l'espace ordinaire. Considérons un rectangle $ABCD$ et un tore, sur lequel nous tracerons deux coupures, à savoir~: une circonférence méridienne et un parallèle; soit $H$ le point d'intersection de ce parallèle et de ce méridien. La surface du tore ainsi découpée sera homéomorphe au rectangle; les deux lèvres de la coupure formée par le méridien correspondront aux deux côtés $AB$ et $CD$ du rectangle; les deux lèvres de la coupure formée par le parallèle correspondront aux deux autres côtés $AD$ et $BC$. On voit l'analogie avec ce qui précède~: le rectangle correspond au polyèdre $P_{1}$, le tore découpé à la variété $Q_{1}$, les côtés $AB$ et $CD$ aux deux faces $F_{1}$ et $F_{2}$, les deux lèvres de la coupure méridienne aux deux variétés $\Phi_{1}$ et $\Phi_{2}$ qui, comme on le voit, coïncident.

Cela posé, imaginons que parmi les faces des $n$ polyèdres $P_{i}$ il y en ait
\origpage{48}
un certain nombre qui soient conjuguées deux à deux, et un certain nombre qui restent \emph{libres}.

Considérons la variété totale $V$ formée par l'ensemble des variétés $Q_{i}$. Comme quelques-unes de ces variétés $Q_{i}$ sont contiguës, il pourra se faire que la variété totale $V$ soit continue~: c'est ce que je supposerai.

Si parmi les faces des $P_{i}$ il n'y en a aucune qui reste \emph{libre}, la variété $V$ sera fermée. Dans le cas contraire, les points qui correspondront aux faces restées libres formeront la frontière complète de $V$.

On conçoit que la connaissance des polyèdres $P_{i}$ et celle du mode de conjugaison de leurs faces nous fournissent, dans l'espace ordinaire, une image de la variété $V$ et que cette image suffise pour l'étude de ses propriétés au point de vue de l'\emph{Analysis Situs}.

Voyons comment doit être défini le mode de conjugaison des faces. Il est clair d'abord que, pour que deux faces puissent être conjuguées, il faut qu'elles aient le même nombre de côtés. Ensuite, pour connaître complètement le mode de conjugaison, il ne suffit pas de savoir que telle face est conjuguée de telle autre; il faut savoir, en outre, quel est le sommet ou le côté de cette face qui correspond à tel sommet ou à tel côté de sa conjuguée.

Alors seulement le mode de conjugaison sera complètement défini.

A deux faces conjuguées correspond, par définition, une même variété à deux dimensions intérieure à $V$. Il peut se faire aussi qu'à plusieurs arêtes des polyèdres $P$ corresponde une même ligne intérieure à $V$, ou qu'à plusieurs sommets de ces polyèdres corresponde un même point intérieur à $V$. Nous dirons alors que ces arêtes ou ces sommets appartiennent à un même cycle.

Voici comment on peut former les cycles d'arêtes et les cycles de sommets.

Soit $A_{1}$ une arête, $F'_{1}$ l'une des deux faces qui passent par $A_{1}$, $F_{2}$ la conjuguée de $F'_{1}$, $A_{2}$ l'arête de $F_{2}$ qui correspond à $A_{1}$; $F'_{2}$ l'autre face qui passe par $A_{2}$; $F_{3}$ la conjuguée de $F'_{2}$, $A_{3}$ l'arête correspondant à $A_{2}$, etc.

On s'arrêtera quand on arrivera à une face libre ou quand on sera
\origpage{49}
revenu à l'arête $A_{1}$. Toutes les arêtes $A_{1}$, $A_{2}$, $A_{3}$, \ldots\ formeront un cycle.

De même pour les sommets. Soient $S_{1}$ un sommet, $F'_{1}$ l'une des faces qui passent par $S_{1}$, $F_{2}$ la conjuguée de $F'_{1}$; $S_{2}$ le sommet correspondant à $S_{1}$; $F'_{2}$ \emph{une} des faces qui passent par $S_{2}$, \ldots.

$S_{1}$, $S_{2}$, $S_{3}$, \ldots\ appartiendront à un même cycle.

Seulement ici il passe par chaque sommet plus de deux faces; alors $F'_{2}$, par exemple, peut être choisie de plusieurs manières; on ne devra s'arrêter qu'après avoir épuisé tous les choix possibles.

L'analogie avec la formation des cycles dans la théorie des groupes fuchsiens est évidente. Elle le devient plus encore, si l'on suppose qu'il n'y a qu'un seul polyèdre $P_{1}$.

\subsection*{Premier exemple.}

L'exemple le plus simple est celui où l'on n'a qu'un seul polyèdre $P_{1}$ et où ce polyèdre est un cube $ABCDA'B'C'D'$ dont les sommets ont respectivement pour coordonnées
\[ \begin{matrix} A & 0 & 0 & 0 & \qquad & A' & 0 & 0 & 1 \\[1ex] B & 0 & 1 & 0 & \qquad & B' & 0 & 1 & 1 \\[1ex] C & 1 & 0 & 0 & \qquad & C' & 1 & 0 & 1 \\[1ex] D & 1 & 1 & 0 & \qquad & D' & 1 & 1 & 1 \end{matrix} \]

Je suppose que les faces opposées soient conjuguées et cela de la façon suivante~:
\[ \left\{ \begin{aligned} &ABDC \equiv A'B'D'C',\\ &ACC'A' \equiv BDD'A',\\ &CDD'C' \equiv ABB'A'. \end{aligned} \right. \tag{1} \]
\ednote{In the first row only broken fragments of the sign $\equiv$ printed; the sign is read as $\equiv$, as in the other rows. The second row prints $BDD'A'$, which is not a face of the cube; apparently for $BDD'B'$, the face opposite $ACC'A'$, as in the restatement of this example on p. 65.}

Voici ce que j'entends par cette notation~: la relation
\[ ABDC \equiv A'B'D'C' \]
signifie~:
\origpage{50}

1° Que les faces $ABDC$ et $A'B'D'C'$ sont conjuguées;

2° Que les sommets se rencontrant sur la première de ces faces se rencontrent dans l'ordre circulaire $ABDC$;

3° Que les sommets $A$ et $A'$, $B$ et $B'$, $D$ et $D'$, $C$ et $C'$ se correspondent.

Observons en passant que les sommets $ABDC$ se rencontrent sur la première de ces faces dans un ordre circulaire tel qu'en parcourant le périmètre dans le sens correspondant on ait la face à sa gauche. Au contraire, les sommets correspondants $A'B'D'C'$ se rencontrent sur la seconde face dans un ordre tel qu'en parcourant le périmètre dans ce sens on ait la face à sa droite.

Cette condition doit toujours être remplie \emph{si l'on veut que $V$ soit bilatère}. Si deux faces $F_{1}$ et $F_{2}$ sont conjuguées et si un point décrit le périmètre de $F_{1}$ de façon à laisser la face à sa gauche, le point correspondant sur $F_{2}$ devra décrire le périmètre de cette face en la laissant à sa droite.

Cela posé, supposons donc que le mode de conjugaison soit défini par les relations (1). Il est aisé de voir alors que tous les sommets forment un seul cycle et qu'il y a trois cycles d'arêtes, comprenant respectivement celles des arêtes qui sont parallèles à l'axe des $x$, à celui des $y$ et à celui des $z$.

\subsection*{Deuxième exemple.}

Conservons toujours notre cube, mais changeons le mode de conjugaison et définissons-le par les relations
\[ \left\{ \begin{aligned} &ABDC \equiv B'D'C'A',\\ &ABB'A' \equiv DD'C'C,\\ &ACC'A' \equiv DD'B'B; \end{aligned} \right. \tag{2} \]
il y aura alors deux cycles d'arêtes et deux cycles de sommets et je résume les résultats par les relations suivantes; j'ai mis le signe $\equiv$ entre deux arêtes (ou deux sommets) pour exprimer qu'elles font partie d'un même cycle.
\origpage{51}

Deux cycles d'arêtes~:
\[ \begin{aligned} &AB \equiv B'D' \equiv C'C \equiv B'A' \equiv AC \equiv DD',\\ &AA' \equiv DC \equiv C'A' \equiv B'B \equiv C'D' \equiv DB. \end{aligned} \]

Deux cycles de sommets~:
\[ \begin{aligned} &A \equiv B' \equiv C' \equiv D,\\ &B \equiv D' \equiv C \equiv A'. \end{aligned} \]
\ednote{In the second row only two strokes of the last sign $\equiv$, between $C$ and $A'$, printed; it is read as $\equiv$, as elsewhere in these cycles.}

Nous verrons plus tard pour quelles raisons ce mode de conjugaison est inadmissible.

\subsection*{Troisième exemple.}

Conservons notre cube avec le mode de conjugaison suivant~:
\[ \left\{ \begin{aligned} &ABDC \equiv B'D'C'A',\\ &ABB'A' \equiv DD'C'C,\\ &ACC'A' \equiv DD'B'B. \end{aligned} \right. \tag{3} \]
\ednote{As printed, system (3) repeats system (2) letter for letter, yet the cycles stated below differ from those of the second example. The stated cycles fit a second row $ABB'A' \equiv C'CDD'$; the printed $DD'C'C$ is apparently a misprint.}

On trouve alors~:

Quatre cycles d'arêtes~:
\[ \begin{aligned} &AB \equiv B'D' \equiv C'A', &&AA' \equiv C'D' \equiv DB,\\ &AC \equiv DD' \equiv B'A', &&CD \equiv BB' \equiv A'C'. \end{aligned} \]
\ednote{In the first row only two strokes of the sign $\equiv$ between $AA'$ and $C'D'$ printed; it is read as $\equiv$. The first cycle prints $C'A'$, so that $A'C'$ occurs twice and $CC'$ never; apparently for $C'C$.}

Deux cycles de sommets~:
\[ \begin{aligned} &A \equiv B' \equiv C' \equiv D,\\ &B \equiv D' \equiv A' \equiv C. \end{aligned} \]
\ednote{In the first row only two strokes of the sign $\equiv$ between $C'$ and $D$ printed; it is read as $\equiv$.}

\subsection*{Quatrième exemple.}

Soit maintenant
\[ \left\{ \begin{aligned} &ABDC \equiv B'D'C'A',\\ &ABB'A' \equiv CDD'C',\\ &ACC'A' \equiv BDD'B'. \end{aligned} \right. \tag{5} \]
\ednote{The print numbers this system (5), directly after (3); no equation (4) precedes it.}
\origpage{52}

On trouve~:

Trois cycles d'arêtes~:
\[ \begin{aligned} &AA' \equiv CC' \equiv BB' \equiv DD',\\ &AB \equiv CD \equiv B'D' \equiv A'C',\\ &AC \equiv BD \equiv D'C' \equiv B'A', \end{aligned} \]
et un seul cycle de sommets.

\subsection*{Cinquième exemple.}

Soit un octaèdre régulier; il aura six sommets dont quatre formeront un carré $BCED$ (les lettres sont rangées dans l'ordre circulaire dans lequel on rencontre les quatre sommets en parcourant le périmètre du carré) et dont deux, $A$ et $P$\ednote{Printed $P$; apparently for $F$, the vertex named in the system (6) below.}, sont en dehors du carré. Soit
\[ \left\{ \begin{aligned} &ABC \equiv FED,\\ &ACE \equiv FDB,\\ &AED \equiv FBC,\\ &ADB \equiv FCE \end{aligned} \right. \tag{6} \]
le mode de conjugaison; on trouvera six cycles d'arêtes et trois cycles de sommets; chaque arête formera un cycle avec l'arête opposée, c'est-à-dire avec sa symétrique par rapport au centre de symétrie de l'octaèdre, et chaque sommet formant un cycle avec le sommet opposé.

Il est inutile de multiplier davantage les exemples; je me propose de reconnaître maintenant si tous les modes de conjugaison sont admissibles.

Appelons \emph{aster} la figure formée par un certain nombre d'angles solides ayant même sommet et disposés autour de ce sommet de telle sorte que tout point de l'espace appartienne à un et à un seul de ces angles solides.

Dans la figure, on distinguera donc, outre les angles solides eux-mêmes et leur sommet, un certain nombre de faces communes à deux angles solides et un certain nombre d'arêtes communes à plusieurs angles solides.
\origpage{53}

Je pourrai supposer que les arêtes et les faces sont prolongées indéfiniment, ou bien qu'on les arrête à la surface d'une sphère ayant son centre au sommet commun des angles solides. Alors ces divers angles solides vont découper sur la sphère un certain nombre de polygones sphériques, de sorte que la surface de la sphère sera subdivisée en polygones sphériques. Cette surface ainsi subdivisée pourra être regardée comme homéomorphe à un polyèdre convexe dont les faces correspondront aux polygones sphériques que je viens de définir.

Alors aux angles solides de l'aster correspondront les faces du polyèdre; aux faces de l'aster correspondront les arêtes du polyèdre et, aux arêtes de l'aster, ses sommets.

Soient $S$, $F$ et $A$ le nombre des angles solides, des faces et des arêtes de l'aster. Comme le polyèdre doit satisfaire au théorème d'Euler, on devra avoir
\[ S - F + A = 2. \]

Revenons maintenant à nos polyèdres
\[ P_{1}, \qquad P_{2}, \qquad \ldots, \qquad P_{n} \]
et envisageons un cycle de sommets; à tous les sommets de ce cycle correspondra un même point de $V$ que j'appelle $M$. Parmi les variétés
\[ Q_{1}, \qquad Q_{2}, \qquad \ldots, \qquad Q_{n}, \]
il y en aura un certain nombre qui auront le point $M$ sur leur frontière; je les appellerai les variétés $Q_{\alpha}$ et ce seront celles qui correspondront à des polyèdres $P_{\alpha}$ auxquels appartiendront les divers sommets du cycle.

Considérons maintenant deux des variétés $Q_{\alpha}$ qui soient contiguës, et la variété à deux dimensions qui leur sert de frontière commune. Les variétés à deux dimensions ainsi définies, je les appellerai les $\Phi_{\alpha}$; elles correspondront à celles des faces des polyèdres $P_{\alpha}$ auxquelles appartiennent les divers sommets du cycle.

Nous envisagerons enfin les variétés à une dimension qui sont l'intersection
\origpage{54}
de deux variétés $\Phi_{\alpha}$ et que j'appellerai les $L_{\alpha}$. Elles correspondront à celles des arêtes des polyèdres $P_{\alpha}$ auxquelles appartiennent les divers sommets du cycle.

Envisageons la figure formée par les variétés $Q_{\alpha}$, $\Phi_{\alpha}$, $L_{\alpha}$ ou plutôt par les points de ces variétés qui satisfont à l'inégalité
\[ (x_{1} - x^{0}_{1})^{2} + (x_{2} - x^{0}_{2})^{2} + (x_{3} - x^{0}_{3})^{2} + (x_{4} - x^{0}_{4})^{2} < \varepsilon^{2}. \tag{7} \]

$\varepsilon$ est une très petite quantité; $x^{0}_{1}$, $x^{0}_{2}$, $x^{0}_{3}$, $x^{0}_{4}$ sont les coordonnées du point $M$.

\emph{Cette figure est évidemment homéomorphe à un aster} dont les faces et les arêtes seraient limitées à une sphère. Soit $A$ cet aster.

Considérons une quelconque des variétés $Q_{\alpha}$ et, en outre, la variété $W$ formée par ceux de ses points qui satisfont à l'inégalité (7). Cette variété sera continue ou ne le sera pas. Si elle ne l'est pas, je la décomposerai en plusieurs variétés continues, et les diverses variétés continues que l'on pourra ainsi former, je les appellerai les $Q'_{\alpha}$. (Le nombre des $Q'_{\alpha}$ peut ainsi être plus grand que celui des $Q_{\alpha}$ si l'une des variétés $W$ n'est pas continue.)

Je définirai de même les $\Phi'_{\alpha}$ et les $L'_{\alpha}$.

Soient $q_{\alpha}$, $\varphi_{\alpha}$, $l_{\alpha}$ le nombre des $Q'_{\alpha}$, des $\Phi'_{\alpha}$ et des $L'_{\alpha}$; ce sera aussi celui des angles solides, des faces et des arêtes de l'aster $A$.

D'où cette conclusion~:

\emph{Pour qu'un mode de conjugaison soit admissible, il faut que, pour tous les cycles de sommets, on ait}
\[ q_{\alpha} - \varphi_{\alpha} + l_{\alpha} = 2. \]

Voyons maintenant comment on pourra calculer les nombres $q_{\alpha}$, $l_{\alpha}$\ednote{Printed $l_{\alpha}$; apparently for $\varphi_{\alpha}$, since the three numbers treated in 1°, 2°, 3° below are $q_{\alpha}$, $\varphi_{\alpha}$, $l_{\alpha}$.}, $l_{\alpha}$ pour un cycle de sommets quelconque $\alpha$.

1° $q_{\alpha}$ sera le nombre des sommets du cycle.

2° Pour avoir $\varphi_{\alpha}$, on n'a qu'à faire la somme des nombres des faces qui vont passer par les divers sommets du cycle et à diviser par deux. Si,
\origpage{55}
par exemple, le cycle $\alpha$ se compose de trois sommets $a$, $b$, $c$ appartenant respectivement aux polyèdres $P_{1}$, $P_{2}$ et $P_{3}$; si le point $a$ est le sommet d'un angle trièdre de sorte que trois faces de $P_{1}$ aillent passer par $a$; si quatre faces de $P_{2}$ vont passer par $b$ et cinq faces de $P_{3}$ par $c$, on aura
\[ \varphi_{\alpha} = \frac{3 + 4 + 5}{2} = 6. \]

3° Pour avoir $l_{\alpha}$, on énumérera les arêtes qui aboutissent aux divers sommets du cycle de la manière suivante~: toutes les arêtes d'un même cycle d'arêtes ne compteront que pour une seule, si l'une des extrémités appartient au cycle $\alpha$; elles compteront pour deux si les deux extrémités appartiennent au cycle $\alpha$.

Appliquons ces règles aux six exemples\ednote{Printed "six"; the Tableau that follows lists five examples, and the fifth example of the preceding pages is the last one treated there.} dont nous nous sommes occupés plus haut, nous aboutirons au Tableau suivant~:
\[ \begin{matrix} \text{Exemple.} & q_{\alpha}. & \varphi_{\alpha}. & l_{\alpha}. \\[1ex] \text{Premier}\ldots\ldots & 8 & 12 & 6 \\[1ex] \text{Deuxième}\ldots\ldots & 4 & 6 & 2 \\[1ex] \text{Troisième}\ldots\ldots & 4 & 6 & 4 \\[1ex] \text{Quatrième}\ldots\ldots & 8 & 12 & 6 \\[1ex] \text{Cinquième}\ldots\ldots & 2 & 4 & 4 \end{matrix} \]

Il est à remarquer que, dans les exemples 2, 3 et 5, il y a plusieurs cycles de sommets, mais il arrive que les trois nombres $q_{\alpha}$, $\varphi_{\alpha}$ et $l_{\alpha}$ ont mêmes valeurs pour tous les cycles d'un même exemple.

Ce Tableau montre que la relation
\[ q_{\alpha} - \varphi_{\alpha} + l_{\alpha} = 2 \]
est satisfaite pour tous nos exemples, sauf pour le second. Le mode de conjugaison du second exemple doit donc être rejeté.

\subsection*{\S~11. --- Représentation par un groupe discontinu.}

Voici un autre mode de représentation qui peut aussi être employé dans certains cas.
\origpage{56}

Soit $(x, y, z)$ un point de l'espace ordinaire; considérons une série de substitutions qui changent respectivement $x$, $y$ et $z$ en
\[ \begin{matrix} \varphi_{1}(x, y, z), & \psi_{1}(x, y, z), & \chi_{1}(x, y, z), \\[1ex] \varphi_{2}(x, y, z), & \ldots\ldots, & \ldots\ldots, \\[1ex] \ldots\ldots, & \ldots\ldots, & \ldots\ldots, \\[1ex] \varphi_{n}(x, y, z), & \psi_{n}(x, y, z), & \chi_{n}(x, y, z), \\[1ex] \ldots\ldots, & \ldots\ldots, & \ldots\ldots \end{matrix} \]

Supposons que l'ensemble de ces substitutions forme un groupe proprement discontinu. L'espace va se trouver partagé en une infinité de domaines
\[ D_{0}, \qquad D_{1}, \qquad D_{2}, \qquad \ldots, \]
et cela de telle sorte qu'à chaque domaine $D_{i}$ corresponde une substitution du groupe $S_{i}$ qui change $D_{0}$ en $D_{i}$.

Considérons une surface $\Sigma$ formant une portion de la frontière de $D_{0}$ séparant ce domaine $D_{0}$ du domaine $D_{i}$; la substitution $S^{-1}_{i}$ inverse de $S_{i}$ change $D_{i}$ en $D_{0}$ et, comme les points de $\Sigma$ appartiennent à la frontière de $D_{i}$, la transformée de la surface $\Sigma$ sera une autre partie de la frontière de $D_{0}$.

La frontière de $D_{0}$ va ainsi se trouver partagée en portions de surface qui seront conjuguées deux à deux, de telle façon que chacune d'elles sera la transformée de sa conjuguée par une des substitutions du groupe.

Le domaine $D_{0}$, avec sa frontière ainsi subdivisée, sera homéomorphe à un polyèdre dont les faces seront conjuguées deux à deux, comme dans le numéro précédent. On pourra alors, comme dans le numéro précédent, faire correspondre à ce polyèdre ou, par conséquent, au domaine $D_{0}$, une variété fermée à trois dimensions située dans l'espace à quatre dimensions, et que l'on obtiendra en transportant $D_{0}$ dans l'espace à quatre dimensions, puis en déformant ce domaine jusqu'à ce qu'on puisse \emph{coller} l'une contre l'autre les portions conjuguées de sa frontière.

L'analogie avec la théorie des groupes fuchsiens est trop évidente pour qu'il soit nécessaire d'insister; je me bornerai à un seul exemple~:
\origpage{57}

\subsection*{Sixième exemple.}

Considérons le groupe dérivé des trois substitutions
\[ \left\{ \begin{aligned} &(x, y, z; x + 1, y, z),\\ &(x, y, z; x, y + 1, z),\\ &(x, y, z; \alpha x + \beta y, \gamma x + \delta y, z + 1), \end{aligned} \right. \tag{1} \]
$\alpha$, $\beta$, $\gamma$, $\delta$ étant des entiers tels que
\[ \alpha\delta - \beta\gamma = 1. \]

J'appellerai ce groupe le groupe $(\alpha, \beta, \gamma, \delta)$.

Je dis qu'il est proprement discontinu.

Pour nous en rendre compte, cherchons comment sont distribués dans l'espace les transformés d'un même point
\[ x = a, \qquad y = b, \qquad z = c. \]

D'abord, tous les transformés de ce point, par une combinaison quelconque des deux premières substitutions, seront compris dans la formule
\[ x = a + k, \qquad y = b + k_{1}, \qquad z = c, \tag{2} \]
$k$ et $k_{1}$ étant deux entiers; il est facile de voir d'ailleurs que ces deux substitutions sont permutables.

Transformons maintenant l'ensemble des points (2) par la troisième substitution; nous trouverons
\[ x = \alpha(a + k) + \beta(b + k_{1}), \qquad y = \gamma(a + k) + \delta(b + k_{1}), \qquad z = c + 1. \tag{3} \]

Si nous posons, pour abréger,
\[ \alpha a + \beta b = a_{1}, \qquad \gamma a + \delta b = b_{1}, \]
de sorte que le point $(a_{1}, b_{1})$ soit le transformé du point $(a, b)$ par la substitution
\origpage{58}
\[ (x, y; \alpha x + \beta y, \gamma x + \delta y), \]
que j'appellerai $s$.\ednote{In the display the first $x$ after $\alpha$ prints like an $r$; apparently damaged type or a misprint for $x$.}

Nous pourrons alors remplacer les équations (3) par les suivantes~:
\[ x = a + k', \qquad y = b_{1} + k'_{1}, \qquad z = c + 1, \tag{4} \]
$k'$ et $k'_{1}$ étant deux entiers nouveaux; en appliquant à ces points les deux premières substitutions, on retombe donc toujours sur les mêmes points.\ednote{In (4) the print has $x = a + k'$ with no subscript; apparently for $x = a_{1} + k'$, either a misprint or a subscript that did not print.}

Appliquons-leur la troisième substitution.

Soit
\[ a_{2} = \alpha a_{1} + \beta b_{1}, \qquad b_{2} = \gamma a_{1} + \delta b_{1}, \]
de telle sorte que le point $a_{2}$, $b_{2}$ soit le transformé du point $a_{1}$, $b_{1}$ par $s$.

Alors les transformés des points (4) par la troisième substitution seront compris dans la formule
\[ x = a_{2} + k'', \qquad y = b_{2} + k''_{1}, \qquad z = c + 2, \tag{5} \]
$k''$ et $k''_{1}$ étant deux entiers.

En général, supposons que les transformés successifs du point $(a, b)$ par la substitution $s$ soient $a_{1}, b_{1}$; $a_{2}, b_{2}$; \ldots; $a_{n}, b_{n}$; \ldots\ et de même que les transformés successifs par la substitution inverse soient $a_{-1}$, $b_{-1}$; $a_{-2}$, $b_{+2}$\ednote{Printed $b_{+2}$; apparently a misprint for $b_{-2}$.}; \ldots.

Alors tous les transformés du point $(a, b, c)$ par les substitutions du groupe (1) seront donnés par la formule
\[ x = a_{n} + k, \qquad y = b_{n} + k_{1}, \qquad z = c + n, \]
où $n$, $k$ et $k_{1}$ sont trois entiers arbitraires.

On voit aisément d'ailleurs que le groupe dérivé des deux premières substitutions est permutable à la troisième.

Le domaine fondamental $D_{0}$ est un cube dont le côté est égal à 1 et limité par les six plans
\[ x; y; z = 0; 1. \]
\origpage{59}

Le cas le plus simple est celui où
\[ \alpha = \delta = 1, \qquad \beta = \gamma = 0, \]
parce qu'alors nos trois substitutions se réduisent à
\[ (x, y, z; x + 1, y, z; x, y + 1, z; x, y, z + 1). \]

Chacune d'elles change une des faces du cube en la face opposée; nous retombons donc tout simplement sur notre premier exemple.

Mais la manière dont la superficie du cube $D_{0}$ se subdivise en régions conjuguées n'est pas toujours aussi simple.

Supposons, par exemple,
\[ \alpha = \beta = \delta = 1, \qquad \gamma = 0; \]
chacune des faces parallèles à l'axe des $z$ sera encore conjuguée de la face opposée; mais pour les faces $z = 0$, $z = 1$ perpendiculaires à l'axe des $z$, cela sera plus compliqué.

Supposons que les points $ABCD$, $A'B'C'D'$ aient mêmes coordonnées que dans nos quatre premiers exemples. Chacune des faces $ABCD$, $A'B'C'D'$ devra être partagée en deux triangles, à savoir~: $ABC$ et $BCD$ d'une part, $A'D'C'$ et $A'B'D'$ d'autre part, et la loi de conjugaison des faces sera exprimée par les relations
\[ \begin{aligned} &ACC'A' \equiv BDD'A',\\ &CDD'C' \equiv ABB'A',\\ &ABC \equiv A'D'C',\\ &BCD \equiv B'A'D'. \end{aligned} \]
\ednote{The first relation prints $BDD'A'$; apparently for $BDD'B'$, the face opposite $ACC'A'$.}

Plus généralement les faces parallèles à l'axe des $z$ resteront conjuguées deux à deux, mais les faces perpendiculaires à l'axe des $z$ devront être décomposées en polygones plus petits, d'autant plus nombreux que les nombres $\alpha$, $\beta$, $\gamma$, $\delta$ seront plus grands et qui seront conjugués deux à deux suivant une loi plus ou moins compliquée.
\origpage{60}

Un cas simple est celui-ci~:
\[ \alpha = \delta = 0, \qquad \beta = 1, \qquad \gamma = -1. \]

Dans ce cas, le mode de conjugaison est le même que dans notre quatrième exemple.

\subsection*{\S~12. --- Groupe fondamental.}

Nous sommes ainsi conduit à la notion du groupe fondamental d'une variété.

Soient
\[ F_{1}, \qquad F_{2}, \qquad \ldots, \qquad F_{\lambda}, \]
$\lambda$ fonctions des coordonnées $x_{1}, x_{2}, \ldots, x_{n}$ d'un point $M$ d'une variété $V$ définie par les relations
\[ f_{\alpha} = 0, \qquad \varphi_{\beta} > 0. \]

Je ne suppose pas que les fonctions $F$ soient uniformes. Lorsque le point $M$, partant de sa position initiale $M_{0}$, reviendra à cette position, après avoir parcouru un chemin quelconque, il pourra se faire que les fonctions $F$ ne reviennent pas à leurs valeurs primitives.

Mais pour mieux fixer les idées et bien que cela n'ait rien d'essentiel, supposons que les fonctions $F$ soient définies de la manière suivante. Elles devront satisfaire à des équations différentielles de la forme
\[ dF_{i} = X_{i.1}\,dx_{1} + X_{i.2}\,dx_{2} + \ldots + X_{i.n}\,dx_{n}, \tag{1} \]
où les coefficients $X_{i.k}$ seront des fonctions données des $x_{k}$ et des $F_{i}$. Ces fonctions devront être uniformes, finies et continues, ainsi que leurs dérivées pour toutes les valeurs des $F$ ainsi que pour tous les points de $V$ et même pour tous les points suffisamment voisins de $V$.

Je suppose également que, pour tous les points suffisamment voisins de $V$, les équations (1) satisfassent aux conditions d'intégrabilité, qui
\origpage{61}
peuvent s'écrire
\[ \begin{aligned} &\frac{dX_{i.k}}{dx_{q}} + \frac{dX_{i.k}}{dF_{1}}X_{1.q} + \frac{dX_{i.k}}{dF_{2}}X_{2.q} + \ldots + \frac{dX_{i.k}}{dF_{\lambda}}X_{\lambda.k}\\ &= \frac{dX_{i.q}}{dx_{k}} + \frac{dX_{i.q}}{dF_{1}}X_{1.k} + \frac{dX_{i.q}}{dF_{2}}X_{2.k} + \ldots + \frac{dX_{i.q}}{dF_{k}}X_{\lambda.k}. \end{aligned} \]\ednote{The last term of the first line is printed $X_{\lambda.k}$, apparently for $X_{\lambda.q}$; the last denominator of the second line is printed $dF_{k}$, apparently for $dF_{\lambda}$. Both are reproduced as printed.}

Si alors le point $M$ décrit sur la variété $V$ un contour infiniment petit, les fonctions $F$ reviendront à leurs valeurs primitives. Il en sera encore de même, si le point $M$ décrit sur la variété $V$ un \emph{lacet}, c'est-à-dire s'il va d'abord de $M_{0}$ en $M_{1}$ par un chemin quelconque $M_{0}BM_{1}$, s'il décrit ensuite un contour infiniment petit, et si enfin il revient de $M_{1}$ à $M_{0}$ par le \emph{même} chemin $M_{1}BM_{0}$.

Mais il pourra n'en plus être de même s'il décrit un contour fermé fini.

Supposons, par exemple, que nous soyons dans l'espace ordinaire et que notre variété $V$ soit un tore. On pourra évidemment imaginer que les fonctions $F$ reviennent à leurs valeurs primitives, quand le point $M$ décrira un lacet sur ce tore, mais qu'il n'en soit plus de même si le point $M$ décrit une circonférence méridienne ou un parallèle.

Les valeurs finales des fonctions $F$, quand le mobile $M$, partant d'un point initial $M_{0}$, y reviendra après avoir décrit un contour fermé; ces valeurs finales, dis-je, dépendront naturellement des valeurs initiales.

Soient donc $F^{0}_{a}$ et $F^{1}_{a}$ les valeurs initiale et finale de $F_{a}$. Les $F^{1}_{a}$ seront des fonctions des $F^{0}_{a}$, ou, en d'autres termes, les fonctions $F$ subiront, quand le point $M$ décrira le contour fermé considéré, une certaine substitution, qui changera les $F^{0}_{a}$ en $F^{1}_{a}$.

\emph{L'ensemble de toutes les substitutions que les fonctions $F$ subiront ainsi, quand le point $M$ décrira tous les contours fermés que l'on peut tracer sur la variété $V$ en partant du point initial $M_{0}$, formera évidemment un groupe que j'appelle $g$}.

Maintenant considérons un contour fermé $M_{0}BM_{0}$ tracé sur $V$ et partant du point initial $M_{0}$. Si ce contour fermé se réduit à un lacet, je conviendrai d'écrire
\origpage{62}
\[ M_{0}BM_{0} \equiv 0. \]\ednote{Only two strokes of the sign are printed here, so it reads as $=$; the same relation is written $M_{0}BM_{0} \equiv 0$ below and on p. 63, and a stroke of the $\equiv$ apparently did not print.}

Si maintenant $M_{0}AM_{1}$, $M_{0}BM_{1}$, $M_{0}CM_{1}$ sont trois chemins différents tracés sur $V$ et allant de $M_{0}$ à $M_{1}$, je conviendrai d'écrire
\[ M_{0}AM_{1}CM_{0} \equiv M_{0}AM_{1}BM_{0} + M_{0}BM_{1}CM_{0}. \]

Il importe de remarquer que $M_{0}AM_{1}CM_{0}$ n'est pas la même chose que $M_{0}CM_{1}AM_{0}$, ni $M_{0}AM_{1}BM_{0} + M_{0}BM_{1}CM_{0}$ la même chose que
\[ M_{0}BM_{1}CM_{0} + M_{0}AM_{1}BM_{0}; \]
on ne peut pas intervertir l'ordre des termes d'une somme.

Il résulte de cette convention que l'on aura
\[ M_{0}BM_{0} \equiv 0, \]
si le contour fermé $M_{0}BM_{0}$ constitue la frontière complète d'une variété et\ednote{Printed ``et''; apparently for ``à'', a variety à deux dimensions.} deux dimensions faisant partie de $V$; et, en effet, ce contour fermé pourra alors être décomposé en un nombre très grand de lacets.

Nous sommes ainsi conduit à envisager des relations de la forme
\[ k_{1}C_{1} + k_{2}C_{2} \equiv k_{3}C_{3} + k_{4}C_{4}, \]
où les $k$ sont des entiers et les $C$ des contours fermés tracés sur $V$ et partant de $M_{0}$. Ces relations, que j'appellerai des \emph{équivalences}, ressemblent aux homologies que j'ai étudiées plus haut. Elles en diffèrent~:

1° Parce que, dans les homologies, les contours peuvent partir d'un point initial quelconque.

2° Parce que, dans les homologies, on a le droit d'intervertir l'ordre des termes d'une somme.

On pourra ajouter deux équivalences, l'une à l'autre, mais à la condition de ne pas intervertir l'ordre des termes, si donc on a
\[ A \equiv B \]
\origpage{63}
et
\[ C \equiv D, \]
on pourra en conclure
\[ A + C \equiv B + D, \]
mais non
\[ C + A \equiv B + D. \]
Ajoutons que que\ednote{The word ``que'' is printed twice; an apparent misprint.} de
\[ 2A \equiv 0 \]
on n'a pas le droit de conclure
\[ A \equiv 0. \]

Cela posé, il est clair que l'on peut imaginer un groupe $G$ satisfaisant aux conditions suivantes~:

1° A chaque contour fermé, $M_{0}BM_{0}$ correspondra une substitution $S$ du groupe;

2° La condition nécessaire et suffisante pour que $S$ se réduise à la substitution identique, c'est que
\[ M_{0}BM_{0} \equiv 0; \]\ednote{The strokes of the $\equiv$ are only partly printed here.}

3° Si $S$ et $S'$ correspondent aux contours $C$ et $C'$ et si
\[ C'' \equiv C + C', \]
la substitution correspondant à $C''$ sera $SS'$.

Le groupe $G$ s'appellera le \emph{groupe fondamental} de la variété $V$.

Comparons-le au groupe $g$ des substitutions subies par les fonctions $F$.

Le groupe $g$ sera isomorphe à $G$.

L'isomorphisme pourra être holoédrique.

Il pourra être mériédrique si un contour fermé $M_{0}BM_{0}$ non décomposable en lacets ramène les fonctions $F$ à leurs valeurs primitives.

\subsection*{\S~13. --- Équivalences fondamentales.}

Le groupe $G$ sera dérivé d'un certain nombre de substitutions principales
\origpage{64}
$S_{1}, S_{2}, \ldots, S_{p}$. A chacune d'elles correspondra un contour fermé, de sorte que nous aurons \emph{$p$ contours fermés fondamentaux} $C_{1}, C_{2}, \ldots, C_{p}$ et qu'un contour fermé quelconque soit équivalent à une combinaison des contours fondamentaux se succédant dans un certain ordre.

Ces contours fondamentaux ne sont pas en général tout à fait indépendants et il y a entre eux certaines relations que j'appellerai \emph{équivalences fondamentales}.

Supposons, par exemple, qu'on ait entre les contours fondamentaux, l'équivalence
\[ k_{1}C_{1} + k_{2}C_{2} + k'_{1}C_{1} + k_{3}C_{3} \equiv 0. \]

Cela signifie que la substitution
\[ S^{k_{1}}_{1}S^{k_{2}}_{2}S^{k'_{1}}_{1}S^{k_{3}}_{3} \]
du groupe $G$ se réduit à la substitution identique.

Les équivalences fondamentales nous font donc connaître la forme du groupe $G$.

Supposons que la variété $V$ ait été définie par le mode de représentation du \no 10 et que nous n'ayons qu'un seul polyèdre $P_{1}$. Il est clair qu'on obtiendra tous les contours fondamentaux de la manière suivante. Soient $M_{0}$ un point intérieur à $P_{1}$, $A$ un point d'une des faces de $P_{1}$, $A'$ le point correspondant sur la face conjuguée. On ira de $M_{0}$ à $A$, puis de $A'$ à $M_{0}$ sans sortir de $P_{1}$; le chemin correspondant sur la variété $V$ sera fermé.

Il y aura donc autant de contours fondamentaux que de paires de faces.

Voici maintenant comment on formera les équivalences fondamentales.

Considérons un cycle d'arêtes. Soit, par exemple, une arête, intersection des faces $F_{1}$ et $F'_{\mu}$, et que je désignerai pour cette raison sous le nom d'\emph{arête} $F_{1}F'_{\mu}$; soient $F'_{1}$ la face conjuguée de $F_{1}$ et $F_{2}$, $F'_{1}$\ednote{A comma is printed between $F_{2}$ and $F'_{1}$; the edge is named $F_{2}F'_{1}$ two lines below.} l'arête correspondante à $F_{1}F'_{\mu}$ sur cette face; soient $F'_{2}$ la face conjuguée de $F_{1}$\ednote{Printed $F_{1}$; apparently for $F_{2}$, since $F'_{1}$ is already the face conjugate to $F_{1}$.} et $F_{3}F'_{2}$ l'arête correspondante à $F_{2}F'_{1}$ sur cette face; et ainsi de suite jusqu'à ce qu'on retombe sur la face $F'_{\mu}$ et l'arête $F_{1}F'_{\mu}$.
\origpage{65}

Remarquons qu'en faisant cette opération l'on pourra retomber, plusieurs fois, sur la même face.

Soient $A_{i}$ un point de $F_{i}$ et $A'_{i}$ le point correspondant de $F'_{i}$; soit $C_{i}$ le contour fondamental
\[ M_{0}A_{i} + A'_{i}M_{0}. \]

Alors, on aura l'équivalence fondamentale
\[ C_{1} + C_{2} + \ldots + C_{\mu} \equiv 0. \]

Il y aura donc autant d'équivalences fondamentales que de cycles d'arêtes.

Quand on aura formé ainsi les équivalences fondamentales, on en déduira les homologies fondamentales qui n'en diffèrent que parce que l'ordre des termes est indifférent. La connaissance de ces homologies fera immédiatement connaître le nombre de Betti $P_{1}$.

Appliquons ces principes aux exemples cités plus haut, et remarquons que, toutes les variétés citées dans ces exemples étant fermées et à trois dimensions, on aura
\[ P_{1} = P_{2}. \]

\subsection*{Premier exemple.}

\[ \begin{aligned} &(C_{1}) &&ABDC \equiv A'B'D'C',\\ &(C_{2}) &&ABB'A' \equiv CDD'C',\\ &(C_{3}) &&ACC'A' \equiv BDD'B'. \end{aligned} \]

Voici ce que j'entends par la notation
\[ \begin{aligned} &(C_{1}) &&ABDC \equiv A'B'D'C'. \end{aligned} \]

Je veux dire que la face $ABDC$ est conjuguée de $A'B'D'C'$ et que $\alpha$ désignant un point de $ABDC$ et $\alpha'$ un point de $A'B'D'C'$, le contour fondamental $M_{0}\alpha + \alpha'M_{0}$ est désigné par $C_{1}$.
\origpage{66}

\subsection*{Équivalences fondamentales.}

\[ C_{1} + C_{2} \equiv C_{2} + C_{1}, \qquad C_{1} + C_{3} \equiv C_{3} + C_{1}, \qquad C_{2} + C_{3} \equiv C_{3} + C_{2}. \]

Les homologies fondamentales se réduisent à des identités
\[ P_{1} = P_{2} = 4. \]

Je passe tout de suite au troisième exemple, puisque nous avons vu que le second doit être rejeté.

\subsection*{Troisième exemple.}

\[ \begin{aligned} &(C_{1}) &&ABDC \equiv B'D'C'A',\\ &(C_{2}) &&ABB'A' \equiv DD'C'C,\\ &(C_{2}) &&ACC'A' \equiv DD'B'B. \end{aligned} \]
\ednote{The third row is labelled $C_{2}$ in the print, repeating the second label; apparently for $C_{3}$.}

\subsection*{Equivalences fondamentales.}

\[ \begin{aligned} &C_{1} + C_{3} + C_{2} \equiv 0, &&C_{1} - C_{3} - C_{2} \equiv 0,\\ &C_{2} - C_{1} - C_{3} \equiv 0, &&C_{3} - C_{2} - C_{1} \equiv 0, \end{aligned} \]
\ednote{The second minus sign in $C_{3} - C_{2} - C_{1} \equiv 0$ did not print; a blank gap stands in its place.}
ce qui peut aussi s'écrire
\[ 2C_{1} \equiv 2C_{2} \equiv 2C_{3}, \qquad 4C_{1} = 0. \]

\subsection*{Homologies fondamentales.}

\[ \begin{aligned} &C_{1} \backsim C_{2} \backsim C_{3} \backsim 0,\\ &P_{1} = P_{2} = 1. \end{aligned} \]

On peut donner de ce résultat une interprétation géométrique simple.

Le groupe $G$ est d'ordre fini et se compose seulement de huit substitutions
\origpage{67}
distinctes correspondant aux contours suivants~:
\[ 0, \qquad C_{1}, \qquad C_{2}, \qquad C_{3}, \qquad 2C_{1}, \qquad 3C_{1}, \qquad 3C_{2}, \qquad 3C_{1}. \]
\ednote{The last term reads $3C_{1}$ in the print, repeating the sixth; apparently for $3C_{3}$. Its subscript is faintly printed.}

Le groupe est isomorphe au groupe suivant
\[ \begin{aligned} &(x, y, z, t; -y, x, -t, z;\\ &-t, -z, y, x; z, -t, -x, y;\\ &-x, -y, -z, -t; y, -x, t. -z;\\ &t, z, -y, -x; -z, t, x, -y). \end{aligned} \]
\ednote{The print sets a period after $t$ in the third line, $t. -z$; apparently for a comma.}

Ce groupe, qui transforme en lui-même l'hypercube à quatre dimensions, dont les huit faces ont pour équations
\[ x = \pm 1, \qquad y = \pm 1, \qquad z = \pm 1, \qquad t = \pm 1, \]
\ednote{The last $\pm$ is imperfectly printed; its strokes are broken apart.}
pourrait être appelé le groupe \emph{hypercubique}.

\subsection*{Quatrième exemple.}

\[ \begin{aligned} &(C_{1}) &&ABDC \equiv B'D'C'A',\\ &(C_{2}) &&ABB'A' \equiv CDD'C',\\ &(C_{3}) &&ACC'A' \equiv BDD'B'. \end{aligned} \]

\subsection*{Equivalences fondamentales.}

\[ C_{2} + C_{3} \equiv C_{3} + C_{2}, \qquad C_{3} + C_{1} \equiv C_{1} + C_{2}, \qquad -C_{2} + C_{1} \equiv C_{1} + C_{3}. \]

\subsection*{Homologies fondamentales.}

\[ C_{2} \backsim C_{3} \backsim 0, \]
d'où
\[ P_{1} = P_{2} = 2. \]
\origpage{68}

\subsection*{Cinquième exemple.}

\[ \begin{aligned} &(C_{1}) &&ABC \equiv FED,\\ &(C_{2}) &&ACE \equiv FDB,\\ &(C_{3}) &&AED \equiv FBC,\\ &(C_{4}) &&ADB \equiv FCE. \end{aligned} \]

\subsection*{Equivalences fondamentales.}

\[ C_{1} \equiv C_{2} \equiv C_{3} \equiv C_{4}, \qquad 2C_{1} \equiv 0, \]
d'où
\[ C_{1} \backsim C_{2} \backsim C_{3} \backsim C_{4} \backsim 0, \]
et
\[ P_{1} = P_{2} = 1. \]

Le groupe $G$ se réduit à deux substitutions correspondant aux contours
\[ 0 \qquad \text{et} \qquad C_{1}. \]

\subsection*{Sixième exemple.}

Le groupe $G$ est évidemment holomorphe au groupe $(\alpha, \beta, \gamma, \delta)$.

Les trois substitutions
\[ \begin{aligned} &(x, y, z; x + 1, y, z),\\ &(x, y, z; x, y + 1, z),\\ &(x, y, z; \alpha x + \beta y, \gamma x + \delta y, z + 1) \end{aligned} \]
correspondront respectivement aux contours fondamentaux $C_{1}$, $C_{2}$ et $C_{3}$.

Nous avons d'abord les équivalences fondamentales
\[ \begin{aligned} &C_{1} + C_{2} \equiv C_{2} + C_{1},\\ &C_{1} + C_{3} \equiv C_{3} + \alpha C_{1} + \gamma C_{2},\\ &C_{2} + C_{3} \equiv C_{1} + \beta C_{1} + \delta C_{2}, \end{aligned} \]
\ednote{The third row reads $C_{1} + \beta C_{1}$ on the right in the print; apparently for $C_{3} + \beta C_{1}$, as in the row above.}
\origpage{69}
d'où il résulte d'abord que toute combinaison des contours fondamentaux peut se mettre sous la forme
\[ m_{3}C_{3} + m_{1}C_{1} + m_{2}C_{2}, \]
les $m$ étant des entiers. Et comme cette expression ne pourrait être équivalente à 0, que si les trois entiers $m$ étaient nuls, il en résulte que nous possédons là toutes les équivalences fondamentales.

\subsection*{Homologies fondamentales.}

\[ \begin{aligned} &(\alpha - 1)C_{1} + \gamma C_{2} \backsim 0,\\ &\beta C_{1} + (\delta - 1) C_{1} \backsim 0. \end{aligned} \]
\ednote{The second homology reads $(\delta - 1) C_{1}$ in the print; apparently for $(\delta - 1) C_{2}$.}

Si ces deux homologies ne sont pas distinctes, on aura
\[ C_{1} \backsim C_{2} \backsim 0; \]
d'où
\[ P_{1} = P_{2} = 2; \]
c'est ce qui arrive dans le cas général et particulièrement pour notre quatrième exemple.

Si le déterminant de ces homologies est nul, c'est-à-dire que
\[ (\alpha - 1)(\delta - 1) - \beta\gamma = 0 \]
ou
\[ \alpha + \delta = 2, \]
on aura
\[ P_{1} = P_{2} = 3, \]
sauf dans le cas où les deux homologies se réduiraient à des identités. C'est ce qui arrive pour
\[ \alpha = \delta = 1, \qquad \beta = \gamma = 0, \]
c'est-à-dire pour notre premier exemple; on a alors
\[ P_{1} = P_{2} = 4. \]
\origpage{70}

\subsection*{\S~14. --- Conditions de l'homéomorphisme.}

On sait que deux variétés fermées à deux dimensions qui ont même nombre de Betti sont homéomorphes. C'est ce qui résulte, par exemple, de l'étude des périodes des fonctions abéliennes. Considérons une surface de Riemann $R$ et soit $z$ la variable imaginaire correspondante; on pourra introduire une variable imaginaire nouvelle $t$, telle que $z$ soit une fonction fuchsienne de $t$ et que $t$, considéré comme fonction de $z$, n'ait aucun point singulier sur la surface $R$. Tous les groupes fuchsiens correspondant à diverses surfaces de Riemann ayant même ordre de connexion seront isomorphes.

Ce groupe fuchsien ne sera, d'ailleurs, évidemment autre chose que le groupe fondamental $g$, relatif à la surface $R$ considérée comme une variété à deux dimensions.

On remarquera que tous les groupes fuchsiens ne sont pas susceptibles de définir ainsi une variété fermée à deux dimensions. Considérons le polygone fuchsien fondamental $R_{0}$ auquel, si la fonction fuchsienne existe dans tout le plan, il faudra adjoindre son symétrique $R'_{0}$ par rapport à l'axe des quantités réelles; mais alors le domaine $R_{0} + R'_{0}$ ne sera plus toujours simplement connexe. A chaque point de la variété fermée $V$ devrait correspondre un point de $R_{0}$ (ou de $R_{0} + R'_{0}$) et un seul, et réciproquement. Supposons qu'il existe un ou plusieurs cycles de sommets et que la somme des angles de ce cycle soit zéro ou $\frac{2\pi}{n}$, $n$ étant un entier plus grand que 1; soit alors $M$ le point de $V$ qui correspond à ce cycle de sommets et envisageons un lacet infiniment petit enveloppant $M$. D'après la définition du groupe $g$, à ce lacet doit correspondre, dans le groupe $g$, la substitution identique et, dans le groupe fuchsien, une substitution non identique. Le groupe fuchsien ne peut donc être isomorphe à $g$.

Restent alors les groupes fuchsiens de la première famille tels que la somme des angles d'un cycle quelconque soit $2\pi$, et ceux de la troisième famille. Mais ceux-ci doivent également être rejetés. En effet, si le groupe est de la troisième famille, le domaine $R_{0} + R'_{0}$ n'est plus simplement
\origpage{71}
connexe. Soit $C$ un contour fermé tracé dans ce domaine et tel que l'on n'ait pas
\[ C \backsim 0. \]

A ce contour correspondra, dans le groupe fuchsien, une substitution identique (puisque la variable $z$ sera revenue à son point de départ) et, dans le groupe $g$, une substitution non identique. Ici encore, le groupe fuchsien ne peut être isomorphe à $g$.

Il ne reste donc que les groupes de la première famille et tels que la somme des angles de chaque cycle soit égale à $2\pi$.

Tous ceux de ces groupes qui sont du même genre sont isomorphes entre eux, et c'est pour cette raison que toutes les variétés fermées à deux dimensions qui ont même nombre de Betti sont homéomorphes.

En est-il de même quand le nombre des dimensions est plus grand; deux variétés fermées à $h$ dimensions ($h > 2$) qui ont mêmes nombres de Betti sont-elles homéomorphes?

Nous allons voir que non et c'est ce qui nous fait voir combien les questions d'\emph{Analysis Situs} se compliquent quand le nombre des dimensions s'accroît.

Il est clair d'abord que, si deux variétés sont homéomorphes, leurs groupes $g$ seront isomorphes.

Revenons maintenant à notre sixième exemple et recherchons si deux groupes $(\alpha, \beta, \gamma, \delta)$ peuvent être isomorphes.

Soient $(\alpha, \beta, \gamma, \delta)$ et $(\alpha', \beta', \gamma', \delta')$ ces deux groupes; soient $C_{1}$, $C_{2}$, $C_{3}$; $C'_{1}$, $C'_{2}$, $C'_{3}$ les trois contours fondamentaux de chacun de ces groupes;
\[ \left\{ \begin{aligned} &C_{1} + C_{2} \equiv C_{2} + C_{1},\\ &C_{1} + C_{3} \equiv C_{3} + \alpha C_{1} + \gamma C_{2},\\ &C_{2} + C_{3} \equiv C_{3} + \beta C_{1} = \delta C_{2}, \end{aligned} \right. \tag{1} \]
\ednote{The third row prints $= \delta C_{2}$ at its end, with two clear strokes; apparently for $+ \delta C_{2}$, as in the corresponding row of (1 bis).}
\[ \left\{ \begin{aligned} &C'_{1} + C'_{2} \equiv C'_{2} + C'_{1},\\ &C'_{1} + C'_{3} \equiv C'_{3} + \alpha' C'_{1} + \gamma' C'_{2},\\ &C'_{2} + C'_{3} \equiv C'_{3} + \beta' C'_{1} + \delta' C'_{2}, \end{aligned} \right. \tag{1 \textit{bis}} \]
les équivalences fondamentales des deux groupes\ednote{No final period is printed after ``groupes''; the next page opens a new paragraph.}
\origpage{72}

Supposons les deux groupes isomorphes et soient
\[ \begin{aligned} &a_{3}C_{3} + a_{1}C_{1} + a_{2}C_{2},\\ &b_{3}C_{3} + b_{1}C_{1} + b_{2}C_{2},\\ &c_{3}C_{3} + c_{1}C_{1} + c_{2}C_{2} \end{aligned} \]
les contours du premier groupe qui correspondent respectivement aux contours $C'_{1}$, $C'_{2}$ et $C'_{3}$ du second groupe; les $a$, les $b$ et les $c$ sont des entiers; nous avons vu en effet, plus haut, que tout contour du premier groupe peut se mettre sous cette forme.

Pour que l'isomorphisme ait lieu, il faut qu'en substituant, dans les équivalences (1~\emph{bis}), $a_{3}C_{3} + a_{1}C_{1} + a_{2}C_{2} + \ldots$ à la place de $C'_{1}$, $C'_{2}$, $C'_{3}$, on retrouve les équivalences (1).

Il faut donc d'abord que les substitutions (que je désigne par les mêmes signes que les contours correspondants)
\[ \begin{aligned} &a_{3}C_{3} + a_{1}C_{1} + a_{2}C_{2},\\ &b_{3}C_{3} + b_{1}C_{1} + b_{2}C_{2} \end{aligned} \]
soient permutables. Pour simplifier l'écriture, je vais employer la notation suivante~:

Je poserai
\[ a_{3} = h, \qquad b_{3} = k, \qquad a_{1}C_{1} + a_{2}C_{2} = S_{0}; \qquad b_{1}C_{1} + b_{2}C_{2} = T_{0}, \]
de sorte que nos deux premières substitutions se réduisent à $hC_{3} + S_{0}$, $kC_{3} + T_{0}$.

Je désignerai par $S_{1}$ ce que devient $S_{0}$ quand on y fait subir aux coefficients la substitution linéaire $(\alpha, \beta, \gamma, \delta)$, c'est à-dire quand on remplace $a_{1}$ et $a_{2}$ par
\[ \alpha a_{1} + \beta a_{2} \qquad \text{et} \qquad \gamma a_{1} + \delta a_{2}. \]

$S_{2}$ sera ce que devient $S_{1}$ quand on fait subir à ses coefficients la même substitution et ainsi de suite; de même $T_{1}$, $T_{2}$, \ldots, seront les transformés successifs de $T_{0}$.
\origpage{73}

Cela posé, on aura, en vertu des équivalences (1),
\[ S_{0} + hC_{3} = hC_{3} + S_{h}. \]

Pour que nos substitutions soient permutables, il faut donc que
\[ hC_{3} + S_{0} + kC_{3} + T_{0} = kC_{3} + T_{0} + hC_{3} + S_{0} \]
ou
\[ (k + h)C_{3} + S_{k} + T_{0} \equiv (k + h)C_{3} + T_{h} + S_{0} \]\ednote{The subscript of the second $C_{3}$ prints like $1$; apparently a damaged $3$.}
ou
\[ S_{k} + T_{0} \equiv T_{h} + S_{0}. \tag{$\lambda$} \]

Supposons d'abord que $h$ soit égal à $k$ et différent de 0; l'équivalence précédente pourra être remplacée par les égalités
\[ \begin{aligned} &(\alpha_{h} - 1)(a_{1} - b_{1}) + \beta_{h}(a_{2} - b_{2}) = 0,\\ &\gamma_{h}(a_{1} - b_{1}) + (\delta_{h} - 1)(a_{2} - b_{2}) = 0, \end{aligned} \]
où j'ai désigné par
\[ \begin{vmatrix} \alpha_{h} & \beta_{h} \\[1ex] \gamma_{h} & \delta_{h} \end{vmatrix} \]
les coefficients de la puissance $h^{\text{ième}}$ de la substitution linéaire
\[ \begin{vmatrix} \alpha & \beta \\[1ex] \gamma & \delta \end{vmatrix}. \]

Ces égalités peuvent être satisfaites de deux manières~:

1° Ou bien si
\[ a_{1} = b_{1}, \qquad a_{2} = b_{2}, \]
auquel cas les deux substitutions
\[ hC_{3} + S_{0}, \qquad kC_{3} + T_{0} \]
seraient identiques;
\origpage{74}

2° Ou bien si le déterminant
\[ \begin{vmatrix} \alpha_{h} - 1 & \beta_{h} \\[1ex] \gamma_{h} & \delta_{h} - 1 \end{vmatrix} \]
est nul. Mais cela ne peut arriver que si l'on a
\[ s^{h} = 1, \]
$s$ étant une des racines de l'équation
\[ \begin{vmatrix} \alpha - s & \beta \\[1ex] \gamma & \delta - s \end{vmatrix} = 0, \tag{2} \]
c'est-à-dire si la substitution $(\alpha, \beta, \gamma, \delta)$ est \emph{elliptique}; ce qui veut dire que les racines de l'équation (2) en $s$ sont imaginaires (dans l'espèce, elles doivent être égales à une racine $h^{\text{ième}}$ de l'unité) ou \emph{parabolique}, ce qui veut dire que les racines de l'équation (2) sont égales.

Supposons donc que $(\alpha, \beta, \gamma, \delta)$ soit \emph{hyperbolique}, c'est-à-dire que les racines de l'équation en $s$ soient réelles; et ne supposons plus $h = k$.

Alors la puissance $k^{\text{ième}}$ de
\[ hC_{3} + S_{0} \]
et la puissance $h^{\text{ième}}$ de
\[ kC_{3} + T_{0} \]
devront être permutables; soient
\[ \begin{aligned} &h'C_{3} + S'_{0},\\ &h'C_{3} + T'_{0} \end{aligned} \]\ednote{The second line prints $h'C_{3}$; apparently for $k'C_{3}$, since $h' = k'$ is stated just below.}
ces deux puissances; comme on aura
\[ h' = k' = hk, \]
ces deux puissances, d'après ce qui précède, devraient être identiques. Dans le groupe $(\alpha', \beta', \gamma', \delta')$ supposé isomorphe, la puissance $k^{\text{ième}}$ de $C'_{1}$
\origpage{75}
devrait être identique à la puissance $h^{\text{ième}}$ de $C'_{2}$. Comme il ne saurait en être ainsi, on devra avoir
\[ h = k = 0, \]
c'est-à-dire
\[ a_{3} = b_{3} = 0. \]

Passons au cas des substitutions elliptiques auxquelles nous assimilerons la substitution
\[ \begin{vmatrix} -1 & 0 \\[1ex] 0 & -1 \end{vmatrix}. \]

Alors le déterminant
\[ \begin{vmatrix} \alpha_{h} - 1 & \beta_{h} \\[1ex] \gamma_{h} & \delta_{h} - 1 \end{vmatrix} \]
peut s'annuler pour une certaine valeur de $h$ que j'appelle $\nu$.

Mais alors la substitution
\[ \begin{vmatrix} \alpha_{\nu} & \beta_{\nu} \\[1ex] \gamma_{\nu} & \delta_{\nu} \end{vmatrix} \]
se réduit à la substitution identique, c'est-à-dire que l'on a
\[ \alpha_{\nu} = \delta_{\nu} = 1, \qquad \beta_{\nu} = \gamma_{\nu} = 0. \]

Alors il arrive d'abord que $\nu C_{3}$ est permutable à $C_{1}$ et à $C_{2}$ et plus généralement que deux substitutions
\[ \begin{aligned} &a_{3}C_{3} + a_{1}C_{1} + a_{2}C_{2},\\ &b_{3}C_{3} + b_{1}C_{1} + b_{2}C_{2} \end{aligned} \]
sont permutables, pourvu que $a_{3}$ et $b_{3}$ soient divisibles par $\nu$. Cette condition suffisante est d'ailleurs nécessaire, en excluant, comme nous l'avons fait plus haut, le cas où deux puissances de nos deux substitutions seraient identiques.

Pour le démontrer, je vais encore employer une notation plus abrégée; le symbole
\[ m_{1}C_{1} + m_{2}C_{2} \]
\origpage{76}
n'a de sens que si $m_{1}$ et $m_{2}$ sont entiers; mais le suivant
\[ \mu(a'_{1}C_{1} + a'_{2}C_{2}) + \rho(b'_{1}C_{1} + b'_{2}C_{2}) \]
peut avoir un sens alors même que $\mu$, $\rho$, les $a'$ et les $b'$ ne sont pas entiers, pourvu que
\[ \mu a'_{1} + \rho b'_{1}, \qquad \mu a'_{1} + \rho b'_{2} \]\ednote{The second expression prints $\mu a'_{1}$; apparently for $\mu a'_{2}$.}
soient entiers. Il est évident que ce qui nous permet d'agir ainsi, c'est que $C_{1}$ et $C_{2}$ sont permutables.

Choisissons alors
\[ \begin{aligned} &\xi_{0} = a'_{1}C_{1} + a'_{2}C_{2},\\ &\eta_{0} = b'_{1}C_{1} + b'_{2}C_{2}, \end{aligned} \]
de telle sorte que
\[ \xi_{1} = \xi_{0}s, \qquad \eta_{1} = \eta_{0}s^{-1}; \]
$\xi_{1}$ et $\eta_{1}$ sont, d'après nos conventions, les transformées de $\xi_{0}$ et $\eta_{0}$ par la substitution linéaire $(\alpha, \beta, \gamma, \delta)$; $s$ est une des racines de l'équation (2). D'après la théorie des substitutions linéaires, on peut toujours choisir les nombres $a'$ et $b'$ (qui seront généralement irrationnels et même imaginaires) de façon qu'il en soit ainsi; posons alors
\[ S_{0} = \mu\xi_{0} + \rho\eta_{0}, \qquad T_{0} = \mu'\xi_{0} + \rho'\eta_{0}. \]

On aura
\[ \mu\xi_{0} + \rho\eta_{0} + hC_{3} \equiv hC_{3} + \mu s^{h}\xi_{0} + \rho s^{-h}\eta_{0} \]
et
\[ k(hC_{3} + \mu\xi_{0} + \rho\eta_{0}) \equiv khC_{3} + \mu\frac{s^{kh} - 1}{s^{h} - 1}\xi_{0} + \rho\frac{s^{-kh} - 1}{s^{-h} - 1}\eta_{1}. \]\ednote{The last subscript prints like $1$; apparently for $\eta_{0}$.}

Cela posé, l'équivalence ($\lambda$) peut s'écrire
\[ \begin{aligned} &\mu(s^{k} - 1) = \mu'(s^{h} - 1),\\ &\rho(s^{-k} - 1) = \rho'(s^{-h} - 1). \end{aligned} \]

Si $k$ et $h$ ne sont pas divisibles par $\nu$, ces relations ne sont pas satisfaites
\origpage{77}
identiquement et nous pourrons poser
\[ \mu = (s^{h} - 1)\varepsilon, \qquad \mu' = (s^{k} - 1)\varepsilon, \qquad \rho = (s^{-h} - 1)\zeta, \qquad \rho' = (s^{-k} - 1)\zeta. \]

Mais alors
\[ \begin{aligned} &k(hC_{3} + S_{0}) \equiv khC_{3} + \varepsilon(s^{kh} - 1)\xi_{0} + \zeta(s^{-kh} - 1)\eta_{0},\\ &h(kC_{3} + T_{0}) \equiv khC_{3} + \varepsilon(s^{kh} - 1)\xi_{0} + \zeta(s^{-kh} - 1)\eta_{0}, \end{aligned} \]
ce qui montre que la puissance $k^{\text{ième}}$ de notre première substitution est identique à la puissance $h^{\text{ième}}$ de la seconde. Comme nous excluons ce cas, nous concluons que l'on doit avoir
\[ a_{3} \equiv b_{3} \equiv 0 \qquad (\text{mod.}\ \nu). \]

Mais on peut aller plus loin; d'abord, si nous appelons $C'_{1}$ et $C_{2}$\ednote{Printed $C_{2}$ without a prime; apparently for $C'_{2}$.} nos deux premières substitutions, nous pouvons les remplacer par
\[ \omega_{1}C'_{1} + \omega_{2}C'_{2}, \qquad \omega'_{1}C'_{1} + \omega'_{2}C'_{3}, \]\ednote{The last term prints $C'_{3}$ here and again in the parenthesis below; apparently for $C'_{2}$, as in the same expression on p. 79.}
les $\omega$ étant des entiers tels que $\omega_{1}\omega'_{2} - \omega_{2}\omega'_{1} = 1$. Nous pouvons donc toujours supposer que
\[ b_{3} = 0 \]
(sans quoi l'on remplacerait $C'_{1}$ et $C'_{2}$ par $\omega_{1}C'_{1} + \omega_{2}C'_{2}$ et $\omega'_{1}C'_{1} + \omega'_{2}C'_{3}$ en choisissant les $\omega$ de façon à annuler le nouveau $b_{3}$).

Mais le sous-groupe dérivé de $C'_{1}$ et de $C'_{2}$, à cause de l'isomorphisme avec $(\alpha', \beta', \gamma', \delta')$, doit être permutable à toutes les substitutions du groupe et en particulier à $C_{3}$.

Écrivons donc que la substitution
\[ -C_{3} + C'_{2} + C_{3} \]
fait partie du sous-groupe dérivé de $C'_{1}$ et de $C'_{2}$. Nous avons
\[ \begin{aligned} &C'_{2} = \mu'\xi_{0} + \rho'\eta_{0},\\ &-C_{3} + C'_{2} + C_{3} = \mu's\xi_{0} + \rho's^{-1}\eta_{0}. \end{aligned} \]
\origpage{78}

Pour que cette seconde substitution fasse partie de son groupe, il faut, si $a_{3}$ n'est pas nul, que cette dernière substitution soit un multiple de $C'_{2}$.

Or, cela ne peut évidemment avoir lieu que si $s = s^{-1} = -1$.

Si on laisse ce cas de côté, on devra avoir
\[ a_{3} = b_{3} = 0. \]

Si donc nous laissons de côté le cas où
\[ \alpha + \delta = \pm 2 \]
[substitutions paraboliques et substitution $(-1, 0, 0, -1)$], nous devrons avoir
\[ a_{3} = b_{3} = 0. \]

J'ajoute que $c_{3}$ devra être égal à 1, sans quoi les combinaisons des trois substitutions fondamentales
\[ \begin{aligned} &a_{1}C_{1} + a_{2}C_{2} \equiv C'_{1},\\ &b_{1}C_{1} + b_{2}C_{2} \equiv C'_{2},\\ &c_{3}C_{3} + c_{1}C_{1} + c_{2}C_{2} \equiv C'_{3} \end{aligned} \]
ne pourraient pas reproduire toutes les substitutions
\[ m_{3}C_{3} + m_{1}C_{1} + m_{2}C_{2} \]
du groupe $(\alpha, \beta, \gamma, \delta)$, mais seulement celles où l'entier $m_{3}$ serait divisible par $c_{3}$.

Maintenant toute substitution du groupe $(\alpha, \beta, \gamma, \delta)$ doit pouvoir être mise sous la forme
\[ m_{3}C'_{3} + m_{1}C'_{1} + m_{2}C'_{2}. \]

Pour que $C_{1}$ puisse être mis sous la forme, il faut d'abord que $m_{3}$ soit nul; et ensuite que l'on ait identiquement
\[ C_{1} = m_{1}(a_{1}C_{1} + a_{2}C_{2}) + m_{2}(b_{1}C_{1} + b_{2}C_{2}), \]
\origpage{79}
et de même
\[ C_{2} = m'_{1}(a_{1}C_{1} + a_{2}C_{2}) + m'_{2}(b_{1}C_{1} + b_{2}C_{2}). \]

Comme les $m$ et les $m'$ sont entiers, j'en conclurai que le déterminant
\[ a_{1}b_{2} - a_{2}b_{1} = 1. \]

Mais j'ai dit plus haut que l'on pouvait remplacer $C'_{1}$ et $C'_{2}$ par
\[ \omega_{1}C'_{1} + \omega_{2}C'_{2}, \qquad \omega'_{1}C'_{1} + \omega'_{2}C'_{2}. \]

Si $a_{1}b_{2} - a_{2}b_{1}$ est égal à 1, nous pourrons choisir les $\omega$ de telle façon que
\[ \omega_{1}C'_{1} + \omega_{2}C'_{2} = C_{1}, \qquad \omega'_{1}C'_{1} + \omega'_{2}C'_{2} = C_{2}, \]
c'est-à-dire que nous pourrons toujours supposer
\[ a_{1} = b_{2} = 1, \qquad a_{2} = b_{1} = 0. \]

Mais alors l'équivalence (1~\emph{bis})~:
\[ C'_{1} + C'_{3} \equiv C'_{3} + \alpha'C'_{1} + \gamma'C'_{2} \]
devient
\[ C_{1} + C_{3} \equiv C_{3} + \alpha'C_{1} + \gamma'C_{2}, \]
d'où
\[ \alpha = \alpha', \qquad \gamma = \gamma'; \]
on trouverait de même
\[ \beta = \beta', \qquad \delta = \delta'. \]

Nous devons donc conclure que les deux groupes $(\alpha, \beta, \gamma, \delta)$ et $(\alpha', \beta', \gamma', \delta')$ ne peuvent être isomorphes que si l'on peut passer de l'un à l'autre en changeant $C'_{1}$ et $C'_{2}$ en
\[ \omega_{1}C'_{1} + \omega_{2}C'_{2}, \qquad \omega'_{1}C'_{1} + \omega'_{2}C'_{2}. \]

Cela peut s'énoncer d'une autre manière.
\origpage{80}

Soit
\[ S = \begin{vmatrix} \alpha & \beta \\[1ex] \gamma & \delta \end{vmatrix} \]
une substitution linéaire à coefficients entiers et telle que
\[ \alpha\delta - \beta\gamma = 1. \]

Soit
\[ T = \begin{vmatrix} \omega_{1} & \omega_{2} \\[1ex] \omega'_{1} & \omega'_{2} \end{vmatrix} \]
une autre substitution linéaire à coefficients entiers et telle que
\[ \omega_{1}\omega'_{2} - \omega'_{1}\omega_{2} = 1. \]

La substitution
\[ T^{-1}ST, \]
qui s'appelle la \emph{transformée} de $S$ par $T$, est aussi linéaire à coefficients entiers et a pour déterminant 1.

Si deux substitutions linéaires $S$ et $S'$ à coefficients entiers et de déterminant 1 sont transformées l'une de l'autre par une substitution de la forme $T$, je dirai que $S$ et $S'$ appartiennent à la même classe.

Il est clair d'abord que $S$ et $S'$ ne peuvent appartenir à la même classe que si la somme $\alpha + \delta$ a la même valeur pour l'une et pour l'autre, mais cette condition n'est pas suffisante et à une même valeur entière de $\alpha + \delta$ correspondront plusieurs classes de substitutions linéaires, de la même façon qu'à un même déterminant correspondent plusieurs classes de formes quadratiques.

Ainsi la condition nécessaire et suffisante pour que les deux groupes $(\alpha, \beta, \gamma, \delta)$, $(\alpha', \beta', \gamma', \delta')$ soient isomorphes, c'est que les deux substitutions linéaires $(\alpha, \beta, \gamma, \delta)$, $(\alpha', \beta', \gamma', \delta')$ appartiennent à la même classe.
\origpage{81}

Nous avons laissé de côté le cas où
\[ \alpha + \delta = \pm 2. \]

Si $\alpha + \delta = 2$, la substitution $(\alpha, \beta, \gamma, \delta)$ sera de la même classe que $(1, h, 0, 1)$; le groupe $(\alpha, \beta, \gamma, \delta)$ sera isomorphe à $(1, h, 0, 1)$.

Celui-ci contient une substitution remarquable $C_{2}$, qui n'est multiple d'aucune autre et qui est permutable à toutes les substitutions du groupe. On voit d'ailleurs sans peine que, si $h$ n'est pas nul, $C_{2}$ est la seule substitution qui jouisse de cette propriété.

Nous pouvons d'ailleurs laisser de côté le cas où $h$ est nul, car le groupe $(1, 0, 0, 1)$, dont toutes les substitutions sont permutables deux à deux, ne peut être évidemment isomorphe à aucun autre groupe $(\alpha, \beta, \gamma, \delta)$.

Si $\alpha + \delta = -2$, le groupe $(\alpha, \beta, \gamma, \delta)$ sera isomorphe à $(-1, h, 0, -1)$. Celui-ci contient une substitution remarquable $C_{2}$, qui n'est multiple d'aucune autre, qui n'est pas permutable à toutes les substitutions du groupe, mais qui est permutable \emph{au double} de toutes ces substitutions. Si $h$ n'est pas nul, $C_{2}$ est la seule substitution qui jouisse de cette propriété. Si, au contraire, $h$ est nul, il y en a une infinité, ce qui prouve déjà que $(-1, 0, 0, -1)$ ne peut être isomorphe à $(-1, h, 0, -1)$.

De même, la présence de cette substitution remarquable $C_{2}$ dans $(1, h, 0, 1)$ et l'absence de toute substitution jouissant de la même propriété dans $(-1, h', 0, -1)$ montrent que ces deux groupes ne sauraient être isomorphes.

Il me reste deux questions à résoudre~:

1° Les groupes $(1, h, 0, 1)$ et $(1, h', 0, 1)$ ou\ednote{Printed ``ou'' without accent; apparently for ``où''.} $|h| \gtrless 0$, $|h'| \gtrless 0$, $|h| \gtrless |h'|$, peuvent-ils être isomorphes?

2° Même question pour les groupes $(-1, h, 0, -1)$, $(-1, h', 0, -1)$.

Commençons par la première question.

J'observe d'abord que, $C_{2}$ étant la seule substitution du premier groupe qui jouisse de la propriété caractéristique énoncée plus haut, on devra avoir
\[ C'_{2} \equiv C_{2} \]
\origpage{82}
(ou $C'_{2} \equiv -C_{2}$, mais alors on changerait $C'_{1}$ et $C'_{2}$ en $-C'_{1}$ et $-C'_{2}$).

Maintenant, pour que l'on puisse, en combinant
\[ \begin{aligned} &C'_{1} \equiv a_{3}C_{3} + a_{1}C_{1} + a_{2}C_{2},\\ &C'_{2} \equiv C_{2},\\ &C'_{3} \equiv c_{3}C_{3} + c_{1}C_{1} + c_{2}C_{2}, \end{aligned} \]
retrouver toutes les substitutions du premier groupe, il faut que
\[ a_{3}c_{1} - c_{3}a_{1} = 1. \]

On démontre alors aisément que
\[ C'_{1} + C'_{3} \equiv C'_{3} + C'_{1} + hC'_{2}, \]
ce qui prouve que les deux groupes ne peuvent être isomorphes, que si
\[ h = \pm h'. \]

Passons à la seconde question.

On verrait, comme tout à l'heure, que l'on doit avoir
\[ \begin{aligned} &C'_{2} \equiv C_{2},\\ &a_{3}c_{1} - c_{3}a_{1} = 1 \end{aligned} \]

Il vient ensuite
\[ C'_{1} + C'_{3} \equiv (a_{3} + c_{3})C_{3} + (c_{1} + a_{1}\varepsilon)C_{1} + (c_{2} + a_{2}\varepsilon - a_{1}c_{3}h\varepsilon)C_{2}, \]
où
\[ \varepsilon = (-1)^{c_{3}} \]
et
\[ \begin{aligned} &C'_{3} - C'_{1} \equiv (c_{3} - a_{3})C_{3} + (c_{1} - a_{1})\varepsilon'C_{1}\\ &+ [(c_{2} - a_{2})\varepsilon' + (a_{1} + c_{1})a_{3}\varepsilon'h]C_{2}, \end{aligned} \]
où
\[ \varepsilon' = (-1)^{a_{3}}. \]
\origpage{83}

Si l'on veut que
\[ C'_{1} + C'_{3} \equiv C'_{3} - C'_{1} + h'C'_{2}, \] \ednote{The minus sign in this display prints as two short broken strokes; read as a single $-$.}
il faut que
\[ a_{3} = 0, \qquad \varepsilon' = 1, \qquad \varepsilon = -1, \qquad a_{1}c_{3} = -1, \qquad h = h'. \]

Donc nos deux groupes ne pourront être isomorphes que si
\[ h = \pm h'. \]

Le résultat énoncé plus haut est donc général et s'étend aux substitutions paraboliques.

Si les deux substitutions linéaires $(\alpha, \beta, \gamma, \delta)$, $(\alpha', \beta', \gamma', \delta')$ ne sont pas de la même classe, les deux groupes correspondants ne peuvent être isomorphes.

Il résulte de cette longue discussion que les différents groupes $(\alpha, \beta, \gamma, \delta)$ peuvent donner naissance à une infinité de variétés fermées $V$ distinctes, c'est-à-dire non homéomorphes. Or, le nombre $P_{1}$ ne peut prendre que l'une des trois valeurs 2, 3 ou 4.

\emph{Pour que deux variétés fermées soient homéomorphes, il ne suffit donc pas qu'elles aient mêmes nombres de Betti}.

C'est ce que nos autres exemples mettent également en évidence.

Dans le troisième exemple, le groupe $G$ se réduit à huit substitutions et dans le cinquième exemple, à deux substitutions seulement.

Considérons, au contraire, l'hypersphère
\[ x^{2}_{1} + x^{2}_{2} + x^{2}_{3} + x^{2}_{4} = 1; \]
c'est une variété dont le groupe $G$ se réduit à une seule substitution, la substitution identique.

Voilà donc trois variétés, dont les groupes $G$ sont d'ordre fini; mais ces groupes ne sont pas isomorphes; de sorte que ces variétés ne seront
\origpage{84}
pas homéomorphes; et cependant elles auront mêmes nombres de Betti~:
\[ P_{1} = P_{2} = 1. \]

Il paraîtra naturel de restreindre le sens du mot \emph{simplement connexe} et de le réserver aux variétés, dont le groupe $G$ se réduit à la substitution identique.

Alors une variété fermée de plus de deux dimensions pourra avoir son groupe $G$ d'ordre fini sans être simplement connexe.

Cela n'arriverait pas avec les variétés de deux dimensions~: leur groupe $G$ ne pouvait être d'ordre fini sans se réduire à une seule substitution.

Nous avons vu pourquoi un groupe fuchsien, qui admet un cycle de sommets dont la somme des angles est égale à $\frac{2n}{n}$\ednote{Printed with an italic $n$ in the numerator; apparently for $\frac{2\pi}{n}$.} $(n > 1)$, ne peut être le groupe $G$ d'une variété fermée à deux dimensions.

Il en est de même pour la même raison de tout groupe d'ordre fini ou plus généralement de tout groupe contenant une substitution, dont une puissance entière se réduit à la substitution identique.

Il pourrait être intéressant de traiter les questions suivantes~:

1° Étant donné un groupe $G$ défini par un certain nombre d'équivalences fondamentales, peut-il donner naissance à une variété fermée à $n$ dimensions?

2° Comment doit-on s'y prendre pour former cette variété?

3° Deux variétés d'un même nombre de dimensions, qui ont même groupe $G$, sont-elles toujours homéomorphes?

Ces questions exigeraient de difficiles études et de longs développements. Je n'en parlerai pas ici.

Je veux toutefois attirer l'attention sur un point.

Riemann a rattaché l'étude des courbes algébriques à celle des variétés à deux dimensions au point de vue de l'\emph{Analysis Situs}.

De même, l'étude des surfaces algébriques se rattache à celle des variétés à quatre dimensions. Ces variétés ont trois nombres de Betti~:
\[ P_{1} = P_{3} \qquad \text{et} \qquad P_{2}. \]
\origpage{85}

M.~Picard a montré que si la surface algébrique est la plus générale de son degré, le nombre $P_{1}$ se réduit à 1; il ne prend une valeur plus grande que 1 que dans certains cas particuliers. La connexion multiple se présentait ainsi comme un cas particulier de la connexion simple.

Ce résultat semble paradoxal; peut-être nous le paraîtra-t-il un peu moins maintenant; un groupe $G$ peut être beaucoup plus complexe qu'un autre groupe $G'$ et cependant correspondre à une valeur plus petite du nombre $P_{1}$.

\subsection*{\S~15. --- Autres modes de génération.}

On peut donner des variétés d'autres définitions, qui sont, pour ainsi dire, intermédiaires entre les deux premières.

Par exemple, on peut envisager une variété $V$ formée de l'ensemble des points qui satisfont à certaines inégalités et aux équations
\[ \left\{ \begin{aligned} &F_{1}(x_{1}, x_{2}, \ldots, x_{n}; y_{1}, y_{2}, \ldots, y_{q})\\ &F_{2}(x_{1}, x_{2}, \ldots, x_{n}; y_{1}, y_{2}, \ldots, y_{q})\\ &\ldots\ldots\ldots\ldots\ldots\ldots\\ &F_{p}(x_{1}, x_{2}, \ldots, x_{n}; y_{1}, y_{2}, \ldots, y_{q}) \end{aligned} \right\} = 0, \tag{1} \]
où les $x$ sont des coordonnées d'un point et les $y$ des paramètres arbitraires; le nombre des dimensions de $V$ est alors $n - p + q$.

On peut aussi adjoindre aux équations (1) les relations
\[ \varphi_{\alpha}(y_{1}, y_{2}, \ldots, y_{q}) = 0 \qquad (\alpha = 1, 2, \ldots, \lambda) \]
entre les $y$. Les $y$ ne sont plus alors des paramètres entièrement arbitraires, et le nombre des dimensions de $V$ devient égal à $n - p + q - \lambda$.

On peut encore envisager une variété $V$ définie par certaines inégalités et par les relations
\[ x_{i} = \theta_{i}(y_{1}, y_{2}, \ldots, y_{p}) \qquad (i = 1, 2, \ldots, n), \tag{2} \]
\origpage{86}
les $y$ étant des paramètres liés par $\lambda$ équations
\[ \varphi_{\alpha}(y_{1}, y_{2}, \ldots, y_{p}) = 0 \qquad (\alpha = 1, 2, \ldots, \lambda). \]

Le nombre des dimensions de $V$ est alors $p - \lambda$.

Considérons un instant $y_{1}, y_{2}, \ldots, y_{p}$ comme les coordonnées d'un point $P$ dans l'espace à $p$ dimensions. Les égalités
\[ \varphi_{\alpha} = 0 \]
définiront alors dans cet espace à $p$ dimensions une certaine variété $W$.

A chaque point de $W$ correspondra un point de $V$, puisque les équations (2) expriment les $x$ en fonction des $y$.

Le cas le plus simple est celui où, réciproquement, à un point de $V$ correspond un seul point de $W$. Mais un cas, qui est aussi très intéressant, est le suivant.

Supposons que la variété $W$ demeure inaltérée, quand les $y$ subissent les substitutions d'un certain groupe $G$. Soient $P$ un point de $W$, $P_{1}, P_{2}, \ldots, P_{h}$ les transformées de $P$ par les substitutions de $G$.

L'ensemble des points $P, P_{1}, P_{2}, \ldots, P_{h}$ formeront ce que j'appellerai un \emph{système de points}.

Si les fonctions $\theta_{i}$ ne sont pas altérées par les substitutions de $G$, il est clair qu'aux divers points d'un même système correspondra un même point de $V$.

Le cas intéressant est celui où à un point de $V$ correspond un seul système de points de $W$.

Étant donnée une variété $W$ et un groupe $G$ qui ne l'altère pas, on peut toujours construire une variété $V$, de telle façon qu'à tout point de $V$ corresponde un système de points de $W$ et un seul.

Pour que $V$ soit bilatère, il faut et il suffit que la variété $W$ soit bilatère et que toutes les substitutions de $G$ jouissent de la propriété suivante.

Soient $y_{1}, y_{2}, \ldots, y_{p}$ les coordonnées de $P$, $y'_{1}, y'_{2}, \ldots, y'_{p}$ celles de son transformé; le déterminant fonctionnel des $y'$ par rapport aux $y$ doit être positif.
\origpage{87}

\subsection*{Septième exemple.}

Soit
\[ y^{2}_{1} + y^{2}_{2} + y^{2}_{3} = 1 \]
l'équation de la variété $W$, qui sera ainsi une sphère dans l'espace ordinaire.

Cette sphère n'est pas altérée, quand on change $y_{1}, y_{2}, y_{3}$ en $-y_{1}$, $-y_{2}$, $-y_{3}$. Ce sera là notre groupe $G$.

Si alors nous posons, par exemple,
\[ \begin{aligned} &x_{1} = y^{2}_{1}, &&x_{2} = y^{2}_{2}, &&x_{3} = y^{2}_{3},\\ &x_{4} = y_{2}y_{3}, &&x_{5} = y_{1}y_{3}, &&x_{6} = y_{1}, y_{2}, \end{aligned} \]
\ednote{Printed with a comma in $x_{6} = y_{1}, y_{2}$; apparently for $x_{6} = y_{1}y_{2}$.}
les $x$ ne changeront pas, quand les $y$ changeront de signe, et nous aurons défini une variété $V$ à deux dimensions dans l'espace à 6 dimensions.

Cette variété sera fermée; elle sera \emph{unilatère}.

Soit, en effet, $P$ un point de $W$, dont les coordonnées seront $y_{1}$, $y_{2}$, $y_{3}$. Pour définir la position de ce point sur la sphère $W$, il suffira de connaître deux de ces coordonnées, par exemple $y_{1}$ et $y_{2}$; car l'équation de la sphère nous fait connaître $y_{3}$ en fonction de $y_{1}$ et $y_{2}$.

Son transformé $P'$, dont les coordonnées sont $-y_{1}$, $-y_{2}$ et $-y_{3}$, est diamétralement opposé. Mais il ne faut pas se servir, pour définir la position du point $P$, de $y_{1}$ et de $y_{2}$, parce que $y_{3}$ n'est pas une fonction uniforme de ces deux variables. Il vaut mieux poser
\[ \begin{aligned} &y_{1} = \cos \varphi \sin \theta,\\ &y_{2} = \cos \varphi \sin \theta,\\ &y_{3} = \cos \theta. \end{aligned} \]
\ednote{The second line is printed $y_{2} = \cos \varphi \sin \theta$, identical to the first; apparently for $y_{2} = \sin \varphi \sin \theta$.}

Les coordonnées du point $P$ dans ce nouveau système sont $\varphi$ et $\theta$; celles de $P'$ sont $\varphi + \pi$ et $\pi - \theta$. On voit alors que
\[ \frac{\partial(\varphi + \pi, \pi - \theta)}{\partial(\varphi, \theta)} = -1 < 0. \]

La variété $V$ est donc unilatère.
\origpage{88}

\subsection*{Huitième exemple.}

Soient $y_{1}, y_{2}, \ldots, y_{q}; z_{1}, z_{2}, \ldots, z_{q}$, $2q$ paramètres liés par les relations
\[ \left\{ \begin{aligned} &y^{2}_{1} + y^{2}_{2} + \ldots + y^{2}_{q} = 1,\\ &z^{2}_{1} + z^{2}_{2} + \ldots + z^{2}_{q} = 1. \end{aligned} \right. \tag{3} \]

Si nous regardons ces $2q$ paramètres comme les coordonnées d'un point dans l'espace à $2q$ dimensions, ces équations (3) représenteront une variété fermée $W$ à $2q - 2$ dimensions.

Si nous regardons $y_{1}, y_{2}, \ldots, y_{q}$ et $z_{1}, z_{2}, \ldots, z_{q}$ comme les coordonnées de deux points $Q$ et $Q'$ dans l'espace à $q$ dimensions, ces deux points devront se trouver tous deux sur l'\emph{hypersphère} $S$ qui a pour équation
\[ y^{2}_{1} + y^{2}_{2} + \ldots + y^{2}_{q} = 1, \]
et qui est une variété fermée à $q - 1$ dimensions.

A chaque couple de points de $S$ correspondra ainsi un point de $W$ et un seul et inversement, \emph{pourvu que je convienne de regarder les deux couples de points $QQ'$ et $Q'Q$ comme distincts}.

Considérons maintenant les $\frac{q(q + 3)}{2}$ combinaisons
\[ y_{i} + z_{i}, \qquad y_{i}z_{i}, \qquad y_{i}z_{k} + y_{k}z_{i} \qquad (i, k = 1, 2, \ldots, q), \]
égalons-les à $\frac{q(q + 3)}{2} = n$ variables
\[ x_{1}, \qquad x_{2}, \qquad \ldots, \qquad x_{n}. \]

Nous aurons ainsi défini une variété $V$ à $2q - 2$ dimensions dans l'espace à $n$ dimensions.

Quand on change $y_{i}$ en $z_{i}$ et $z_{i}$ en $y_{i}$ (c'est-à-dire quand on permute les deux points $Q$ et $Q'$), ces $\frac{q(q + 3)}{2}$ combinaisons ne changent pas.
\origpage{89}

A chaque couple de points de l'hypersphère correspond ainsi un point et un seul de la variété $V$ et inversement, mais \emph{à la condition de ne pas regarder comme distincts les deux couples $QQ'$ et $Q'Q$}.

Cette variété $V$ est-elle fermée? Je vais montrer qu'elle ne l'est pas pour $q = 2$ et qu'elle l'est pour $q > 2$.

On a alors
\[ \begin{aligned} &x_{1} = y_{1} + z_{1}, &&x_{2} = y_{1}z_{1}, &&x_{3} = y_{2} + z_{2},\\ &x_{4} = y_{2}z_{2}, &&x_{5} = y_{1}z_{2} + y_{2}z_{1}. \end{aligned} \]

On reconnaît alors que l'on doit avoir pour que les $y$ et les $z$ soient réels
\[ x^{2}_{1} > 4x_{2}, \qquad x^{2}_{3} > 4x_{4}. \]

On aurait des inégalités analogues pour $q > 2$; mais dans quels cas ces inégalités définissent-elles une véritable frontière pour notre variété $V$?

Pour mieux nous en rendre compte, je vais traiter d'abord un exemple plus simple.

Soit d'abord, dans l'espace ordinaire, le cercle
\[ x^{2} + y^{2} = 1, \qquad z = 0. \]

Si l'on convient de ne conserver que les points de cette droite pour lesquels $y$ est positif, alors nous aurons les relations suivantes~:
\[ x^{2} + y^{2} = 1, \qquad z = 0, \qquad y > 0, \]
qui définiront une variété à une dimension (qui, dans ce cas, est une demi-circonférence).

Cette variété n'est pas fermée, elle admet deux points frontières~:
\[ x = \pm 1, \qquad y = z = 0. \]

Considérons au contraire la surface suivante~:
\[ x^{2} + y^{2} - z^{2} + (x^{2} + y^{2} + z^{2})^{2} = 0. \tag{4} \]
\origpage{90}

C'est la surface engendrée par la révolution d'une lemniscate autour de son grand axe. Elle se compose de deux nappes distinctes $N_{1}$ et $N_{2}$ ayant un point conique commun qui est l'origine; on ne peut aller d'une nappe à l'autre qu'en passant par l'origine. Adjoignons alors à l'équation (4) l'inégalité
\[ z > 0. \tag{5} \]

Les relations (4) et (5) définissent une variété à deux dimensions qui n'est autre chose que la nappe $N_{1}$. Cette variété doit être regardée comme fermée; elle est homéomorphe à une sphère; on ne doit pas regarder le point conique
\[ x = y = z = 0 \]
comme un véritable point frontière.

En général, si une variété à $p$ dimensions n'est pas fermée, sa frontière sera formée par une ou plusieurs variétés à $p - 1$ dimensions. Si l'ensemble des points qu'on peut soupçonner d'être des points frontières forme une ou plusieurs variétés à moins de $p - 1$ dimensions, c'est qu'ils ne sont pas de véritables points frontières au point de vue qui nous occupe actuellement, et que la variété donnée est fermée.

Or, on obtiendra ici les points qu'on peut soupçonner d'être frontières, quand on supposera que les points $Q$ et $Q'$ se confondent, c'est-à-dire que
\[ y_{1} = z_{1}, \qquad y_{2} = z_{2}, \qquad \ldots, \qquad y_{q} = z_{q}. \]

On obtiendra ainsi une variété à $q - 1$ dimensions. Ainsi la frontière complète de $V$ n'aura que $q - 1$ dimensions, tandis que $V$ en a $2q - 2$. $V$ sera donc fermée, à moins que
\[ 2q - 2 = (q - 1) + 1 \]
ou
\[ q = 2. \]
\origpage{91}

Pour mieux nous en rendre compte, comparons les deux exemples $q = 2$ et $q = 3$.

Soit d'abord $q = 2$ et envisageons notre variété dans le voisinage du point
\[ y_{1} = z_{1} = 0, \qquad y_{2} = z_{2} = 1, \]
c'est-à-dire du point
\[ x_{1} = 0, \qquad x_{2} = 0, \qquad x_{3} = 2, \qquad x_{4} = 1, \qquad x_{5} = 0. \]

Observons d'abord que, pour les valeurs petites de $x_{1}$ et de $x_{2}$, les trois autres variables $x_{3}$, $x_{4}$ et $x_{5}$ sont développables suivant les puissances de $x_{1}$ et de $x_{2}$; il nous suffira donc d'étudier les variations de $x_{1}$ et de $x_{2}$.

Nous voyons alors que $x_{1}$ et $x_{2}$ peuvent prendre toutes les valeurs qui satisfont à
\[ x^{2}_{1} > 4x_{2}, \]

Le plan des $x_{1}, x_{2}$ se trouve ainsi partagé en deux régions par la ligne
\[ x^{2}_{1} = 4x_{2}, \]
qui est une véritable ligne frontière.

On obtiendrait le même résultat si l'on étudiait la variété $V$ dans le voisinage d'un autre point frontière quelconque. Cette variété n'est donc pas fermée.

Soit maintenant $q = 3$ et
\[ \begin{aligned} &x_{1} = y_{1} + z_{1}, &&x_{2} = y_{1}z_{1}, &&x_{3} = y_{2} + z_{2}, &&x_{4} = y_{2}z_{2}, &&x_{5} = y_{1}z_{2} + y_{2}z_{1},\\ &x_{6} = y_{3} + z_{3}, &&x_{7} = y_{3}z_{3}, &&x_{8} = y_{1}z_{3} + y_{3}z_{1}, &&x_{9} = y_{2}z_{3} + y_{3}z^{2}. \end{aligned} \]\ednote{The last term is printed $y_{3}z^{2}$, with the 2 set as an exponent; apparently a misprint for $y_{3}z_{2}$, matching the pattern of $x_{5}$ and $x_{8}$.}

Étudions la variété $V$ dans le voisinage du point $P_{0}$ qui est tel que
\[ y_{1} = z_{1} = y_{2} = z_{2} = 0, \qquad y_{3} = z_{3} = 1, \]
\origpage{92}
d'où
\[ x_{1} = x_{2} = x_{3} = x_{4} = x_{5} = x_{8} = x_{9} = 0, \qquad x_{6} = 2, \qquad x_{7} = 1. \]\ednote{Between $x_{2}$ and $x_{3}$ only one short stroke of the $=$ printed.}

Nous voyons que, dans le voisinage de ce point, $x_{6}$, $x_{7}$, $x_{8}$, $x_{9}$ sont développables suivant les puissances de $x_{1}$, $x_{2}$, $x_{3}$, $x_{4}$ et $x_{5}$, de sorte qu'il suffit d'étudier les variations de ces cinq dernières variables.

Pour ne conserver que trois variables et rendre possible une représentation géométrique, je coupe ma variété $V$ par la variété plane
\[ x_{1} = 0, \qquad x_{3} = 0, \]
de sorte que l'intersection sera une variété à deux dimensions $V'$.

Soient $x_{2}$, $x_{4}$ et $x_{5}$ les coordonnées d'un point de $V'$; nous pourrons regarder ces trois variables comme les coordonnées d'un point dans l'espace ordinaire et nous aurons ainsi une représentation géométrique de la variété $V'$ ou plutôt de la portion de cette variété qui est voisine de $P_{0}$.

On trouve alors
\[ y_{1} = -z_{1}, \qquad y_{2} = -z_{2}, \]
puisque $x_{1}$ et $x_{3}$ sont supposés nuls, et par conséquent
\[ x_{2} = -y^{2}_{1}, \qquad x_{4} = -y^{2}_{2}, \qquad x_{5} = -2y_{1}y_{2}, \]
d'où
\[ 4x_{2}x_{4} - x^{2}_{5} = 0. \]

C'est l'équation d'un cône du second degré, mais une seule nappe de ce cône convient parce qu'on doit avoir
\[ x_{2} < 0, \qquad x_{4} < 0. \]

La portion du cône qui convient est ainsi séparée de celle qui ne convient pas par le sommet qui ne saurait être regardé comme un véritable point frontière. C'est ainsi que la variété $V'$, de même que $V$, est encore fermée.
\origpage{93}

On obtiendrait un résultat analogue en étudiant $V$ dans le voisinage d'un autre point frontière, ou en coupant $V$ par d'autres variétés planes que la variété
\[ x_{1} = x_{2} = 0. \]\ednote{Printed $x_{1} = x_{2} = 0$; apparently for $x_{1} = x_{3} = 0$, the plane variety used on p. 92.}

Je n'ai voulu donner qu'un exemple destiné à éclaircir le raisonnement qui précède.

En résumé, \emph{la variété $V$ est fermée si $q > 2$ et ne l'est pas si $q = 2$}.

La variété $V$ est elle\ednote{Printed without the hyphen, for ``est-elle''.} unilatère ou bilatère?

Je me propose d'établir qu'elle est bilatère si $q$ est impair et qu'elle est au contraire unilatère si $q$ est pair.

Posons
\[ \begin{aligned} &y_{1} = \cos \theta_{1}, &&z_{2} = \cos \theta'_{1},\\ &y_{2} = \sin \theta_{1} \cos \theta_{2}, &&z_{2} = \sin \theta'_{1} \cos \theta'_{2},\\ &y_{3} = \sin \theta_{1} \sin \theta_{2} \cos \theta_{3}, &&z_{3} = \sin \theta'_{1} \sin \theta'_{2} \cos \theta'_{3},\\ &\ldots\ldots, &&\ldots\ldots,\\ &y_{q-1} = \sin \theta_{1} \sin \theta_{2} \ldots \sin \theta_{q-2} \cos \theta_{q-1}, &&z_{q-1} = \sin \theta'_{1} \sin \theta'_{2} \ldots \sin \theta'_{q-2} \cos \theta'_{q-1},\\ &y_{q} = \sin \theta_{1} \sin \theta_{2} \ldots \sin \theta_{q-2} \sin \theta_{q-1}, &&z_{p} = \sin \theta'_{1} \sin \theta'_{2} \ldots \sin \theta'_{q-1} \sin \theta'_{q-1}. \end{aligned} \]\ednote{The first row of the right-hand column is printed $z_{2}$, apparently for $z_{1}$. The last row appears to be printed $z_{p}$, apparently for $z_{q}$, and its next-to-last factor is printed $\sin \theta'_{q-1}$, apparently for $\sin \theta'_{q-2}$ as in the left-hand column.}

La position d'un point sur $W$ est définie par les $2q - 2$ coordonnées
\[ \theta_{1}, \qquad \theta_{2}, \qquad \ldots, \qquad \theta_{p-1}, \qquad \theta'_{1}, \qquad \theta'_{2}, \qquad \ldots, \qquad \theta'_{q-1}. \]\ednote{Printed $\theta_{p-1}$, apparently for $\theta_{q-1}$.}

D'autre part, le groupe $G$ se compose (outre la substitution identique) de l'unique substitution qui permute $\theta_{i}$ avec $\theta'_{i}$. Pour que la variété $V$ soit bilatère, il faut donc et il suffit que le déterminant fonctionnel
\[ \frac{\partial(\theta_{i}, \theta'_{i})}{\partial(\theta'_{i}, \theta_{i})}, \qquad (i = 1, 2, \ldots, q - 1) \]
soit positif. Or il est égal à
\[ (-1)^{q-1}, \]
\origpage{94}
c'est-à-dire à $+1$ si $q$ est impair et à $-1$ si $q$ est pair. Donc

$V$ est bilatère si $q$ est impair.

$V$ est unilatère si $q$ est pair. C.~Q.~F.~D.

Occupons-nous maintenant de déterminer le nombre de Betti
\[ P_{q-1}. \]

Déterminons d'abord les nombres de Betti pour $W$.

Nous pouvons construire sur $W$ deux variétés à $q - 1$ dimensions de la manière suivante. On sait que tout point de $W$ correspond à un couple de points $QQ'$ de l'hypersphère $S$. Aux couples $Q_{0}Q'$ où $Q_{0}$ est fixe et où $Q'$ décrit l'hypersphère tout entière correspondra alors une variété fermée $U_{1}$ à $q - 1$ dimensions et faisant partie de $W$. De même aux couples $QQ_{0}$, où $Q$ décrit l'hypersphère entière pendant que $Q_{0}$ est fixe, correspondra une variété fermée $U_{2}$ à $q - 1$ dimensions faisant partie de $W$.

Ces deux variétés seront linéairement indépendantes (au point de vue des homologies).

Considérons, en effet, l'intégrale d'ordre $q - 1$
\[ J = \int \sin^{q-2} \theta_{1} \sin^{q-3} \theta_{2} \ldots \sin^{2} \theta_{q-3} \sin \theta_{q-1}\,d\vartheta_{1}\,d\vartheta_{2} \ldots d\theta_{q-1}, \]\ednote{The factor $\sin \theta_{q-1}$ is printed so, apparently for $\sin \theta_{q-2}$. The first two differentials are printed with the curly $\vartheta$, the last with $\theta$, as reproduced.}
et soit $J'$ l'intégrale d'ordre $q - 1$ obtenue en remplaçant dans $J$ les $\theta_{i}$ par les $\theta'_{i}$.

Envisageons ensuite l'intégrale
\[ J + \lambda J', \]
où $\lambda$ est un nombre incommensurable.

Nous connaissons les $x$ en fonction des $\theta$ et des $\theta'$; réciproquement, nous pouvons calculer les $\theta$ et les $\theta'$ en fonctions des $x$ et cela de telle sorte que l'intégrale $J + \lambda J'$ prenne la forme
\[ \int \sum X \partial, \]
\origpage{95}
$X$ étant une fonction entière par rapport aux $x$, et $\partial$ le produit de $q - 1$ différentielles $dx_{i}$.

J'observe en outre que l'intégrale $J + \lambda J'$ est nulle quand elle est étendue à une variété fermée infiniment petite faisant partie de $W$.

Étendue à $U_{2}$, elle est égale à $\sigma$ (surface de l'hypersphère $S$)~: étendue à $U_{1}$, elle est égale à $\lambda\sigma$; et comme il n'y a aucune relation linéaire entre $\sigma$ et $\lambda\sigma$; c'est que $U_{1}$ et $U_{2}$ sont linéairement indépendantes.

Le nombre de Betti $P_{q-1}$ est donc au moins égal à 3.

Pour aller plus loin, considérons la variété $W'$ obtenue en supprimant dans $W$ tous les points des deux variétés $U_{1}$ et $U_{2}$; je me propose d'établir que $W'$ est simplement connexe.

Considérons de même la variété $S'$ en supprimant, dans l'hypersphère $S$, le point $Q_{0}$. A tout couple de points $QQ'$ de la variété $S'$ correspondra un point de $W'$ et réciproquement.

Mais $S'$ est homéomorphe au domaine $D$,
\[ \eta^{2}_{1} + \eta^{2}_{2} + \ldots + \eta^{2}_{q-1} < 1, \]
faisant partie de l'espace à $q - 1$ dimensions où les coordonnées sont désignées par $\eta_{1}, \eta_{2}, \ldots, \eta_{q-1}$.

C'est ainsi que la surface d'une sphère ordinaire, quand on en supprime un point, devient homéomorphe à la surface d'un cercle.

A chaque couple de points du domaine $D$ correspondra un point de $W'$, d'où il suit que $W'$ est homéomorphe au domaine $\Delta$ défini par les inégalités
\[ \begin{aligned} &\eta^{2}_{1} + \eta^{2}_{2} + \ldots + \eta^{2}_{q-1} < 1,\\ &\zeta^{2}_{1} + \zeta^{2}_{2} + \ldots + \zeta^{2}_{q-1} < 1, \end{aligned} \]
et situé dans l'espace à $2q - 2$ dimensions, où les coordonnées sont désignées par $\eta_{i}$ et $\zeta_{i}$.

Or le domaine $\Delta$ est simplement connexe parce qu'il est convexe. Soit, en effet, $v$ une variété fermée d'un nombre quelconque de dimensions et faisant partie de $\Delta$; je dis qu'on peut, en la déformant d'une manière continue,
\origpage{96}
la réduire à un point sans qu'elle sorte de $\Delta$. Construisons une variété $v'$ en multipliant par un facteur positif $k$ plus petit que 1 les coordonnées de tous les points de $v$ (la variété $v'$ sera homothétique à $v$, le centre d'homothétie étant l'origine, et le rapport d'homothétie étant égal à $k$). La variété $v'$ sera tout entière à l'intérieur de $\Delta$ et, quand $k$ décroîtra de 1 à 0; elle coïncidera d'abord avec $v$ et finira par se réduire à un point. C.~Q.~F.~D.

Donc $\Delta$ et par conséquent $W'$ sont simplement connexes. Déterminons maintenant les nombres de Betti pour $W$ et d'abord le nombre $P_{h}$ où $h < q - 1$. Considérons donc une variété fermée à $h$ dimensions contenue dans $W$; cette variété, qu'elle rencontre ou non $U_{1}$ et $U_{2}$, peut toujours être regardée comme homologue à une variété fermée à $h$ dimensions, ne rencontrant ni $U_{1}$ ni $U_{2}$; celle-ci sera homologue à 0 dans $W'$ (puisque $W'$ est simplement connexe) et \emph{a fortiori} dans $W$.

Ainsi $P_{h}$ est égal à 1.

Comme, d'autre part, les nombres de Betti également distincts\ednote{Printed ``distincts'', apparently for ``distants'', the wording of the theorem elsewhere in the memoir.} des extrêmes sont égaux, $P_{h}$ sera encore égal à 1 si $h$ est plus grand que $q - 1$.

Il reste à déterminer $P_{q-1}$.

Considérons une variété fermée $v$ à $q - 1$ dimensions, faisant partie de $W$.

Si elle ne rencontre ni $U_{1}$, ni $U_{2}$, elle sera homologue à 0; en effet, elle fait partie de $W'$, et par conséquent sera homologue à 0 par rapport à $W'$ et \emph{a fortiori} par rapport à $W$.

Supposons qu'elle rencontre $U_{1}$ et $U_{2}$.

Les points de $W$, qui correspondront au couple $Q_{1}Q'$, où $Q_{1}$ est fixe et où $Q'$ décrit l'hypersphère entière, formeront une variété $U'_{1}$, qui sera homologue à $U_{1}$. Par tout point de $W$ passe une variété $U'_{1}$ et une seule.

Les points de $W$ qui correspondent au couple $QQ_{2}$, où $Q_{2}$ est fixe et où $Q$ décrit l'hypersphère entière formeront une variété $U'_{2}$, qui sera homologue à $U_{2}$. Par tout point de $W$ passe une variété $U'_{2}$ et une seule.

Deux variétés $U'_{1}$ et $U'_{2}$ ont un point commun et un seul (à savoir le point qui correspond au couple $Q_{1}Q_{2}$). Au contraire, deux variétés $U'_{1}$ (ou deux variétés $U'_{2}$) ne se rencontrent pas.
\origpage{97}

Cela posé, revenons à la variété $v$ et supposons, pour fixer les idées, qu'elle coupe $U_{1}$ en deux points $M'$ et $M''$ et $U_{2}$ en un point $N'$. Par le point $M'$ je puis faire passer une variété $U'_{2}$, par le point $M''$ une variété analogue $U''_{2}$; par le point $N'$ une variété $U'_{1}$.

Je puis ensuite, en déformant infiniment peu la variété $v$, m'arranger pour qu'elle ne coupe toujours $U_{1}$ et $U_{2}$ qu'aux points $M'$, $M''$ et $N'$, et pour qu'elle admette autour de $M'$ une petite partie commune $u'_{2}$ avec $U'_{2}$; autour de $M''$ une petite partie commune $u''_{2}$ avec $U''_{2}$; autour de $N'$ une petite partie commune $u'_{1}$ avec $U'_{1}$.

Alors l'ensemble des variétés
\[ U'_{1} - u'_{1}, \qquad U'_{2} - u'_{2}, \qquad U''_{2} - u''_{2}, \qquad v - u'_{1} - u'_{2} - u''_{2} \]
formera une variété fermée, qui ne rencontrera ni $U_{1}$, ni $U_{2}$, et qui sera, par conséquent, homologue à 0; on aura donc
\[ U'_{1} - u'_{1} + U'_{2} - u'_{2} + U''_{2} - u''_{2} \backsim v - u'_{1} - u'_{2} - u''_{2}, \]
d'où
\[ v \backsim U'_{2} + U'_{1} + U''_{2} \backsim U_{1} + 2U_{2}. \]

Il pourrait aussi se faire que, par exemple, le nombre que nous avons appelé plus haut $S(M)$ n'ait pas la même valeur pour le point $M'$ considéré comme point d'intersection de $v$ et de $U_{1}$, ou pour le point $M'$ considéré comme point d'intersection de $U'_{2}$ et de $U_{1}$. Dans ce cas, il faudrait remplacer $U'_{2}$ par la variété opposée et l'on aura
\[ v \backsim U'_{1} - U'_{2} + U''_{2} \backsim U_{1}. \]

Dans tous les cas, $v$, $U_{1}$ et $U_{2}$ ne seront pas linéairement indépendants et l'on aura
\[ P_{q-1} = 3. \]

Déterminons enfin les nombres de Betti pour $V$.

Aux couples $QQ'$ et $Q'Q$ correspond un seul et même point de $V$; il
\origpage{98}
en résulte qu'à $U_{1}$ et à $U_{2}$ correspond dans $V$ une seule et même variété, de sorte que je puis écrire
\[ U_{1} \backsim U_{2}. \]

D'autre part, l'intégrale $J + J'$, étendue à cette variété, n'est pas nulle; ce qui montre que l'on n'a pas
\[ U_{1} \backsim 0. \]

Le nombre $P_{q-1}$ est donc au moins égal à 2.

A toute variété fermée $v$, faisant partie de $V$, correspondra une variété $w$ faisant partie de $W$; mais deux cas peuvent se présenter; nous savons qu'à chaque point de $V$ correspondent deux points de $W$; et je pourrai dire, pour abréger le langage, que ces deux points sont \emph{symétriques}, puisqu'on passe de l'un à l'autre en permutant les $y_{i}$ avec les $z_{i}$.

Nous construirons $w$ en prenant pour chaque point de $v$ \emph{un} des deux points qui lui correspondent dans $w$\ednote{The print has lower-case $w$ here; apparently for $W$, the variety in which each point of $V$ has two corresponding points.}; il pourra se faire alors ou bien que $w$ soit fermé ou bien que $w$ ne soit pas fermé, mais que sa frontière se compose de parties symétriques deux à deux.

Considérons d'abord le cas où $w$ est fermé.

Si le nombre des dimensions est différent de $q - 1$, $w$ (et, par conséquent, $v$) pourra, par déformation continue, se réduire à un point et l'on aura
\[ v \backsim 0. \]

Si le nombre des dimensions est égal à $q - 1$, on aura, par rapport à $W$,
\[ w \backsim mU_{1} + nU_{2}, \]
$m$ et $n$ étant entiers; mais par rapport à $V$, $U_{1}$ est homologue à $U_{2}$.

On aura donc, par rapport à $V$,
\[ v \backsim (m + n)U_{1}. \]

Considérons maintenant le cas où $w$ n'est pas fermé. Nous admettrons
\origpage{99}
que le nombre des dimensions est inférieur ou égal à $q - 1$; le cas où ce nombre serait supérieur à $q - 1$ s'y ramène aisément, puisque les nombres de Betti sont égaux deux à deux. La frontière $f$ de $w$ aura alors moins de $q - 1$ dimensions.

Soit $H$ la variété à $q - 1$ dimensions, faisant partie de $W$, qui se compose des points qui sont leur propre symétrique, c'est-à-dire des points
\[ y_{i} = z_{i}. \]

La frontière $f$ peut se déformer d'une manière continue, sans que ses points cessent d'être symétriques deux à deux, et de telle sorte qu'à la fin de la déformation la frontière déformée fasse partie de $H$.

Donc $w$ peut se déformer d'une manière continue, sans cesser de correspondre à une variété $v$ fermée, et de telle sorte qu'à la fin de la déformation la variété $w$ déformée soit fermée.

Donc $v$ est toujours homologue à une variété $v$ correspondant à une variété $w$ fermée.

Donc on a, suivant les cas,
\[ v \backsim 0, \qquad v \backsim (m + n)U_{1}. \]

Donc tous les nombres de Betti de $V$ sont égaux à 1, sauf $P_{q-1}$ qui est égal à 2.

De là deux conséquences~:

1° Soit $q$ impair, $V$ sera bilatère; d'où il suit que, pour une variété bilatère à $4k$ dimensions, le nombre $P_{2k}$ n'est pas forcément impair.

2° Si $q$ est pair, $V$ sera unilatère; d'où il suit que, pour une variété unilatère à $4k + 2$ dimensions, le nombre $P_{2k+1}$ n'est pas forcément impair, ainsi que cela aurait lieu pour une variété bilatère.

J'avais annoncé ces résultats à la fin du \no 9.
\origpage{100}

\subsection*{\S~16. --- Théorème d'Euler.}

On connaît le théorème d'Euler, d'après lequel, si $S$, $A$ et $F$ sont le nombre des sommets, des arêtes et des faces d'un polyèdre convexe, on doit avoir
\[ S - A + F = 2. \]

Ce théorème a été généralisé par M.~l'amiral de Jonquières, pour le cas des polyèdres non convexes. Si un polyèdre forme une variété fermée à deux dimensions, dont le nombre de Betti est $P_{1}$, on aura
\[ S - A + F = 3 - P_{1}. \]

Le fait que les faces sont planes n'a évidemment aucune importance, le théorème s'applique également aux polyèdres curvilignes; il s'applique encore à la subdivision d'une surface fermée quelconque en régions simplement connexes; ces régions correspondent alors aux faces du polyèdre; les lignes, qui servent de frontière à deux de ces régions, correspondent aux arêtes, et les extrémités de ces lignes aux sommets.

Je me propose de généraliser ces résultats pour un espace quelconque.

Soit donc $V$ une variété fermée à $p$ dimensions. Subdivisons-la en un certain nombre de variétés $v_{p}$ à $p$ dimensions; ces variétés $v_{p}$ ne seront pas fermées et leurs frontières seront formées par un certain nombre de variétés $v_{p-1}$ à $p - 1$ dimensions; les frontières des $v_{p-1}$ seront formées à leur tour par un certain nombre de variétés $v_{p-2}$ à $p - 2$ dimensions, et ainsi de suite; j'arriverai ensuite à un certain nombre de variétés $v_{1}$ à une dimension, qui auront pour frontières un certain nombre de points isolés ou de variétés à 0 dimension que j'appellerai $v_{0}$.

La variété $V$ peut avoir des nombres de Betti quelconques, mais je suppose expressément que les variétés $v_{p}, v_{p-1}, \ldots, v_{1}$ sont simplement connexes.

J'appellerai $\alpha_{p}, \alpha_{p-1}, \ldots, \alpha_{1}$ et $\alpha_{0}$ le nombre des $v_{p}$, des $v_{p-1}$, \ldots, des $v_{1}$ et des $v_{0}$.
\origpage{101}

La figure formée par toutes ces variétés pourra s'appeler un polyèdre; car l'analogie avec les polyèdres ordinaires est évidente. Un polyèdre ordinaire est, en effet, une variété fermée $V$ à 2 dimensions, qui est subdivisée en un certain nombre de variétés $v_{2}$, qui sont les faces. Les faces ont pour frontières un certain nombre de variétés $v_{1}$, qui sont les arêtes et qui admettent à leur tour pour frontières un certain nombre de variétés $v_{0}$ appelées \emph{sommets}.

Je me propose de calculer le nombre
\[ N = \alpha_{p} - \alpha_{p-1} + \alpha_{p-2} - \ldots \mp \alpha_{1} \pm \alpha_{0}. \]

J'introduirai ici quelques dénominations nouvelles, peut-être mal justifiées, mais commodes.

Si deux polyèdres sont obtenus en subdivisant de deux manières différentes une variété $V$, je dirai qu'ils sont \emph{congruents}.

Soit maintenant le polyèdre $P$ formé par la variété $V$, par les régions $v_{p}$ et par leurs frontières successives $v_{p-1}, \ldots, v_{1}, v_{0}$.

Subdivisons les $v_{p}$ en régions plus petites $v'_{p}$; les frontières complètes des $v'_{p}$ se composeront d'un certain nombre de régions nouvelles $v''_{p-1}$ et, en outre, de régions $v'_{p-1}$ obtenues en subdivisant les $v_{p-1}$; les frontières complètes des $v'_{p-1}$ et des $v''_{p-1}$ se composeront d'un certain nombre de régions nouvelles $v''_{p-2}$ et, en outre, de régions $v'_{p-2}$ obtenues en subdivisant les $v_{p-2}$, et ainsi de suite; nous arriverons enfin aux $v'_{1}$ et aux $v''_{1}$, dont les frontières complètes se composeront d'un certain nombre de points nouveaux $v''_{0}$ et, en outre, des points $v_{0}$.

Soit $P'$ le polyèdre formé par l'ensemble des régions $v'_{p}$, $v'_{p-1}$, $v''_{p-1}$, $v'_{p-2}$, $v''_{p-2}$, \ldots, $v'_{1}$, $v''_{1}$, $v_{0}$, $v''_{0}$.

Je dirai alors que le polyèdre $P'$ est \emph{dérivé} du polyèdre $P$.

J'éclaircirai cette définition par un exemple emprunté à la Géométrie ordinaire. Considérons un tétraèdre régulier $T$. Dans chacune des faces, je joins chaque sommet au milieu du côté opposé. Chaque face se trouvera ainsi décomposée en six triangles; en tout vingt-quatre triangles. Le polyèdre à vingt-quatre faces ainsi obtenu sera \emph{dérivé} de $T$.
\origpage{102}

Soient maintenant $P$ et $P'$ deux polyèdres congruents, c'est-à-dire obtenus par deux modes de décomposition différents d'une même variété $V$; il existera toujours un polyèdre $P''$ qui sera dérivé à la fois de $P$ et de $P'$ et qu'on obtiendra en combinant les deux modes de décomposition; de telle façon que, si nous appelons $v_{p}$, $v'_{p}$ et $v''_{p}$ les subdivisions de $V$ dans les trois modes de décomposition qui correspondent aux trois polyèdres $P$, $P'$ et $P''$, la condition nécessaire et suffisante pour que deux points appartiennent à la même région $v''_{p}$, c'est qu'ils appartiennent à la fois à la même région $v_{p}$ et à la même région $v'_{p}$.

Je me propose d'établir que le nombre $N$ est le même pour deux polyèdres congruents et, comme nous venons de voir que deux polyèdres congruents ont toujours un dérivé commun, il nous suffira de montrer que le nombre $N$ est le même pour un polyèdre et pour tous ses dérivés.

Si, dans le polyèdre $P$, nous envisageons une des régions $v_{p-1}$, elle appartiendra toujours à deux régions $v_{p}$ et à deux seulement qu'elle séparera l'une de l'autre. Au contraire, une région $v_{p-2}$ appartiendra, en général, à plus de deux régions $v_{p}$ et à plus de deux régions $v_{p-1}$. C'est ainsi que dans les polyèdres ordinaires, une arête sépare toujours deux faces l'une de l'autre, tandis qu'un sommet appartient en général à plus de deux faces et à plus de deux arêtes.

Nous n'excluons pas toutefois le cas où une région $v_{p-2}$ appartiendrait à deux régions $v_{p-1}$ seulement. Ainsi pour un polyèdre ordinaire, nous n'exclurions pas le cas où le milieu d'une arête serait regardé comme un sommet et où cette arête serait par conséquent regardée comme formée de deux arêtes juxtaposées.

Les régions $v_{p-2}$, qui n'appartiendraient ainsi qu'à deux régions $v_{p-1}$, seront dites \emph{singulières}. Soit donc $v_{p-2}$ une région singulière appartenant à deux régions $v_{p-1}$ que j'appelle $v'_{p-2}$\ednote{Printed $v'_{p-2}$; apparently for $v'_{p-1}$, since the two regions named here are $v_{p-1}$.} et $v''_{p-1}$. Il est clair que $v'_{p-1}$ séparera l'une de l'autre \emph{les deux mêmes régions} $v_{p}$ qui sont séparées par $v''_{p-1}$; de sorte que $v_{p-2}$ n'appartiendra aussi qu'à deux régions $v_{p}$ seulement.

De même je dirai que la variété $v_{q}$ est \emph{singulière} si elle n'appartient qu'à deux variétés $v_{q+1}$; et dans ce cas elle n'appartiendra non plus qu'à deux variétés $v_{q+2}$, $v_{q+3}$, \ldots, $v_{p}$.
\origpage{103}

Soit maintenant une variété $v_{h}$; à cette variété appartiendront un certain nombre de variétés $v_{h-1}$; si l'une d'elles est singulière, je dirai que la variété $v_{h}$ est \emph{irrégulière}; elle sera régulière dans le cas contraire.

Considérons donc le polyèdre $P$ avec les régions $v_{p}$, $v_{p-1}$, \ldots\ et le polyèdre dérivé $P'$ avec les régions $v'_{p}$, $v'_{p-1}$, \ldots. Cherchons à remonter du polyèdre $P'$ au polyèdre $P$. Soient deux régions $v'_{p}$ que j'appellerai $\alpha$ et $\beta$; je suppose qu'elles soient séparées l'une de l'autre par une région $v'_{p-1}$ que j'appellerai $\gamma$; qu'elles soient par conséquent contiguës et qu'elles fassent partie d'une même région $v_{p}$. (Comme $\alpha$ et $\beta$ font partie d'une même région $v_{p}$, la région $\gamma$ n'est pas une subdivision de l'une des régions $v_{p-1}$ qui séparent les régions $v_{p}$ les unes des autres; cette région $\gamma$ est donc une de celles que j'avais appelées $v''_{p-1}$ dans la définition des polyèdres dérivés; mais ici je ne fais plus cette distinction et je désigne les variétés que j'avais appelées alors $v''_{p-1}$ aussi bien que celles que j'avais appelées alors $v'_{p-1}$ par la même notation $v'_{p-1}$.)

Cela posé, supprimons la région $\gamma$, qui sert de frontière à $\alpha$ et à $\beta$, et \emph{annexons} la région $\alpha$ à la région $\beta$. Nous aurons ainsi supprimé une région $v'_{p}$ et une région $v'_{p-1}$. D'autre part, nous n'aurons supprimé aucune région $v'_{p-2}$ si la région $\gamma$ est régulière; si, en effet, aucune des régions $v'_{p-2}$ n'est singulière, chacune d'elles appartiendra au moins à trois régions $v'_{p-1}$ et, après la suppression de $\gamma$, elle appartiendra encore au moins à deux régions $v'_{p-1}$. De même toute région $v'_{q}$ (où $q < p - 2$) faisant partie de $\gamma$ appartiendra au moins à trois régions $v'_{p-1}$ et, après la suppression de $\gamma$, elle appartiendra encore au moins à deux régions $v'_{p-1}$. La suppression de $\gamma$ ne supprime donc aucune des régions $v'_{q}$; \emph{elle ne change donc pas la valeur du nombre $N$}.

Si au contraire la région $\gamma$ est irrégulière, nous n'avons plus le droit de la supprimer, car il existera alors une région $v'_{p-2}$ qui appartiendra seulement à $\gamma$ et à une autre région $v'_{p-1}$; après la suppression de $\gamma$, elle n'appartiendrait plus qu'à une seule région $v'_{p-1}$, ce qui est inadmissible.

Que faudra-t-il faire alors? La région $\gamma$ sépare deux régions $v_{p}$ que j'ai appelées $\alpha$ et $\beta$; mais elle ne constitue pas à elle seule la frontière entre $\alpha$ et $\beta$; en effet, comme $\gamma$ est irrégulière, il existera une région $v'_{p-2}$ singulière
\origpage{104}
que j'appellerai $\delta$ et qui appartiendra à $\gamma$ et à une autre région $v'_{p-1}$ que j'appellerai $\gamma'$. Cette région $\gamma'$, d'après ce que nous avons vu plus haut, sépare les mêmes régions que $\gamma$, c'est-à-dire $\alpha$ et $\beta$.

Si la région $\delta$ est régulière, nous pourrons la supprimer et annexer $\gamma$ à $\gamma'$. La région $\gamma + \gamma'$ séparera alors $\alpha$ de $\beta$. Nous aurons ainsi diminué $\alpha_{p-1}$ et $\beta_{p-2}$\ednote{Printed $\beta_{p-2}$; apparently for $\alpha_{p-2}$, the number of the $v_{p-2}$.} d'une unité, tandis que les autres nombres $\alpha_{i}$ n'auront pas changé. $N$ n'aura donc pas changé non plus.

Si $\delta$ est irrégulière, il y aura une région $v'_{p-3}$ singulière que j'appellerai $\varepsilon$ et qui la séparera d'une autre région $\delta'$; nous supprimerons $\varepsilon$ et annexerons $\delta$ à $\delta'$, et ainsi de suite.

Nous pourrons ainsi supprimer une région $v'_{q}$, qui sépare l'une de l'autre deux régions $v'_{q+1}$, et annexer ces deux régions $v'_{q+1}$ l'une à l'autre, mais à deux conditions~:

1° Si $q$ est plus petit que $p - 1$, il faut que la région $v'_{q}$ soit singulière;

2° Et dans tous les cas, il faut qu'elle soit régulière.

Cela posé, voici dans quel ordre nous ferons les opérations~:

Je veux remonter du polyèdre $P'$ au polyèdre $P$. Je puis sans inconvénient supposer que le polyèdre $P$ n'admet pas de région singulière, mais le polyèdre $P'$ et les polyèdres intermédiaires en admettront.

Par une série de suppressions et d'annexions successives, nous remonterons de $P'$ à $P$ en passant par une série de polyèdres intermédiaires que j'appellerai
\[ P_{0} = P', \qquad P_{1}, \qquad P_{2}, \qquad \ldots, \qquad P_{m-1}, \qquad P_{m}. \]

Comment passerons-nous du polyèdre $P_{i}$ au polyèdre $P_{i+1}$?

Si dans $P_{i}$ il y a des $v'_{0}$ singuliers, j'en supprimerai un. S'il n'y en a pas, tous les $v'_{1}$ seront réguliers; s'il y a des $v'_{1}$ singuliers, j'en supprimerai un.

S'il n'y a pas de $v'_{1}$ singulier, tous les $v'_{2}$ seront réguliers; s'il y a des $v'_{2}$ singuliers, j'en supprimerai un.

Et ainsi de suite.

Si enfin il n'y a pas de $v'_{p-2}$ singulier, tous les $v'_{p-1}$ seront réguliers et l'on aura le droit de supprimer un quelconque d'entre eux; si l'une des régions $v_{p}$ est subdivisée en plusieurs régions $v'_{p}$, je choisirai deux de ces
\origpage{105}
régions $v'_{p}$ qui seront contiguës et séparées l'une de l'autre par une région $v'_{p-1}$ qui sera leur frontière commune; je les annexerai l'une à l'autre en supprimant cette frontière commune.

Aucune de ces opérations ne peut altérer le nombre $N$.

On ne sera arrêté que quand il n'y aura plus de région singulière et que, d'autre part, aucune des régions $v_{p}$ ne sera plus subdivisée en plusieurs régions $v'_{p}$. Mais alors on sera arrivé au polyèdre $P$.

Aucune de ces opérations ne peut altérer le nombre $N$.

Ce nombre a donc la même valeur pour $P$ et pour $P'$. C.~Q.~F.~D.

Cette démonstration pourrait donner lieu à certaines objections, car on pourrait se demander si, dans cette série d'opérations, toutes les régions resteront simplement connexes; mais, avant de modifier notre démonstration de façon à me mettre à l'abri de ces objections, je veux déterminer quelle doit être la valeur du nombre $N$ pour un polyèdre simplement connexe.

Si notre théorème est vrai, le nombre $N$ doit avoir la même valeur pour deux polyèdres obtenus en subdivisant deux variétés homéomorphes; il a donc même valeur pour deux polyèdres simplement connexes quelconques.

Il nous suffira donc de faire cette détermination pour un polyèdre simplement connexe arbitrairement choisi.

Je choisirai le tétraèdre généralisé.

J'appelle ainsi le polyèdre formé par la frontière complète du domaine
\[ \begin{aligned} &x_{1} > 0, &&x_{2} > 0, &&\ldots, &&x_{p} > 0,\\ &x_{p+1} > 0, &&x_{1} + x_{2} + \ldots + x_{p} + x_{p+1} < 0. \end{aligned} \]\ednote{Printed $< 0$; apparently for $< 1$, since with every $x_{i} > 0$ the sum cannot be negative.}

On a alors
\[ \begin{aligned} &\alpha_{p} = p + 2, &&\alpha_{p-1} = \frac{(p + 1)(p + 2)}{2}, &&\ldots,\\ &\alpha_{q} = \frac{(p + 2)!}{(q + 2)!(p - q)!}, &&\ldots, &&\alpha_{1} = \frac{(p + 1)(p + 2)}{2}, &&\alpha_{0} = p + 2, \end{aligned} \]\ednote{The denominator is printed $(q + 2)!(p - q)!$; the values $\alpha_{p} = p + 2$ and $\alpha_{1}$ printed alongside apparently call for $(q + 1)!(p - q + 1)!$.}
c'est-à-dire que les nombres $\alpha_{q}$ sont égaux aux coefficients du binome.
\origpage{106}
On aura donc
\[ (1 - 1)^{p+2} = 1 - \alpha_{p} + \alpha_{p-1} - \ldots \pm \alpha_{1} \mp \alpha_{0} \pm 1 = 1 - N \pm 1. \]

On doit prendre devant le dernier terme le signe $+$ si $p$ est pair, et le signe $-$ si $p$ est impair.

On aura donc $N = 2$ si $p$ est pair, $N = 0$ si $p$ est impair.

Je serais arrivé au même résultat en choisissant le cube généralisé.

J'appelle ainsi le polyèdre formé par la frontière complète du domaine
\[ -1 < x_{i} < 1, \qquad (i = 1, 2, \ldots, p + 1). \]

On a alors
\[ \begin{aligned} &\alpha_{p} = 2(p + 1), &&\alpha_{p-1} = 2^{2}\frac{p(p + 1)}{2}, &&\ldots, &&\alpha_{q} = 2^{p-p+1}\frac{(p + 1)!}{(q + 1)!(p - q)!}, &&\ldots,\\ &\alpha_{1} = 2^{p}(p + 1), &&\alpha_{0} = 2^{p+1}. \end{aligned} \]\ednote{The exponent is printed $p-p+1$, apparently for $p-q+1$; the factor $\frac{(p + 1)!}{(q + 1)!(p - q)!}$ as printed does not reproduce the $\alpha_{p}$ and $\alpha_{p-1}$ printed alongside, which apparently call for $\frac{(p + 1)!}{q!(p - q + 1)!}$.}
d'où
\[ (1 - 2)^{p+1} = 1 - \alpha_{p} + \alpha_{p-1} - \ldots \pm \alpha_{1} \mp \alpha_{0} = 1 - N, \]
d'où
\[ N = 1 - (-1)^{p+1}, \]
c'est-à-dire

$N = 2$ si $p$ est pair,

$N = 0$ si $p$ est impair.

Ainsi, pour un polyèdre simplement connexe, le nombre $N$ est égal à 2 si $p$ est pair et à 0 si $p$ est impair.

Cela posé, pour établir notre théorème d'une façon complète et rigoureuse, je vais supposer qu'il soit vrai pour les variétés de moins de $p$ dimensions.

Considérons donc notre polyèdre $P$ et une région $v_{q}$ à $q$ dimensions appartenant à ce polyèdre. Cette région $v_{q}$ fera partie d'un certain nombre de régions $v_{q+1}$, d'un certain nombre de régions $v_{q+2}$, \ldots, d'un certain nombre de régions $v_{p}$. L'ensemble de toutes ces régions formera ce que j'appellerai l'\emph{aster} de $v_{q}$.
\origpage{107}

Je désignerai par $\gamma_{h}$ le nombre des régions $v_{h}$ $(h > q)$ qui font ainsi partie de l'aster de $v_{q}$.

Soit $(x^{0}_{1}, x^{0}_{2}, \ldots, x^{0}_{n})$ un point de $v_{q}$; considérons l'hypersphère $S$ qui a pour équation
\[ (x_{1} - x^{0}_{1})^{2} + (x_{2} - x^{0}_{2})^{2} + \ldots + (x_{n} - x^{0}_{n})^{2} = \varepsilon^{2}, \]
$\varepsilon$ étant très petit.

Soit $\Pi$ la variété plane définie par les $q$ équations
\[ \begin{aligned} &A^{1}_{i}(x_{1} - x^{0}_{1}) + A^{2}_{i}(x_{2} - x^{0}_{2}) + \ldots + A^{n}_{i}(x_{n} - x^{0}_{n}) = 0,\\ &(i = 1, 2, \ldots, q), \end{aligned} \]
où les $A^{k}_{i}$ sont des constantes quelconques.

L'intersection de $S$, $\Pi$ et $V$ sera une variété à $p - q - 1$ dimensions que j'appellerai $W$ et qui sera simplement connexe. Les intersections de $S$ et $\Pi$ avec les diverses régions qui forment l'aster de $v_{q}$ formeront par leur ensemble un polyèdre dû à la subdivision de $W$ et qui sera simplement connexe.

Pour ce polyèdre, le nombre $\alpha_{h}$ sera égal à $\gamma_{h+q+1}$; comme il a moins de $p$ dimensions, notre théorème lui sera applicable, de sorte que nous pourrons écrire
\[ \gamma_{p} - \gamma_{p-1} + \ldots \pm \gamma_{q+2} \mp \gamma_{q+1} = 2 \text{ ou } 0, \tag{A} \]
suivant que $p - q$ sera pair ou impair\ednote{Printed ``pair ou impair''. Since $W$ has $p - q - 1$ dimensions, the value 2 apparently goes with $p - q$ odd; compare p. 112, ``impair ou pair''.}.

Définissons maintenant un polyèdre $Q$ qui sera formé par une opération que l'on pourra appeler un \emph{quadrillage}.

Soit $V$ une variété à $p$ dimensions située dans l'espace à $n$ dimensions. Construisons une infinité de variétés planes définies par les équations
\[ \begin{aligned} &x_{i} = a_{k.i},\\ &(i = 1, 2, \ldots, n; \qquad K = -\infty, \ldots, -1, 0, +1, +2, \ldots, +\infty). \end{aligned} \tag{B} \]\ednote{Printed capital $K$, apparently for the index $k$ of $a_{k.i}$. The first $\infty$ is printed in an open, damaged-looking form.}
\origpage{108}

Ces variétés planes décomposeront l'espace en une infinité de domaines $D_{n}$ assimilables à des parallélépipèdes rectangles. Les frontières des $D_{n}$ seront formées par un certain nombre de domaines $D_{n-1}$ à $n - 1$ dimensions faisant partie des diverses variétés planes $x_{i} = a_{k.i}$ et également assimilables à des parallélépipèdes rectangles. Les frontières des $D_{n-1}$ seront formées par un certain nombre de domaines $D_{n-2}$ assimilables à des parallélépipèdes rectangles dans l'espace à $n - 2$ dimensions, et ainsi de suite.

Alors le polyèdre $Q$ sera défini de la façon suivante~: les régions $v_{p}$ seront formées par les intersections de $V$ avec les domaines $D_{n}$, les régions $v_{p-1}$ par les intersections de $V$ avec les domaines $D_{n-1}$, et ainsi de suite; enfin, les régions $v_{0}$ par les intersections de $V$ avec les domaines $D_{n-1}$\ednote{Printed $D_{n-1}$; the construction apparently calls for $D_{n-p}$, since $V$ has $p$ dimensions.}.

Il résulte de cette définition que le polyèdre $Q$ n'admet pas de région singulière.

Je considérerai en outre un polyèdre $P$ quelconque congruent à $Q$, et un polyèdre $P'$ dérivé à la fois de $P$ et de $Q$.

Je vais remonter, d'une part, de $P'$ à $P$ et, d'autre part, de $P'$ à $Q$, et j'établirai que, dans ces deux opérations, le nombre $N$ n'a pas changé.

Remontons d'abord de $P'$ à $P$.

Soit $x_{i} = a$ une des variétés planes définies par l'équation (B). Classons les régions $v'_{q}$ d'un nombre quelconque de dimensions qui composent le polyèdre $P'$ en quatre sortes.

Celles de la première sorte sont celles qui font partie de la variété
\[ x_{i} = a. \]

Celles de la deuxième sorte sont celles qui admettent des points tels que
\[ x_{i} = a + \varepsilon, \]
$\varepsilon$ étant positif et très petit.

Celles de la troisième sorte sont celles qui admettent des points tels que
\[ x_{i} = a - \varepsilon. \]

Toutes les autres sont de la quatrième sorte.
\origpage{109}

Soient $\delta_{q}$, $\delta'_{q}$, $\delta''_{q}$ le nombre des variétés à $q$ dimensions qui sont respectivement de la première, de la deuxième et de la troisième sorte.

Toute variété de la deuxième sorte sera contiguë à une variété de la troisième sorte et leur frontière commune sera une variété de la première sorte qui aura une dimension de moins. Les variétés des trois premières sortes se correspondent donc chacune à chacune et l'on a
\[ \delta'_{q} = \delta''_{q} = \delta_{q-1}. \]

Il est clair d'ailleurs que $\delta'_{0} = \delta''_{0} = \delta_{p} = 0$.

Si dans l'ensemble des variétés planes (B) qui constituent le quadrillage et qui ont donné naissance aux polyèdres $Q$ et $P'$, on avait supprimé la variété $x_{i} = a$, on aurait obtenu deux polyèdres $Q_{1}$ et $P'_{1}$ plus simples que les premiers. Comparons $P'_{1}$ à $P'$.

Quand on supprime la variété plane $x_{i} = a$, on supprime les variétés de la première sorte et l'on \emph{annexe} chaque variété de la troisième sorte à la variété correspondante de la deuxième sorte. Donc en passant de $P'$ à $P'_{1}$, le nombre $\alpha_{q}$ diminue de
\[ \delta''_{q} + \delta_{q} = \delta_{q} + \delta_{q-1}. \]

En particulier, les nombres extrêmes $\alpha_{p}$ et $\alpha_{0}$ diminuent de $\delta_{p-1}$ et $\delta_{0}$. Il résulte de là que le nombre $N$ diminue de
\[ \delta_{p-1} - (\delta_{p-1} + \delta_{p-2}) + (\delta_{p-2} + \delta_{p-3}) - \ldots \pm (\delta_{1} + \delta_{0}) \mp \delta_{0} = 0. \]

Donc $N$ ne change pas. Ainsi, en supprimant la variété $x_{i} = a$, on ne change pas $N$, mais, en supprimant de la sorte toutes les variétés planes définies par (B), on remontera au polyèdre $P$. Le nombre $N$ est donc le même pour $P'$ et pour $P$.

Remontons maintenant de $P'$ à $Q$.

Soient $w_{p}$, $w_{p-1}$, \ldots, $w_{1}$, $w_{0}$ les variétés, dont l'ensemble constitue le polyèdre $Q$; soient de même $v'_{p}$, $v'_{p-1}$, \ldots, $v'_{1}$, $v'_{0}$ les variétés dont l'ensemble constitue le polyèdre $P'$.

Nous répartirons les variétés $v'_{q}$ en $p + 1$ classes.
\origpage{110}

Celles de la première classe seront celles qui feront partie d'une des régions $w_{p}$ sans faire partie d'une des régions $w_{p-1}$. Le polyèdre $P'$ étant dérivé de $Q$, nous rangerons dans cette première classe toutes les variétés $v'_{p}$ (qui sont toutes des subdivisions des $w_{p}$); celles des variétés $v'_{p-1}$, qui séparent l'une de l'autre deux variétés $v'_{p}$ faisant partie d'une même région $w_{p}$; et leurs intersections.

Celles de la deuxième classe seront celles qui feront partie d'une des régions $w_{p-1}$; sans faire partie d'une des régions $w_{p-2}$.

Celles de la troisième classe seront celles qui feront partie d'une des régions $w_{p-2}$, sans faire partie d'une des régions $w_{p-3}$, etc.

Celles de la $p^{\text{ième}}$ classe seront celles qui feront partie d'une des régions $w_{1}$, sans être un des points $w_{0}$.

Enfin, la $(p + 1)^{\text{ième}}$ classe comprendra les points $w_{0}$.

Pour remonter de $P'$ à $Q$, voici ce que je vais faire~: je commencerai par supprimer toutes les variétés de la première classe qui ont moins de $p$ dimensions, ce qui a pour effet de réunir en une seule, par annexion, toutes les régions $v'_{p}$, qui sont les subdivisions d'une même région $w_{p}$.

Je dis que cette opération n'altère pas le nombre $N$.

En effet, je puis supposer que les mailles du quadrillage qui a donné naisasnce\ednote{Printed ``naisasnce''; apparently a misprint for ``naissance''.} à $Q$ sont assez serrées, pour qu'à l'intérieur d'une de ces mailles $D_{p}$, c'est-à-dire à l'intérieur d'une des régions $w_{p}$, on ne puisse trouver des points appartenant à deux variétés différentes $v_{p-1}$, sauf dans le cas où l'on y trouve des points appartenant à l'intersection de ces deux variétés. (Je désigne toujours par $v_{q}$ les variétés dont l'ensemble forme $P$.) Plus généralement, je puis supposer qu'on ne peut pas trouver dans une région $w_{p}$ des points appartenant à plusieurs variétés $v_{q} (q < p)$ à moins qu'on n'y puisse trouver des points appartenant à l'intersection de ces diverses variétés.

Dans une région $w_{p}$, nous pourrons donc avoir des points d'une région $v_{q}$ et de toutes les régions $v_{h} (h > q)$, qui font partie de son aster.

Mais nous ne pourrons avoir des points de deux régions $v_{q}$, ou des points de $v_{q}$ et d'une région $v_{h} (h > q)$ ne faisant partie de l'aster de $v_{q}$, sans avoir, en outre, des points d'une région $v_{q-1}$.
\origpage{111}

Si alors je suppose que $q$ est la plus petite valeur que l'on puisse donner au nombre des dimensions d'une région $v_{q}$, pour que cette région ait des points à l'intérieur de $w_{p}$, si je suppose cela, dis-je, nous n'aurons dans $w_{p}$ que des points d'une \emph{seule} région $v_{q}$ et des régions de son aster.

Considérons alors une région $w_{p}$, contenant des points de $v_{q}$ et des régions de l'aster de $v_{q}$. Soient
\[ \gamma_{q+1}, \qquad \gamma_{q+2}, \qquad \ldots, \qquad \gamma_{p} \]
les nombres que nous avons définis plus haut en définissant l'aster.

Il en résulte qu'à l'intérieur de $w_{p}$ nous aurons
\[ \begin{matrix} 1 & \text{région} & v'_{q} & \text{de la première classe,} \\[1ex] \gamma_{q+1} & \text{»} & v'_{q+1} & \text{»} \\[1ex] \gamma_{q+2} & \text{»} & v'_{q+2} & \text{»} \\[1ex] \ldots\ldots & \ldots\ldots & \ldots\ldots & \ldots\ldots \\[1ex] \gamma_{p} & \text{»} & v'_{p} & \text{»} \end{matrix} \]

En supprimant les régions de la première classe de moins de $p$ dimensions et en annexant les unes aux autres les $\gamma_{p}$ régions $v'_{p}$, dont l'ensemble constitue $w_{p}$, nous diminuons $\alpha_{p}$ de $\gamma_{p} - 1$, $\alpha_{p-1}$ de $\gamma_{p-1}$, \ldots, $\alpha_{q+1}$ de $\gamma_{q+1}$, $\alpha_{q}$ de 1. Donc, en vertu de l'équation (A), le nombre $N$ ne change pas.

Cela fait, nous supprimerons toutes les variétés de la deuxième classe de moins de $p = 1$\ednote{Printed $p = 1$ with two clean strokes; apparently for $p - 1$, as the parallel $p - 2$ below and the $v'_{p-1}$ that follows indicate.} dimensions, en annexant les unes aux autres toutes les variétés $v'_{p-1}$ de la deuxième classe, qui font partie d'une même région $w_{p-1}$. Ensuite, nous supprimerons toutes les variétés de la troisième classe de moins de $p - 2$ dimensions, en annexant les unes aux autres toutes les variétés $v'_{p-2}$ qui font partie d'une même région $w_{p-2}$; et ainsi de suite.

Nous arriverons enfin ainsi au polyèdre $Q$.

On montrerait de la même manière que plus haut qu'aucune de ces opérations n'altère le nombre $N$.

Le nombre $N$ est donc le même pour $P'$ et $Q$; il est donc aussi le même
\origpage{112}
pour $P$ et $Q$ et, par conséquent, pour deux polyèdres congruents quelconques. C.~Q.~F.~D.

\subsection*{\S~17. --- Cas où $p$ est impair.}

Je vais définir à l'égard d'un polyèdre quelconque $P$ de nouveaux nombres remarquables que j'appellerai $\beta_{\lambda.\mu}$.

Soit d'abord $\lambda > \mu$; je considérerai toutes les variétés $v_{\lambda}$; pour chacune d'elles, j'envisagerai le nombre des variétés $v_{\mu}$ qui en font partie; je ferai la somme de tous ces nombres relativement aux diverses variétés $v_{\lambda}$ et ce sera cette somme que j'appellerai $\beta_{\lambda.\mu}$.

Comme les variétés $v_{\lambda}$ sont toutes par hypothèse simplement connexes, on aura
\[ \beta_{\lambda.\lambda-1} - \beta_{\lambda.\lambda-2} + \ldots \pm \beta_{\lambda.1} \mp \beta_{\lambda.0} = 2\alpha_{\lambda} \text{ ou } 0, \]
selon que $\lambda$ sera impair ou pair.

Soit maintenant $\lambda < \mu$; je considérerai toutes les variétés $v_{\lambda}$; pour chacune d'elles, j'envisagerai le nombre des variétés $v_{\mu}$, dont elle fait partie (c'est-à-dire le nombre $\gamma_{\mu}$ relatif à l'aster de $v_{\lambda}$); je ferai la somme de tous ces nombres relatifs aux diverses variétés $v_{\lambda}$ et ce sera cette somme que j'appellerai $\beta_{\lambda.\mu}$.

En vertu de l'équation (A) du numéro précédent, on aura
\[ \beta_{\lambda.p} - \beta_{\lambda.p-1} + \ldots \pm \beta_{\lambda.\lambda+2} \mp \beta_{\lambda.\lambda+1} = 2\alpha_{\lambda} \text{ ou } 0, \]
selon que $p - \lambda$ sera impair ou pair.

Il résulte de cette définition que
\[ \beta_{\lambda.\mu} = \beta_{\mu.\lambda}. \]

Cela posé, formons le Tableau suivant
\[ \begin{matrix} +\beta_{p.p-1}, & -\beta_{p.p-2}, & +\beta_{p.p-3}, & \ldots, & \pm\beta_{p.1}, & \mp\beta_{p.0}, \\[1ex] & +\beta_{p-1.p-2}, & -\beta_{p-1.p-3}, & \ldots, & \mp\beta_{p-1.1}, & \pm\beta_{p-1.0}, \\[1ex] & \ldots\ldots, & \ldots\ldots, & \ldots, & \ldots\ldots, & \ldots\ldots, \\[1ex] & & \ldots\ldots, & \ldots, & \ldots\ldots, & \ldots\ldots, \\[1ex] & & & \ldots, & \ldots\ldots, & \ldots\ldots, \\[1ex] & & & & +\beta_{2.1}, & -\beta_{2.0}, \\[1ex] & & & & & +\beta_{1.0}. \end{matrix} \]
\origpage{113}

On voit que, dans chaque ligne et dans chaque colonne, chaque terme est affecté alternativement du signe $+$ et du signe $-$, de telle façon que $\beta_{\lambda.\mu}$ est affecté du signe $+$, si $\lambda - \mu$ est impair et du signe $-$ dans le cas contraire.

Faisons la somme des termes de ce Tableau; cette somme peut s'effectuer de deux manières; en sommant d'abord les lignes et en sommant d'abord les colonnes.

La somme des termes des diverses lignes du Tableau est, en commençant par le haut~:
\[ 2\alpha_{p}, \qquad 0, \qquad 2\alpha_{p-2}, \qquad 0, \qquad \ldots, \qquad 2\alpha_{3}, \qquad 0, \qquad 2\alpha_{1} \]
si $p$ est impair, et
\[ 0, \qquad 2\alpha_{p-1}, \qquad 0, \qquad 2\alpha_{p-3}, \qquad \ldots, \qquad 2\alpha_{3}, \qquad 0, \qquad 2\alpha_{1} \]
si $p$ est pair. La somme des termes du Tableau est donc~:
\[ 2\alpha_{1} + 2\alpha_{3} + \ldots + 2\alpha_{p} \]
si $p$ est impair, et
\[ 2\alpha_{1} + 2\alpha_{3} + \ldots + 2\alpha_{p-1} \]
si $p$ est pair.

La somme des termes des diverses colonnes du Tableau est, en commençant par la gauche~:
\[ 2\alpha_{p-1}, \qquad 0, \qquad 2\alpha_{p-3}, \qquad 0, \qquad \ldots, \qquad 2\alpha_{2}, \qquad 0, \qquad 2\alpha_{0} \]
si $p$ est impair, et
\[ 2\alpha_{p-1}, \qquad 0, \qquad 2\alpha_{p-3}, \qquad 0, \qquad \ldots, \qquad 0, \qquad 2\alpha_{1}, \qquad 0 \]
si $p$ est pair.

La somme des termes du Tableau est donc~:
\[ 2\alpha_{0} + 2\alpha_{2} + \ldots + 2\alpha_{p-1} \]
si $p$ est impair et
\[ 2\alpha_{1} + 2\alpha_{3} + \ldots + 2\alpha_{p-1} \]
si $p$ est pair.
\origpage{114}

En égalant les deux expressions de cette somme, on obtiendra une identité, si $p$ est pair, et l'équation
\[ N = 0 \]
si $p$ est impair.

D'où cette conséquence~:

Le nombre $N$ est nul et indépendant des nombres de Betti si $p$ est impair; il dépend, au contraire, des nombres de Betti si $p$ est pair.

\subsection*{\S~18. --- Deuxième démonstration.}

La deuxième démonstration va nous apprendre comment il en dépend.

Pour bien la faire comprendre, je vais d'abord l'exposer pour les polyèdres ordinaires avec $\alpha_{0}$ sommets, $\alpha_{1}$ arêtes et $\alpha_{2}$ faces.

A chacun des $\alpha_{0}$ sommets, je fais correspondre un nombre arbitraire; à chacune des $\alpha_{1}$ arêtes, je fais correspondre un nombre $\delta$ égal à la différence des nombres correspondant à ses deux sommets.

Nous aurons ainsi $\alpha_{1}$ différences $\delta$; mais elles ne pourront pas toutes être choisies arbitrairement; en effet, elles seront déterminées, quand on connaîtra les $\alpha_{0}$ nombres affectés aux divers sommets et même quand on connaîtra les excès de $\alpha_{0} - 1$ de ces nombres sur le dernier. Il y aura donc en tout $\alpha_{0} - 1$ différences $\delta$, qui pourront être choisies arbitrairement.

Il doit donc y avoir entre les différences $\delta$
\[ \alpha_{1} - \alpha_{0} + 1 \]
relations linéaires.

Il est clair que l'on pourra obtenir toutes ces relations linéaires de la manière suivante~: Considérons une suite d'arêtes formant un contour fermé; la somme algébrique des différences $\delta$ relations\ednote{Printed ``relations''; apparently for ``relatives''.} aux diverses arêtes de cette suite devra être nulle.

Cherchons donc à construire des contours fermés formés d'arêtes.

Nous avons d'abord les contours polygonaux de chaque face; ils sont au nombre de $\alpha_{2}$.
\origpage{115}

Si ensuite le polyèdre n'est pas simplement connexe, nous pourrons tracer à sa surface $P_{1} - 1$ contours fermés linéairement indépendants au sens donné à ce mot dans le paragraphe sur les homologies. Soit $C$ l'un de ces contours, il traversera successivement différentes faces; soient $a_{1}a_{2}, a_{2}a_{3}, \ldots, a_{n-1}a_{n}, a_{n}a_{1}$ les arcs de ce contour, qui sont dans chacune de ces faces.

Prenons le premier de ces arcs $a_{1}a_{2}$ et $F$ la face correspondante.

Le point $a_{1}$ et le point $a_{2}$ sont sur le périmètre de cette face; nous pouvons alors aller de $a_{1}$ en $a_{2}$ en suivant ce périmètre par un chemin que j'appellerai $a_{1}m_{1}a_{2}$. Nous aurons alors
\[ a_{1}m_{1}a_{2} + a_{2}a_{1} \backsim 0, \]
c'est-à-dire que, dans le contour $C$, on peut remplacer l'arc $a_{1}a_{2}$ par l'arc $a_{1}m_{1}a_{2}$; opérons de même pour les autres arcs de $C$; nous aurons finalement remplacé $C$ par le contour homologue
\[ a_{1}m_{1}a_{2} + a_{2}m_{2}a_{3} + \ldots + a_{n}m_{n}a_{1}, \]
que j'appellerai $C'$. Ce contour $C'$ se composera d'un certain nombre d'arêtes et de portions d'arêtes. Par exemple, l'arc $a_{1}m_{1}a_{2}$ se composera d'un segment d'arête joignant $a_{1}$ au sommet le plus voisin, puis d'un certain nombre d'arêtes complètes, puis d'un segment $Sa_{2}$ joignant $a_{2}$ au sommet le plus voisin.

Mais ce segment $Sa_{2}$ se retrouvera parcouru en sens contraire dans l'arc $a_{2}m_{2}a_{3}$. Les portions d'arêtes, qui font partie de $C'$, sont donc parcourues deux fois en sens contraires; nous pouvons donc les supprimer et nous obtenons ainsi un contour $C''$ exclusivement formé d'arêtes complètes et homologue à $C$ ou à $C'$.

Nous aurons $P_{1} - 1$ contours analogues à $C''$, qui nous donneront $P_{1} - 1$ relations entre les $\delta$.

Nous obtenons ainsi $\alpha_{2} + P_{1} - 1$ contours fermés formés d'arêtes; je dis que tous les autres contours possibles n'en sont que des combinaisons.

Soit d'abord un contour fermé formé d'arêtes et homologue à 0, il
\origpage{116}
partagera la surface du polyèdre en deux régions. Soit $R$ l'une d'elles, elle sera évidemment formée d'un certain nombre $q$ de faces, puisque le contour est exclusivement formé d'arêtes; nous pourrons donc remplacer le contour donné par $q$ contours partiels formés par les périmètres de ces $q$ faces.

Si maintenant le contour donné n'est pas homologue à 0, nous pourrons toujours le remplacer par une combinaison des contours $C''$ et par un contour homologue à 0.

Nous avons donc
\[ \alpha_{2} + P_{1} - 1 \]
relations entre nos $\delta$ et nous n'en avons pas d'autres; mais sont-elles toutes distinctes?

Pour le reconnaître, il faut voir si l'on peut former une combinaison linéaire de ces relations qui se réduise à une identité, ou, ce qui revient au même, si l'on peut former une combinaison de nos $\alpha_{2} + P_{1} - 1$ contours, une combinaison telle que chaque arête soit parcourue deux fois en sens contraire (ou, si elle l'est plus de deux fois, qu'elle le soit autant de fois dans un sens que dans l'autre).

Peut-on d'abord former une telle combinaison avec les $\alpha_{2}$ périmètres $\Pi$? Dire que chaque arête est parcourue ainsi deux fois en sens contraire, c'est dire que l'ensemble des polygones, dont on parcourt ainsi les périmètres, forme un polyèdre fermé.

Or, on ne peut évidemment construire de la sorte qu'un seul polyèdre fermé, qui est le polyèdre donné.

Avec les $\alpha_{2}$ relations correspondant aux $\alpha_{2}$ périmètres $\Pi$, nous pouvons donc former une combinaison linéaire identique et une seule.

Peut-on former maintenant de pareilles combinaisons avec les périmètres $\Pi$ et les contours $C''$? Si cela était, l'ensemble des polygones, dont on parcourt ainsi les périmètres, formerait une surface polyédrique non fermée, dont la frontière complète serait formée par une combinaison des contours $C''$. Or, cela est impossible, puisque les contours $C''$ sont linéairement indépendants.
\origpage{117}

On peut donc former avec nos $\alpha_{2} + P_{1} - 1$ relations une combinaison identique et une seule. Il y a donc entre les $\delta$
\[ (\alpha_{2} + P_{1} - 1) - 1 \]
relations distinctes. Le nombre des $\delta$ arbitraires est donc
\[ \alpha_{1} - (\alpha_{2} + P_{1} - 1) + 1, \]
de sorte qu'il vient
\[ \alpha_{0} - 1 = \alpha_{1} - (\alpha_{2} + P_{1} - 1) + 1, \]
ou
\[ N = 3 - P_{1}. \]

Étendons cette démonstration au cas où nous avons un polyèdre à trois dimensions, où l'on distinguera~:
\[ \begin{matrix} \text{Les sommets} & v_{0} & \text{au nombre de} & \alpha_{0}, \\[1ex] \text{les variétés} & v_{1} & \text{»} & \alpha_{1}, \\[1ex] \text{»} & v_{2} & \text{»} & \alpha_{2}, \\[1ex] \text{»} & v_{3} & \text{»} & \alpha_{3}. \end{matrix} \]

A chacun des $v_{0}$, faisons correspondre un nombre, et à chacun des $v_{1}$ la différence $\delta$ des nombres correspondant à ces deux sommets. Nous aurons ainsi $\alpha_{1}$ différences $\delta$, dont $\alpha_{0} - 1$ seront arbitraires.

On obtiendra les relations linéaires qui ont lieu entre les $\delta$, en construisant tous les contours fermés exclusivement formés avec des $v_{1}$.

Nous aurons d'abord les périmètres $\Pi$ des $v_{2}$, qui seront au nombre de $\alpha_{2}$; soit ensuite $C$ un contour fermé quelconque non homologue à 0, il traversera successivement différents $v_{3}$; soient $a_{1}a_{2}$, $a_{2}a_{3}$, \ldots\ les arcs de $C$ qui se trouvent dans chacun de ces $v_{3}$.

Considérons l'arc $a_{1}a_{2}$ qui se trouve dans un des $v_{3}$, que j'appelle $\Phi$, et dont les extrémités $a_{1}$ et $a_{2}$ se trouvent dans deux des $v_{2}$, qui servent de frontière à $\Phi$, à savoir $a_{1}$ dans $\varphi_{1}$ et $a_{2}$ dans $\varphi_{2}$.

Soient $b_{1}$ un sommet de $\varphi_{1}$, $b_{2}$ un sommet de $\varphi_{2}$; joignons $a_{1}$ à $b_{1}$ par
\origpage{118}
une ligne quelconque $a_{1}b_{1}$, puis $b_{1}$ à $b_{2}$ par une ligne $b_{1}b_{2}$ formée exclusivement par des $v_{1}$ appartenant à la frontière de $\Phi$, et $b_{2}$ à $a_{2}$ par une ligne quelconque $a_{2}b_{2}$; on aura
\[ a_{1}a_{2} \backsim a_{1}b_{1} + b_{1}b_{2} + b_{2}a_{2}, \]
de sorte que nous pourrons remplacer l'arc $a_{1}a_{2}$ par l'arc $a_{1}b_{1}b_{2}a_{2}$; opérons de même pour les autres arcs de $C$. Nous aurons remplacé $C$ par le contour homologue
\[ a_{1}b_{1}b_{2}a_{2} + a_{2}b_{2}b_{3}a_{3} + \ldots \]
que j'appelle $C'$. Le contour $C'$ se compose d'un certain nombre de $v_{1}$ et d'arcs analogues à $a_{1}b_{1}$, $a_{2}b_{2}$, \ldots, parcourus chacun deux fois en sens contraire. On peut donc supprimer ces arcs, et il reste un contour $C''$ homologue à $C$ et exclusivement formé de $v_{1}$.

Il y aura $P_{1} - 1$ contours $C''$.

Je dis maintenant que tout contour fermé $K$ formé de $v_{1}$ est une combinaison des $\Pi$ et des $C''$. Si
\[ K \backsim 0, \]
le contour $K$ sera la frontière complète d'une certaine région $R$ à deux dimensions. Cette région $R$ pourra être subdivisée en un certain nombre de variétés $r$, chacune des régions $r$ étant la portion de $R$ qui se trouve à l'intérieur de l'un des $v_{3}$. Considérons une des variétés $r$, et soit $\varphi$ la région $v_{3}$ dans laquelle elle se trouve; la frontière complète de $r$ sera une variété fermée $u$ à une dimension, qui fera partie de la frontière complète de $\varphi$. Comme $\varphi$ est simplement connexe, $u$ partagera la frontière complète de $\varphi$ en deux régions; soit $r'$ l'une de ces régions; elle se composera d'un certain nombre de $v_{2}$ qui en feront partie tout entières et d'un certain nombre de portions de $v_{2}$ (puisque, parmi les $v_{2}$ qui forment la frontière complète de $\varphi$, il y en a qui sont partagées en deux parties par $u$). On voit que $r$ est homologue à $r'$; nous pourrons remplacer $r$ par $r'$; si nous opérons de même pour toutes les régions $r$, nous aurons obtenu une variété
\origpage{119}
$R'$ homologue à $R$ et qui se composera d'un certain nombre de $v_{2}$ complètes et d'un certain nombre de portions de $v_{2}$ prises chacune deux fois en sens contraire. Nous pourrons supprimer ces portions de $v_{2}$ et nous obtiendrons ainsi une variété $R''$ homologue à $R$ et limitée par le même contour $K$. Cette variété $R''$ pourra être décomposée en un certain nombre de polygones $v_{2}$, de sorte que le contour $K$ pourra être remplacé par les périmètres $\Pi$ de ces polygones.

Si $K$ n'est pas homologue à 0, on pourra le remplacer par un certain nombre de contours $C''$ et par un contour analogue\ednote{Printed ``analogue''; apparently for ``homologue'', as in the corresponding sentence on p. 116.} à 0.

Nous avons donc
\[ \alpha_{2} + P_{1} - 1 \]
relations entre les $\delta$, que j'écrirai
\[ \varepsilon = 0, \]
et nous n'en aurons pas d'autres. Mais sont-elles toutes distinctes?

En d'autres termes, pouvons-nous former une combinaison linéaire des $\varepsilon$ qui soit identiquement nulle, ou encore une combinaison des $\Pi$ et des $C''$ où chaque $v_{1}$ figure deux fois en sens contraire?

Si dans cette combinaison devaient figurer des $C''$, l'ensemble des polygones $v_{2}$, dont les périmètres $\Pi$ figureraient dans la combinaison, formerait une variété à deux dimensions, qui ne serait pas fermée et dont la frontière complète serait formée par un certain nombre de contours $C''$. Or cela n'est pas possible, puisque les $C''$ sont linéairement indépendants.

Il me reste donc à examiner les combinaisons où ne figureraient que des $\Pi$; l'ensemble des polygones $v_{2}$ dont les périmètres $\Pi$ y figureraient formerait alors une variété fermée.

Nous sommes ainsi conduits à examiner les variétés fermées exclusivement formées de $v_{2}$.

Nous avons d'abord les frontières complètes des $v_{3}$, frontières complètes que j'appellerai $\Phi$ et qui sont au nombre de $\alpha_{3}$.

Soit ensuite $D$ une variété quelconque à deux dimensions, non homologue
\origpage{120}
à 0. Traitons-la comme nous avons traité $R$; nous verrons qu'elle est homologue à une variété $D''$ fermée et à deux dimensions, exclusivement formée de $v_{2}$.

Le nombre des $D''$ est $P_{2} - 1$.

Je dis maintenant que toute variété fermée $K$ formée de $v_{2}$ est une combinaison des $\Pi$ et des $D''$. Si, en effet, elle est homologue à 0, elle limite une région $S$ à trois dimensions qui se composera d'un certain nombre de $v_{3}$, puisque $K$ se compose d'un certain nombre de $v_{2}$; on pourra donc remplacer $K$ par les frontières complètes $\Phi$ de ces $v_{3}$. Si $K$ n'est pas homologue à 0, on pourra la remplacer par un certain nombre de $D''$ et par une variété homologue à 0 et exclusivement formée de $v_{2}$.

Nous avons donc, entre les $\varepsilon$,
\[ \alpha_{3} + P_{2} - 1 \]
relations linéaires, que j'écrirai
\[ \zeta = 0, \]
et nous n'en aurons pas d'autres. Sont-elles distinctes?

Pour former une combinaison linéaire des $\zeta$ qui soit identiquement nulle, il faut former une combinaison des $\Phi$ et des $D''$ telle que chaque $v_{2}$ y figure deux fois en sens contraire. Nous verrions, comme plus haut, que les $D''$ ne peuvent pas figurer dans cette combinaison et que l'ensemble des $v_{3}$, dont les frontières $\Phi$ y figurent, doit former une variété fermée à trois dimensions. Or, nous ne pouvons construire qu'une seule variété de cette sorte~: c'est le polyèdre donné lui-même.

Il y a donc une combinaison linéaire des $\zeta$ qui s'annule identiquement.

Il y aura donc
\[ (\alpha_{3} + P_{2} - 1) - 1 \]
relations linéaires distinctes entre les $\varepsilon$,
\[ (\alpha_{2} + P_{1} - 1) - (\alpha_{3} + P_{2} - 1) + 1 \]
relations linéaires distinctes entre les $\delta$.
\origpage{121}

Parmi les $\delta$, il en restera donc
\[ \alpha_{1} - (\alpha_{2} + P_{1} - 1) + (\alpha_{3} + P_{2} - 1) - 1 \]
qui seront arbitraires; on trouve ainsi
\[ \alpha_{0} - 1 = \alpha_{1} - (\alpha_{2} + P_{1} - 1) + (\alpha_{3} + P_{2} - 1) - 1 \]
et l'on trouverait
\[ \alpha_{0} - 1 = \alpha_{1} - (\alpha_{2} + P_{1} - 1) + (\alpha_{3} + P_{2} - 1) - (\alpha_{4} + P_{3} - 1) + 1 \]
pour les polyèdres à quatre dimensions.

On a donc
\[ N = P_{2} - P_{1} \]
pour les polyèdres à trois dimensions, et
\[ N = 3 - P_{1} + P_{2} - P_{3} \]
pour les polyèdres à quatre dimensions.

En général, on aura
\[ N = P_{p-1} - P_{p-2} + \ldots + P_{2} - P_{1} \]
si $p$ est impair, et
\[ N = 3 - P_{1} + P_{2} - \ldots + P_{p-1} \]\ednote{The last sign is printed $+$; since $p-1$ is odd when $p$ is even, apparently for $-$, as in the case $p = 4$ above.}
si $p$ est pair.

Si nous observons maintenant que les nombres de Betti, également distants des extrêmes, sont égaux, on verra que l'on doit avoir
\[ N = 0 \]
si $p$ est impair, comme nous l'avions déjà vu dans le numéro précédent\ednote{No final period is printed after ``précédent''; the memoir ends here.}

\end{document}
